% Generated mechanically from frozen input inputs/p12/ja/stacks-more-morphisms.tex.
% Do not edit this generated file; edit the bound locale source instead.

% OK, start here.
%

\setcounter{chapter}{105}
\chapter{スタックの射についてさらに}



\phantomsection
\label{stacks-more-morphisms-section-phantom}


\section{序論}
\label{stacks-more-morphisms-section-introduction}

\noindent
本章では，代数スタックの射の性質の研究を続ける．
表現可能な対角射をもつ準分離代数スタックの場合については
\cite{LM-B} が参考文献である．




\section{規約と用語の濫用}
\label{stacks-more-morphisms-section-conventions}

\noindent
引き続き，Properties of Stacks, Section
\ref{stacks-properties-section-conventions} で導入された規約と
用語の濫用を用いる．






\section{厚化}
\label{stacks-more-morphisms-section-thickenings}

\noindent
以下の用語は完全に標準的とは限らないが，便利である．
$\mathcal{Y}$ が代数スタック $\mathcal{X}$ の閉部分スタックならば，
射 $\mathcal{Y} \to \mathcal{X}$ は表現可能である．

\begin{definition}
\label{stacks-more-morphisms-definition-thickening}
厚化を次のように定める．
\begin{enumerate}
\item 代数スタック $\mathcal{X}'$ が代数スタック $\mathcal{X}$ の
{\it 厚化} であるとは，$\mathcal{X}$ が $\mathcal{X}'$ の
閉部分スタックであり，付随する位相空間が等しいことをいう．
\item 二つの厚化 $\mathcal{X} \subset \mathcal{X}'$ と
$\mathcal{Y} \subset \mathcal{Y}'$ が与えられたとき，
{\it 厚化の射} とは，代数スタックの射
$f' : \mathcal{X}' \to \mathcal{Y}'$ であって，
$f'|_\mathcal{X}$ が閉部分スタック $\mathcal{Y}$ を経由するものをいう．
この状況では
$f = f'|_\mathcal{X} : \mathcal{X} \to \mathcal{Y}$ とおき，
$(f, f') : (\mathcal{X} \subset \mathcal{X}') \to
(\mathcal{Y} \subset \mathcal{Y}')$ を厚化の射という．
\item $\mathcal{Z}$ を代数スタックとする．同様に，
{\it $\mathcal{Z}$ 上の厚化} および
{\it $\mathcal{Z}$ 上の厚化の射} を定義する．
これは，代数スタック $\mathcal{X}'$ と $\mathcal{Y}'$ に
$\mathcal{Z}$ への構造射が備わり，$f'$ が代数スタックの適切な
$2$-可換図式に組み込まれることを意味する．
\end{enumerate}
\end{definition}

\noindent
$\mathcal{X} \subset \mathcal{X}'$ を代数スタックの厚化とする．
$U'$ をスキームとし，$U' \to \mathcal{X}'$ を全射かつ平滑な射とする．
$U = \mathcal{X} \times_{\mathcal{X}'} U'$ とおけば，厚化の射
$$
(U \subset U') \longrightarrow (\mathcal{X} \subset \mathcal{X}')
$$
が得られ，$U \to \mathcal{X}$ は全射かつ平滑である．厚化
$\mathcal{X} \subset \mathcal{X}'$ の性質は，しばしば対応する厚化
$U \subset U'$ の性質から導ける．用語を濫用して，射
$\mathcal{X} \to \mathcal{X}'$ が，$|\mathcal{X}| \to
|\mathcal{X}'|$ に全単射を誘導する閉埋め込みであるとき，
これを厚化ということもある．

\begin{lemma}
\label{stacks-more-morphisms-lemma-thickening}
$i : \mathcal{X} \to \mathcal{X}'$ を代数スタックの射とする．
次は同値である．
\begin{enumerate}
\item $i$ は代数スタックの厚化である（上記の用語の濫用による）．
\item $i$ は代数空間で表現可能であり，Properties of Stacks, Section
\ref{stacks-properties-section-properties-morphisms} の意味で厚化である．
\end{enumerate}
この場合，$i$ は閉埋め込みかつ普遍同相射である．
\end{lemma}

\begin{proof}
More on Morphisms of Spaces, Lemmas
\ref{spaces-more-morphisms-lemma-descending-property-thickening} および
\ref{spaces-more-morphisms-lemma-base-change-thickening} により，
代数空間の射が（一次）厚化であるという性質 $P$ は，基底上 fpqc 局所的で，
基底変換のもとで安定である．したがって，Properties of Stacks, Section
\ref{stacks-properties-section-properties-morphisms} の議論が実際に適用できる．
このことを踏まえると，(1) と (2) の同値性は
$P = P_1 + P_2$ という事実から従う．ここで $P_1$ は閉埋め込みであるという
性質，$P_2$ は全射であるという性質である．
（厳密には，More on Morphisms of Spaces, Definition
\ref{spaces-more-morphisms-definition-thickening}，
Properties of Stacks, Definition \ref{stacks-properties-definition-immersion}
とそれに続く議論，Morphisms of Spaces, Lemma
\ref{spaces-morphisms-lemma-surjective-representable}，および
Properties of Stacks, Section \ref{stacks-properties-section-surjective}
も参照し，すべての概念が整合することを確認すべきである．）
% P12-ERR-0062: Japanese removes the duplicated “all” in the frozen English clause.
最後の主張は以上から明らかである．
\end{proof}

\noindent
以後，この補題を断りなく用いる．補題で用いたのと同じ
More on Morphisms of Spaces, Lemmas
\ref{spaces-more-morphisms-lemma-descending-property-thickening} および
\ref{spaces-more-morphisms-lemma-base-change-thickening} により，
一次厚化を次のように定義できる．

\begin{definition}
\label{stacks-more-morphisms-definition-first-order-thickening}
代数スタック $\mathcal{X}'$ が代数スタック $\mathcal{X}$ の
{\it 一次厚化} であるとは，$\mathcal{X}$ が $\mathcal{X}'$ の
閉部分スタックであり，$\mathcal{X} \to \mathcal{X}'$ が
Properties of Stacks, Section
\ref{stacks-properties-section-properties-morphisms} の意味で
一次厚化であることをいう．
\end{definition}

\noindent
上のように $(U \subset U') \to (\mathcal{X} \subset \mathcal{X}')$
がスキームによる平滑被覆ならば，これは単に $U \subset U'$ が
一次厚化であることを意味する．次に，当然必要となる補題を述べる．

\begin{lemma}
\label{stacks-more-morphisms-lemma-base-change-thickening}
$\mathcal{Y} \subset \mathcal{Y}'$ を代数スタックの厚化とする．
$\mathcal{X}' \to \mathcal{Y}'$ を代数スタックの射とし，
$\mathcal{X} = \mathcal{Y} \times_{\mathcal{Y}'} \mathcal{X}'$ とおく．
このとき
$(\mathcal{X} \subset \mathcal{X}') \to (\mathcal{Y} \subset \mathcal{Y}')$
は厚化の射である．$\mathcal{Y} \subset \mathcal{Y}'$ が一次厚化ならば，
$\mathcal{X} \subset \mathcal{X}'$ も一次厚化である．
\end{lemma}

\begin{proof}
上の議論，Properties of Stacks, Section
\ref{stacks-properties-section-properties-morphisms}，および
More on Morphisms of Spaces, Lemma
\ref{spaces-more-morphisms-lemma-base-change-thickening} を参照せよ．
\end{proof}

\begin{lemma}
\label{stacks-more-morphisms-lemma-composition-thickening}
$\mathcal{X} \subset \mathcal{X}'$ と $\mathcal{X}' \subset \mathcal{X}''$
が代数スタックの厚化ならば，$\mathcal{X} \subset \mathcal{X}''$ も
代数スタックの厚化である．
\end{lemma}

\begin{proof}
上の議論，Properties of Stacks, Section
\ref{stacks-properties-section-properties-morphisms}，および
More on Morphisms of Spaces, Lemma
\ref{spaces-more-morphisms-lemma-composition-thickening}
を参照せよ．
\end{proof}

\begin{example}
\label{stacks-more-morphisms-example-reduction-thickening}
$\mathcal{X}'$ を代数スタックとする．このとき $\mathcal{X}'$ は
その被約化 $\mathcal{X}'_{red}$ の厚化である．Properties of Stacks,
Definition \ref{stacks-properties-definition-reduced-induced-stack} を参照せよ．
さらに，$\mathcal{X} \subset \mathcal{X}'$ が代数スタックの厚化ならば，
$\mathcal{X}'_{red} = \mathcal{X}_{red} \subset \mathcal{X}$ である．
言い換えれば，$\mathcal{X} = \mathcal{X}'_{red}$ であることと，
$\mathcal{X}$ が被約代数スタックであることとは同値である．
\end{example}

\begin{lemma}
\label{stacks-more-morphisms-lemma-reduced-diagonal}
$(f, f') : (\mathcal{X} \subset \mathcal{X}') \to
(\mathcal{Y} \subset \mathcal{Y}')$ を代数スタックの厚化の射とする．
このとき
$\mathcal{X} \times_\mathcal{Y} \mathcal{X} \to
\mathcal{X}' \times_{\mathcal{Y}'} \mathcal{X}'$
は厚化であり，標準図式
$$
\xymatrix{
\mathcal{X} \ar[r]_-\Delta \ar[d] &
\mathcal{X} \times_\mathcal{Y} \mathcal{X} \ar[d] \\
\mathcal{X}' \ar[r]^-{\Delta'} &
\mathcal{X}' \times_{\mathcal{Y}'} \mathcal{X}'
}
$$
はカルテシアンである．
\end{lemma}

\begin{proof}
$\mathcal{X} \to \mathcal{Y}'$ は閉部分スタック $\mathcal{Y}$ を
経由するので，
$\mathcal{X} \times_\mathcal{Y} \mathcal{X} =
\mathcal{X} \times_{\mathcal{Y}'} \mathcal{X}$ である．したがって
$\mathcal{X} \times_\mathcal{Y} \mathcal{X} \to
\mathcal{X}' \times_{\mathcal{Y}'} \mathcal{X}'$
は合成
$$
\mathcal{X} \times_{\mathcal{Y}'} \mathcal{X} \to
\mathcal{X} \times_{\mathcal{Y}'} \mathcal{X}' \to
\mathcal{X}' \times_{\mathcal{Y}'} \mathcal{X}'
$$
と同型である．各射は厚化の基底変換として厚化である
（Lemma \ref{stacks-more-morphisms-lemma-base-change-thickening}）．ゆえに，その合成も厚化である
（Lemma \ref{stacks-more-morphisms-lemma-composition-thickening}）．
$\mathcal{X} \to \mathcal{X}'$ はモノ射なので，補題の最後の主張は
$\mathcal{X} \to \mathcal{X}' \to \mathcal{Y}'$ に
Properties of Stacks, Lemma
\ref{stacks-properties-lemma-monomorphism-diagonal}
を適用すれば従う．
\end{proof}

\begin{lemma}
\label{stacks-more-morphisms-lemma-thickening-diagonals}
$(f, f') : (\mathcal{X} \subset \mathcal{X}') \to
(\mathcal{Y} \subset \mathcal{Y}')$ を代数スタックの厚化の射とする．
$\Delta : \mathcal{X} \to \mathcal{X} \times_\mathcal{Y} \mathcal{X}$ と
$\Delta' : \mathcal{X}' \to \mathcal{X}' \times_{\mathcal{Y}'} \mathcal{X}'$
を対応する対角射とする．このとき，次の各性質を $\Delta$ が満たすことと
$\Delta'$ が満たすこととは同値である：
(a) スキームで表現可能，(b) アフィン，(c) 全射，(d) 準コンパクト，
(e) 普遍閉，(f) 整，(g) 準分離，(h) 分離，(i) 普遍単射，
(j) 普遍開，(k) 局所準有限，(l) 有限，(m) 不分岐，(n) モノ射，
(o) 埋め込み，(p) 閉埋め込み，(q) 固有．
\end{lemma}

\begin{proof}
$$
(\Delta, \Delta') :
(\mathcal{X} \subset \mathcal{X}')
\longrightarrow
(\mathcal{X} \times_\mathcal{Y} \mathcal{X} \subset
\mathcal{X}' \times_{\mathcal{Y}'} \mathcal{X}')
$$
が厚化の射であることに注意する
（Lemma \ref{stacks-more-morphisms-lemma-reduced-diagonal}）．さらに，Morphisms of Stacks,
Lemma \ref{stacks-morphisms-lemma-properties-diagonal} により，
$\Delta$ と $\Delta'$ は代数空間で表現可能である．したがって，
Properties of Stacks, Section
\ref{stacks-properties-section-properties-morphisms} の議論を通じて，
場合 (a), (b), (c), (d), (e), (f), (g), (h), (i), (j) については
More on Morphisms of Spaces, Lemma
\ref{spaces-more-morphisms-lemma-thicken-property-morphisms}
を用いれば補題が従う．

\medskip\noindent
Lemma \ref{stacks-more-morphisms-lemma-reduced-diagonal} により，
$\mathcal{X} = (\mathcal{X} \times_\mathcal{Y} \mathcal{X})
\times_{(\mathcal{X}' \times_{\mathcal{Y}'} \mathcal{X}')} \mathcal{X}'$
である．また，上で挙げた Morphisms of Stacks, Lemma
\ref{stacks-morphisms-lemma-properties-diagonal} により，
$\Delta$ と $\Delta'$ は局所有限型である．したがって，
場合 (k), (l), (m), (n), (o), (p), (q) の結果は
More on Morphisms of Spaces, Lemma
\ref{spaces-more-morphisms-lemma-properties-that-extend-over-thickenings}
を用いれば従う．
% P12-ERR-0063: Japanese supplies the finite verb omitted in the frozen English inference.
\end{proof}

\noindent
その帰結として，次の好ましい結果を得る．

\begin{lemma}
\label{stacks-more-morphisms-lemma-thickening-properties}
\begin{reference}
\cite[Theorem 2.2.5]{Conrad-moduli}
\end{reference}
$\mathcal{X} \subset \mathcal{X}'$ を代数スタックの厚化とする．このとき
\begin{enumerate}
\item $\mathcal{X}$ が代数空間であることと $\mathcal{X}'$ が
代数空間であることとは同値である．
\item $\mathcal{X}$ がスキームであることと $\mathcal{X}'$ が
スキームであることとは同値である．
\item $\mathcal{X}$ が DM であることと $\mathcal{X}'$ が DM であることとは
同値である．
\item $\mathcal{X}$ が準 DM であることと $\mathcal{X}'$ が準 DM であることとは
同値である．
\item $\mathcal{X}$ が分離的であることと $\mathcal{X}'$ が分離的であることとは
同値である．
\item $\mathcal{X}$ が準分離的であることと $\mathcal{X}'$ が
準分離的であることとは同値である．
\item ここにさらに追加する．
% P12-ERR-0064: preserve the frozen unresolved editorial placeholder literally.
\end{enumerate}
\end{lemma}

\begin{proof}
各場合に問題を対角射に関する問題へ帰着し，厚化の射
$$
(\mathcal{X} \subset \mathcal{X}') \to
\left(\Spec(\mathbf{Z}) \subset \Spec(\mathbf{Z})\right)
$$
に Lemma \ref{stacks-more-morphisms-lemma-thickening-diagonals} を適用する．その際，
Algebraic Stacks, Definition \ref{algebraic-definition-viewed-as} により
$\mathcal{X} \subset \mathcal{X}'$ を $\Spec(\mathbf{Z})$ 上の
代数スタックの厚化とみなす．

\medskip\noindent
場合 (1)．代数スタックが代数空間であるための必要十分条件は，その対角射が
モノ射であることである．Morphisms of Stacks, Lemma
\ref{stacks-morphisms-lemma-hierarchy} を参照せよ
（これは Algebraic Stacks, Proposition
\ref{algebraic-proposition-algebraic-stack-no-automorphisms}
からも直ちに従う）．

\medskip\noindent
場合 (2)．(1) により $\mathcal{X}$ と $\mathcal{X}'$ が代数空間であると
仮定してよく，その後 More on Morphisms of Spaces, Lemma
\ref{spaces-more-morphisms-lemma-thickening-scheme} を用いる．

\medskip\noindent
場合 (3)--(6)．これらの各場合は対角射に関する条件に対応する．
Morphisms of Stacks, Definitions
\ref{stacks-morphisms-definition-separated} および
\ref{stacks-morphisms-definition-absolute-separated} を参照せよ．
\end{proof}
\section{厚化の射}
\label{stacks-more-morphisms-section-morphisms-thickenings}

\noindent
$(f, f') : (\mathcal{X} \subset \mathcal{X}') \to
(\mathcal{Y} \subset \mathcal{Y}')$ が代数スタックの厚化の射であるとき，
射 $f$ の性質はしばしば $f'$ に受け継がれる．いくつかの変種がある．

\begin{lemma}
\label{stacks-more-morphisms-lemma-thicken-property-morphisms}
$(f, f') : (\mathcal{X} \subset \mathcal{X}') \to
(\mathcal{Y} \subset \mathcal{Y}')$
を代数スタックの厚化の射とする．このとき
\begin{enumerate}
\item $f$ がアフィン射であることと $f'$ がアフィン射であることは同値である，
\item $f$ が全射であることと $f'$ が全射であることは同値である，
\item $f$ が準コンパクトであることと $f'$ が準コンパクトであることは同値である，
\item $f$ が普遍閉であることと $f'$ が普遍閉であることは同値である，
\item $f$ が整であることと $f'$ が整であることは同値である，
\item $f$ が普遍単射であることと $f'$ が普遍単射であることは同値である，
\item $f$ が普遍開であることと $f'$ が普遍開であることは同値である，
\item $f$ が準 DMであることと $f'$ が準 DMであることは同値である，
\item $f$ がDMであることと $f'$ がDMであることは同値である，
\item $f$ が（準）分離であることと $f'$ が（準）分離であることは同値である，
\item $f$ が表現可能であることと $f'$ が表現可能であることは同値である，
\item $f$ が代数空間によって表現可能であることと $f'$ が
代数空間によって表現可能であることは同値である，
\item ここにさらに追加する．
\end{enumerate}
\end{lemma}

\begin{proof}
Lemma \ref{stacks-more-morphisms-lemma-thickening} により，射 $\mathcal{X} \to \mathcal{X}'$
および $\mathcal{Y} \to \mathcal{Y}'$ は普遍同相射である．したがって，
$|f| : |\mathcal{X}| \to |\mathcal{Y}|$ に関する任意の条件は，
$|f'| : |\mathcal{X}'| \to |\mathcal{Y}'|$ に関する対応する条件と同値であり，
射 $\mathcal{Z}' \to \mathcal{Y}'$ による任意の基底変換の後にも
同じことが成り立つ．これにより (2)，(3)，(4)，(6)，(7) が成り立つ．

\medskip\noindent
場合 (8)，(9)，(10)，(12) では，$f$ と $f'$ に関する条件を
対角射 $\Delta$ と $\Delta'$ に関する条件へと，補題
\ref{stacks-more-morphisms-lemma-thickening-diagonals} と同様に書き換えることができる．
Morphisms of Stacks, Definition \ref{stacks-morphisms-definition-separated} および
Lemma \ref{stacks-morphisms-lemma-hierarchy} を参照せよ．
したがって，これらの場合はLemma \ref{stacks-more-morphisms-lemma-thickening-diagonals} から従う．

\medskip\noindent
(11) の証明．$f'$ が表現可能ならば $f$ も表現可能である．実際，
スキーム $T$ と射 $T \to \mathcal{Y}$ に対して
$\mathcal{X} \times_\mathcal{Y} T =
\mathcal{X} \times_{\mathcal{X}'} (\mathcal{X}' \times_{\mathcal{Y}'} T)$
であり，$\mathcal{X} \to \mathcal{X}'$ は閉埋め込み（したがって表現可能）である．
逆に，$f$ が表現可能であると仮定し，射 $T' \to \mathcal{Y}'$ をとる．
ここで $T'$ はスキームとする．このとき
$$
\mathcal{X} \times_{\mathcal{Y}}
(\mathcal{Y} \times_{\mathcal{Y}'} T') =
\mathcal{X} \times_{\mathcal{X}'}
(\mathcal{X}' \times_{\mathcal{Y}'} T') \to
\mathcal{X}' \times_{\mathcal{Y}'} T'
$$
は（Lemma \ref{stacks-more-morphisms-lemma-base-change-thickening} により）厚化であり，
その始域はスキームである．したがって補題
\ref{stacks-more-morphisms-lemma-thickening-properties} により，その終域もスキームである．

\medskip\noindent
場合 (1) と (5) では，$f$ または $f'$ のいずれかが述べられた性質をもつならば，
(11) により $f$ と $f'$ はともに表現可能である．この場合，
代数空間 $V'$ と全射かつ滑らかな射
$V' \to \mathcal{Y}'$ をとる．$V = \mathcal{Y} \times_{\mathcal{Y}'} V'$，
$U' = \mathcal{X}' \times_{\mathcal{Y}'} V'$，および
$U = \mathcal{X} \times_{\mathcal{Y}'} V'$ とおく．このとき求める
結果は，代数空間の厚化の射
$(U \subset U') \to (V \subset V')$ に対する対応する結果から，
Properties of Stacks, Lemma
\ref{stacks-properties-lemma-check-property-covering} の原理によって従う．
代数空間の場合の対応する結果については，More on Morphisms of Spaces, Lemma
\ref{spaces-more-morphisms-lemma-thicken-property-morphisms}
を参照せよ．
\end{proof}





\section{代数スタックの無限小変形}
\label{stacks-more-morphisms-section-deform}

\noindent
本節はMore on Morphisms of Spaces, Section
\ref{spaces-more-morphisms-section-deform} に対応するものである．

\begin{lemma}
\label{stacks-more-morphisms-lemma-flatness-morphism-thickenings-fp-over-ft}
次の可換図式を考える．
$$
\xymatrix{
(\mathcal{X} \subset \mathcal{X}') \ar[rr]_{(f, f')} \ar[rd] & &
(\mathcal{Y} \subset \mathcal{Y}') \ar[ld] \\
& (\mathcal{B} \subset \mathcal{B}')
}
$$
この図式は代数スタックの厚化からなるものとする．次を仮定する．
\begin{enumerate}
\item $\mathcal{Y}' \to \mathcal{B}'$ は局所有限型である，
\item $\mathcal{X}' \to \mathcal{B}'$ は
平坦かつ局所有限表示である，
\item $f$ は平坦であり，
\item $\mathcal{X} = \mathcal{B} \times_{\mathcal{B}'} \mathcal{X}'$ かつ
$\mathcal{Y} = \mathcal{B} \times_{\mathcal{B}'} \mathcal{Y}'$ である．
\end{enumerate}
このとき $f'$ は平坦であり，任意の $y' \in |\mathcal{Y}'|$ について，それが
$|f'|$ の像に属するならば，射 $\mathcal{Y}' \to \mathcal{B}'$ は $y'$ において平坦である．
\end{lemma}

\begin{proof}
代数空間 $U'$ と全射かつ滑らかな射
$U' \to \mathcal{B}'$ をとる．
代数空間 $V'$ と全射かつ滑らかな射
$V' \to U' \times_{\mathcal{B}'} \mathcal{Y}'$ をとる．
代数空間 $W'$ と全射かつ滑らかな射
$W' \to V' \times_{\mathcal{Y}'} \mathcal{X}'$ をとる．$U, V, W$ をそれぞれ
$U', V', W'$ の $\mathcal{B} \to \mathcal{B}'$ による基底変換とする．
このとき $f'$ の平坦性は $W' \to V'$ の平坦性と同値であり，
$W \to V$ は平坦であると仮定されている．したがって，次の
代数空間の厚化の図式に代数空間の場合の補題を適用できる．
$$
\xymatrix{
(W \subset W') \ar[rr] \ar[rd] & & (V \subset V') \ar[ld] \\
& (U \subset U')
}
$$
More on Morphisms of Spaces, Lemma
\ref{spaces-more-morphisms-lemma-flatness-morphism-thickenings-fp-over-ft}
を参照せよ．$\mathcal{Y}'/\mathcal{B}'$ が
$|f'|$ の像に属する点において平坦であるという主張も同様に従う．
\end{proof}

\begin{lemma}
\label{stacks-more-morphisms-lemma-deform-property-fp-over-ft}
次の可換図式を考える．
$$
\xymatrix{
(\mathcal{X} \subset \mathcal{X}') \ar[rr]_{(f, f')} \ar[rd] & &
(\mathcal{Y} \subset \mathcal{Y}') \ar[ld] \\
& (\mathcal{B} \subset \mathcal{B}')
}
$$
この図式は代数スタックの厚化からなるものとする．
$\mathcal{Y}' \to \mathcal{B}'$ は局所有限型であり，
$\mathcal{X}' \to \mathcal{B}'$ は平坦かつ局所有限表示であり，
$\mathcal{X} = \mathcal{B} \times_{\mathcal{B}'} \mathcal{X}'$ であり，
$\mathcal{Y} = \mathcal{B} \times_{\mathcal{B}'} \mathcal{Y}'$ と仮定する．このとき
\begin{enumerate}
\item $f$ が平坦であることと $f'$ が平坦であることは同値である，
\label{stacks-more-morphisms-item-flat-fp-over-ft}
\item $f$ が同型であることと $f'$ が同型であることは同値である，
\label{stacks-more-morphisms-item-isomorphism-fp-over-ft}
\item $f$ が開埋め込みであることと $f'$ が開埋め込みであることは同値である，
\label{stacks-more-morphisms-item-open-immersion-fp-over-ft}
\item $f$ がモノ射であることと $f'$ がモノ射であることは同値である，
\label{stacks-more-morphisms-item-monomorphism-fp-over-ft}
\item $f$ が局所準有限であることと $f'$ が局所準有限であることは同値である，
\label{stacks-more-morphisms-item-quasi-finite-fp-over-ft}
\item $f$ がシントミックであることと $f'$ がシントミックであることは同値である，
\label{stacks-more-morphisms-item-syntomic-fp-over-ft}
\item $f$ が滑らかであることと $f'$ が滑らかであることは同値である，
\label{stacks-more-morphisms-item-smooth-fp-over-ft}
\item $f$ が非分岐であることと $f'$ が非分岐であることは同値である，
\label{stacks-more-morphisms-item-unramified-fp-over-ft}
\item $f$ が \'etale であることと $f'$ が \'etale であることは同値である，
\label{stacks-more-morphisms-item-etale-fp-over-ft}
\item $f$ が有限であることと $f'$ が有限であることは同値であり，
\label{stacks-more-morphisms-item-finite-fp-over-ft}
\item ここにさらに追加する．
\end{enumerate}
\end{lemma}

\begin{proof}
場合 (\ref{stacks-more-morphisms-item-flat-fp-over-ft}) では，これは補題
\ref{stacks-more-morphisms-lemma-flatness-morphism-thickenings-fp-over-ft} から従う．

\medskip\noindent
場合
(\ref{stacks-more-morphisms-item-syntomic-fp-over-ft})，(\ref{stacks-more-morphisms-item-smooth-fp-over-ft})
では，Lemma \ref{stacks-more-morphisms-lemma-flatness-morphism-thickenings-fp-over-ft} の証明で
用いた方法によって示すことができる．すなわち，代数空間 $U'$ と全射かつ滑らかな射
$U' \to \mathcal{B}'$ をとる．
代数空間 $V'$ と全射かつ滑らかな射
$V' \to U' \times_{\mathcal{B}'} \mathcal{Y}'$ をとる．
代数空間 $W'$ と全射かつ滑らかな射
$W' \to V' \times_{\mathcal{Y}'} \mathcal{X}'$ をとる．$U, V, W$ をそれぞれ
$U', V', W'$ の $\mathcal{B} \to \mathcal{B}'$ による基底変換とする．
このとき $f$（それぞれ\ $f'$）がこの性質をもつことは，
$W' \to V'$（それぞれ\ $W \to V$）がこの性質をもつことと同値である．
したがって，代数空間の場合の補題を次の
代数空間の厚化の図式に適用できる．
$$
\xymatrix{
(W \subset W') \ar[rr] \ar[rd] & & (V \subset V') \ar[ld] \\
& (U \subset U')
}
$$
More on Morphisms of Spaces, Lemma
\ref{spaces-more-morphisms-lemma-deform-property-fp-over-ft}
を参照せよ．

\medskip\noindent
場合 (\ref{stacks-more-morphisms-item-unramified-fp-over-ft}) と (\ref{stacks-more-morphisms-item-etale-fp-over-ft})
では，まず $f$ または $f'$ に関する仮定から，両方の射
$f$ と $f'$ が代数スタックの DM 射であることがわかる．補題
\ref{stacks-more-morphisms-lemma-thicken-property-morphisms} を参照せよ．次に
代数空間 $U'$ と全射かつ滑らかな射
$U' \to \mathcal{B}'$ をとる．
代数空間 $V'$ と全射かつ滑らかな射
$V' \to U' \times_{\mathcal{B}'} \mathcal{Y}'$ をとる．
代数空間 $W'$ と全射な \'etale(!) 射
$W' \to V' \times_{\mathcal{Y}'} \mathcal{X}'$ をとる．$U, V, W$ をそれぞれ
$U', V', W'$ の $\mathcal{B} \to \mathcal{B}'$ による基底変換とする．
このとき $W \to V \times_\mathcal{Y} \mathcal{X}$ も全射かつ
\'etale である．したがって，$f$（それぞれ\ $f'$）がこの性質をもつことは，
$W' \to V'$（それぞれ\ $W \to V$）がこの性質をもつことと同値である．
したがって，代数空間の場合の補題を次の
代数空間の厚化の図式に適用できる．
$$
\xymatrix{
(W \subset W') \ar[rr] \ar[rd] & & (V \subset V') \ar[ld] \\
& (U \subset U')
}
$$
More on Morphisms of Spaces, Lemma
\ref{spaces-more-morphisms-lemma-deform-property-fp-over-ft}
を参照せよ．

\medskip\noindent
場合 (\ref{stacks-more-morphisms-item-isomorphism-fp-over-ft})，
(\ref{stacks-more-morphisms-item-open-immersion-fp-over-ft})，
(\ref{stacks-more-morphisms-item-monomorphism-fp-over-ft})，
(\ref{stacks-more-morphisms-item-finite-fp-over-ft})
では，まずLemma \ref{stacks-more-morphisms-lemma-thicken-property-morphisms} により，
$f$ と $f'$ は代数空間によって表現可能である．したがって，
代数空間 $U'$ と全射かつ滑らかな射
$U' \to \mathcal{B}'$，
代数空間 $V'$ と全射かつ滑らかな射
$V' \to U' \times_{\mathcal{B}'} \mathcal{Y}'$ を順にとることができ，このとき
$W' = V' \times_{\mathcal{Y}'} \mathcal{X}'$ は代数空間となる．
$U, V, W$ をそれぞれ $U', V', W'$ の
$\mathcal{B} \to \mathcal{B}'$ による基底変換とする．
このときやはり $W = V \times_\mathcal{Y} \mathcal{X}$ である．
そこで，$W' \to V'$ が
同型，\ 開埋め込み，\ モノ射，\ 有限のそれぞれであることと，$W \to V$ が同じ性質をもつこととが
同値であることを示せばよい．Properties of Stacks, Lemma
\ref{stacks-properties-lemma-check-property-covering} を参照せよ．
したがって，上記の代数空間に対する結果を適用すれば結論を得る．

\medskip\noindent
場合 (\ref{stacks-more-morphisms-item-quasi-finite-fp-over-ft}) では，まず
Morphisms of Stacks, Lemma \ref{stacks-morphisms-lemma-finite-type-permanence} により，
$f$ と $f'$ は局所有限型であることに注意する．
一方，射 $f$ が準 DMであることと
$f'$ が準 DMであることは，
Lemma \ref{stacks-more-morphisms-lemma-thicken-property-morphisms} により同値である．
$f$ または $f'$ が局所準有限であるかどうかを判定するために最後に確認すべきことは
（Morphisms of Stacks, Definition
\ref{stacks-morphisms-definition-locally-quasi-finite}），
台となる位相空間に関する条件である．この条件は，
証明の最初の段落の議論により，$f$ について成り立つことと $f'$ について成り立つことが同値である．
\end{proof}
\section{アフィンスキームの持ち上げ}
\label{stacks-more-morphisms-section-lifting-affines}

\noindent
次の実線で示された図式を考える．
$$
\xymatrix{
W \ar[d] \ar@{..>}[r] & W' \ar@{..>}[d] \\
\mathcal{X} \ar[r] & \mathcal{X}'
}
$$
ここで，$\mathcal{X} \subset \mathcal{X}'$ は代数スタックの厚化，
$W$ はアフィンスキームであり，$W \to \mathcal{X}$ は滑らかである．
本節で扱う問題は，正方形がカルテジアンとなり，かつ
$W' \to \mathcal{X}'$ が滑らかとなるような $W'$ と点線の矢印を
見いだせるかどうかである．
一般にはその答えは分からないが，
$\mathcal{X} \subset \mathcal{X}'$ が一次厚化ならば答えが肯定的であることを証明する．

\medskip\noindent
この問題を調べるため，次の圏を導入する．

\begin{remark}[持ち上げの圏]
\label{stacks-more-morphisms-remark-gerbe-of-lifts}
次の図式を考える．
$$
\xymatrix{
W \ar[d]_x \\
\mathcal{X} \ar[r] & \mathcal{X}'
}
$$
ここで，$\mathcal{X} \subset \mathcal{X}'$ は代数スタックの厚化，
$W$ は代数空間であり，$W \to \mathcal{X}$ は滑らかである．
以下のように圏 $\mathcal{C}$ と関手
$$
p : \mathcal{C} \longrightarrow W_{spaces, \etale}
$$
を構成する（記法については Properties of Spaces, Definition
\ref{spaces-properties-definition-spaces-etale-site} を参照）．
$\mathcal{C}$ の対象は系 $(U, U', a, i, x', \alpha)$ であり，これは可換図式
\begin{equation}
\label{stacks-more-morphisms-equation-object}
\vcenter{
\xymatrix{
U \ar[d]_a \ar[r]_i & U' \ar[dd]^{x'} \\
W \ar[d]_x & \\
\mathcal{X} \ar[r] & \mathcal{X}'
}
}
\end{equation}
をなす．その可換性が
$2$-射 $\alpha : x \circ a \to x' \circ i$ によって与えられ，
$U$ と $U'$ が代数空間，$a : U \to W$ が \'etale，
$x' : U' \to \mathcal{X}'$ が滑らかであり，さらに
$U = \mathcal{X} \times_{\mathcal{X}'} U'$ を満たすものとする．
とくに $U \subset U'$ は厚化である．
射
$$
(U, U', a, i, x', \alpha) \to (V, V', b, j, y', \beta)
$$
は $(f, f', \gamma)$ によって与えられる．ここで $f : U \to V$ は
$W$ 上の射，$f' : U' \to V'$ は $U$ への制限が $f$ を与える射，
$\gamma : x' \circ f' \to y'$ は次の図式の右側の三角形の可換性を
与える $2$-射である．
\begin{equation}
\label{stacks-more-morphisms-equation-morphism}
\vcenter{
\xymatrix{
& V \ar[ld]_f \ar[ldd]^b \ar[rr]_j & & V' \ar[ld]_{f'} \ar[lddd]^{y'} \\
U \ar[d]_a \ar[rr]_i & & U' \ar[dd]_{x'} \\
W \ar[d]_x & \\
\mathcal{X} \ar[rr] & & \mathcal{X}'
}
}
\end{equation}
最後に，$\gamma$ が $\alpha$ および $\beta$ と両立することを要求する．
Categories, Sections \ref{categories-section-formal-cat-cat} および
\ref{categories-section-2-categories} の $2$-圏の計算では，これは
$$
\beta = (\gamma \star \text{id}_j) \circ (\alpha \star \text{id}_f)
$$
と書かれる（より簡潔には $\beta = j^*\gamma \circ f^*\alpha$）．
別の言い方をすれば，対象は付加的な性質をもつ可換図式
(\ref{stacks-more-morphisms-equation-object}) であり，射は可換図式
(\ref{stacks-more-morphisms-equation-morphism}) である．後者は圏 $\textit{Spaces}/\mathcal{X}'$
における図式であり，この圏は Properties of Stacks, Remark
\ref{stacks-properties-remark-representable-over} で導入された．
この見方から，$\mathcal{C}$ が圏であり，
$p : \mathcal{C} \to W_{spaces, \etale}$，すなわち
$(U, U', a, i, x', \alpha)$ を $a : U \to W$ へ送る規則が
関手であることは明らかである．
\end{remark}

\begin{lemma}
\label{stacks-more-morphisms-lemma-morphisms-lifts-etale}
任意の射 (\ref{stacks-more-morphisms-equation-morphism}) に対し，写像 $f' : V' \to U'$ は \'etale である．
\end{lemma}

\begin{proof}
実際，$f : V \to U$ は \'etale である．これは $W_{spaces, \etale}$ における射である．
また，Lemma \ref{stacks-more-morphisms-lemma-deform-property-fp-over-ft} を適用できる．なぜなら，
$U' \to \mathcal{X}'$ と $V' \to \mathcal{X}'$ は滑らかで，
$U = \mathcal{X} \times_{\mathcal{X}'} U'$ および
$V = \mathcal{X} \times_{\mathcal{X}'} V'$ が成り立つからである．
\end{proof}

\begin{lemma}
\label{stacks-more-morphisms-lemma-gerbe-of-lifts-fibred}
圏 $p : \mathcal{C} \to W_{spaces, \etale}$ は，
Remark \ref{stacks-more-morphisms-remark-gerbe-of-lifts} で構成したものであり，亜群にファイバー化されている．
\end{lemma}

\begin{proof}
$p$ のファイバー圏は亜群であると主張する．
$(f, f', \gamma')$ が (\ref{stacks-more-morphisms-equation-morphism}) における射であり，
$f : U \to V$ が同型ならば，Lemma \ref{stacks-more-morphisms-lemma-deform-property-fp-over-ft}
により $f'$ も同型である．したがって $(f, f', \gamma')$ は同型である．

\medskip\noindent
$f : V \to U$ を $W_{spaces, \etale}$ における射とし，
$\xi = (U, U', a, i, x', \alpha)$ を $\mathcal{C}$ の $U$ 上の対象とする．
「引き戻し」$f^*\xi$ を $V$ 上に構成しよう．
まず $b = a \circ f$ とおく．射 $f' : V' \to U'$ を取る．この射は \'etale で，
その $V$ への制限が $f$ である（More on Morphisms of Spaces,
Lemma \ref{spaces-more-morphisms-lemma-topological-invariance}）．
対応する厚化を $j : V \to V'$ と記す．
$y' = x' \circ f'$ および $\gamma = \text{id} : x' \circ f' \to y'$ とおく．
さらに
$$
\beta = \alpha \star \text{id}_f :
x \circ b = x \circ a \circ f \to
x' \circ i \circ f = x' \circ f' \circ j = y' \circ j
$$
とおく．$(f, f', \gamma) : (V, V', b, j, y', \beta) \to
(U, U', a, i, x', \alpha)$ が (\ref{stacks-more-morphisms-equation-morphism}) の射であることは
明らかである．このように構成された射 $(f, f', \gamma)$ は強カルテジアンである
（Categories, Definition \ref{categories-definition-cartesian-over-C}）．
詳しい証明は省略するが，本質的な理由は次のとおりである．射
$(g, g', \epsilon) : (Y, Y', c, k, z', \delta) \to (U, U', a, i, x', \alpha)$
が $\mathcal{C}$ において与えられ，$g$ が $g = f \circ h$ と分解するとする．
ここで $h : Y \to V$ である．
このとき一意的な分解 $g' = f' \circ h'$ を得る
（More on Morphisms of Spaces,
Lemma \ref{spaces-more-morphisms-lemma-topological-invariance}）．
その後，必要な $\zeta$ を構成して，
$(h, h', \zeta) :  (Y, Y', c, k, z', \delta) \to
(V, V', b, j, y', \beta)$ が $\mathcal{C}$ の射となり，
$(g, g', \epsilon) = (f, f', \gamma) \circ (h, h', \zeta)$ を満たすようにできる．

\medskip\noindent
したがって $p : \mathcal{C} \to W_\etale$ はファイバー化圏である
（Categories, Definition \ref{categories-definition-fibred-category}）．
これと上で見たファイバー圏が亜群であるという事実を合わせると，
$p : \mathcal{C} \to W_\etale$ は亜群にファイバー化されていると結論される．
これは Categories, Lemma \ref{categories-lemma-fibred-groupoids} による．
\end{proof}

\begin{lemma}
\label{stacks-more-morphisms-lemma-gerbe-of-lifts-stack}
圏 $p : \mathcal{C} \to W_{spaces, \etale}$ は，
Remark \ref{stacks-more-morphisms-remark-gerbe-of-lifts} で構成したものであり，亜群のスタックである．
\end{lemma}

\begin{proof}
Lemma \ref{stacks-more-morphisms-lemma-gerbe-of-lifts-fibred} により，Stacks, Definition
\ref{stacks-definition-stack-in-groupoids} の第1の条件が成り立つ．
慣例どおり対象の降下を検証し，射の降下の検証は読者に委ねる．
そこで，$a : U \to W$ を $W_{spaces, \etale}$ における射，
$\{U_k \to U\}_{k \in K}$ を $W_{spaces, \etale}$ における被覆，
$\xi_k = (U_k, U'_k, a_k, i_k, x'_k, \alpha_k)$ を
$\mathcal{C}$ の $U_k$ 上の対象とし，さらに
コサイクル条件を満たす制限間の射
$$
\varphi_{kk'} = (f_{kk'}, f'_{kk'}, \gamma_{kk'}) :
\xi_k|_{U_k \times_U U_{k'}} \to
\xi_{k'}|_{U_k \times_U U_{k'}}
$$
が与えられているとする．有効性を証明するため，まず被覆を細分してよい．
したがって，各 $U_k$ がスキーム（望むならアフィンスキーム）であると仮定してよい．
次のように書く．
$$
\xi_k|_{U_k \times_U U_{k'}} =
(U_k \times_U U_{k'}, U'_{kk'}, a_{kk'}, x'_{kk'}, \alpha_{kk'})
$$
すると，\'etale 射（Lemma \ref{stacks-more-morphisms-lemma-morphisms-lifts-etale} による）
$s_{kk'} : U'_{kk'} \to U'_k$ を得る．これは射
$\xi_k|_{U_k \times_U U_{k'}} \to \xi_k$ の第2成分であり，
$\mathcal{C}$ における射である．
同様に，\'etale 射 $t_{kk'} : U'_{kk'} \to U'_{k'}$ を，合成
$$
\xi_k|_{U_k \times_U U_{k'}} \xrightarrow{\varphi_{kk'}}
\xi_{k'}|_{U_k \times_U U_{k'}} \to \xi_{k'}
$$
の第2成分を見ることにより得る．
次の射を考える．
$$
j :
\coprod\nolimits_{(k, k') \in K \times K} U'_{kk'}
\xrightarrow{(\coprod s_{kk'}, \coprod t_{kk'})}
(\coprod\nolimits_{k \in K} U'_k) \times (\coprod\nolimits_{k \in K} U'_k)
$$
これは \'etale 同値関係であると主張する．
まず，表示された射の成分 $s, t$ が \'etale であることは既に見た．
射 $j$ を
$(\coprod U_k) \times (\coprod U_k) \to (\coprod U'_k) \times (\coprod U'_k)$
で基底変換すると，それは写像
$$
\coprod\nolimits_{(k, k') \in K \times K} U_k \times_U U_{k'}
\longrightarrow
(\coprod\nolimits_{k \in K} U_k) \times (\coprod\nolimits_{k \in K} U_k)
$$
であるから単射である．したがって More on Morphisms, Lemma
\ref{more-morphisms-lemma-properties-that-extend-over-thickenings} により
$j$ は単射である．最後に，関係 $j$ の対称性は
$\varphi_{kk'}^{-1}$ が $\varphi_{k'k}$ の「反転」であること
（Stacks, Remarks \ref{stacks-remarks-definition-descent-datum} を参照）から従い，
推移性はコサイクル条件から従う（詳細は省略する）．
よって $\coprod U'_k$ の $j$ による商は代数空間 $U'$ である
（Spaces, Theorem \ref{spaces-theorem-presentation}）．
また，上で厚化 $i : U \to U'$ が存在することも既に示されている．
実際，$j$ の $\coprod U_k$ への制限が
$(\coprod U_k) \times_U (\coprod U_k)$ を与えることを見た．
最後に，$1$-射 $x'_k : U'_k \to \mathcal{X}'$ を一時的に
スタック $\mathcal{X}'$ の $U'_k$ 上の対象とみなす．すると，これらには
\'etale 被覆 $\{U'_k \to U'\}$ に関する降下データが備わっていることが分かる．
$\gamma_{kk'}$，すなわち射 $\varphi_{kk'}$ の第3成分がこの降下データを与える．
ここで後者は $\mathcal{C}$ の射である．
$\mathcal{X}'$ はスタックなので，この降下データは有効である．元の記述に戻すと，
$1$-射 $x' : U' \to \mathcal{X}'$ を得る．この射について，合成
$U'_k \to U' \to \mathcal{X}'$ は $x'_k$ への同型を備え，
それらは $\gamma_{kk'}$ と両立する．これは，射
$\alpha_k : x \circ a_k \to x'_k \circ i_k$ が射
$\alpha : x \circ a \to x' \circ i$ へ貼り合わされることを意味する．
したがって $\xi = (U, U', a, i, x', \alpha)$ が求める $U$ 上の対象である．
\end{proof}

\begin{lemma}
\label{stacks-more-morphisms-lemma-etale-local-lifts}
$\mathcal{X} \subset \mathcal{X}'$ を代数スタックの厚化とする．
$W$ を代数空間とし，$W \to \mathcal{X}$ を滑らかな射とする．
\'etale 被覆 $\{W_i \to W\}_{i \in I}$ が存在し，各 $i$ に対して
カルテジアン図式
$$
\xymatrix{
W_i \ar[r] \ar[d] & W_i' \ar[d] \\
\mathcal{X} \ar[r] & \mathcal{X}'
}
$$
であって，$W_i' \to \mathcal{X}'$ が滑らかなものが存在する．
\end{lemma}

\begin{proof}
スキーム $U'$ と全射かつ滑らかな射 $U' \to \mathcal{X}'$ を選ぶ．
通常どおり $U = \mathcal{X} \times_{\mathcal{X}'} U'$ とおく．すると
$U \to \mathcal{X}$ は全射かつ滑らかな射である．したがって基底変換
$$
V = W \times_{\mathcal{X}} U \longrightarrow W
$$
は代数空間の全射かつ滑らかな射である．
Topologies on Spaces, Lemma
\ref{spaces-topologies-lemma-etale-dominates-smooth} により，\'etale 被覆
$\{W_i \to W\}$ であって，$W_i \to W$ が $V \to W$ を経由して
分解するものを見いだせる．$W_i$ をアフィンで被覆すれば
（Properties of Spaces, Lemma
\ref{spaces-properties-lemma-cover-by-union-affines}），
各 $W_i$ がアフィンであると仮定してよい．$W$ を $W_i$ で置き換えることができ，
実際そうすることで，次の段落で論じる状況へ帰着される．

\medskip\noindent
$W$ はアフィンであり，与えられた射 $W \to \mathcal{X}$ は $U$ を経由して
分解すると仮定する．図式は
$$
W \xrightarrow{i} U \to \mathcal{X}
$$
である．$W$ と $U$ は $\mathcal{X}$ 上滑らかなので，$i$ は局所有限型である
（Morphisms of Stacks, Lemma
\ref{stacks-morphisms-lemma-finite-type-permanence}）．
$U$ を $\mathbf{A}^n_U$ で置き換えることにより，$i$ が埋め込みであると
仮定してよい．Morphisms, Lemma
\ref{morphisms-lemma-quasi-affine-finite-type-over-S} を参照されたい．
Morphisms of Stacks, Lemma
\ref{stacks-morphisms-lemma-lci-permanence} により，射 $i$ は局所完全交叉である．
したがって $i$ は Koszul 正則埋め込み（Divisors, Definition
\ref{divisors-definition-regular-immersion} で定義される意味で）である．
これは More on Morphisms, Lemma \ref{more-morphisms-lemma-lci} による．

\medskip\noindent
なおも $W$ をアフィン開被覆で置き換えてよい．
各点 $w \in W$ に対し，アフィン開集合 $U'_w \subset U'$ を選べる．
ここで，対応するアフィン開集合を $U_w \subset U$ とすると，
$w \in i^{-1}(U_w)$ であり，$i^{-1}(U_w) \to U_w$ が
Koszul 正則列
$f_1, \ldots, f_r \in \Gamma(U_w, \mathcal{O}_{U_w})$
によって切り出される閉埋め込みとなるように選ぶ．これは Koszul 正則埋め込みの定義と
Divisors, Lemma \ref{divisors-lemma-regular-ideal-sheaf-scheme} から従う．
$W_w = i^{-1}(U_w)$ とおく．これは $w \in W$ のアフィン開近傍である．
$f'_1, \ldots, f'_r \in \Gamma(U'_w, \mathcal{O}_{U'_w})$ を，
$f_1, \ldots, f_r$ の持ち上げとして選ぶ．
$U_w \to U'_w$ はアフィンスキームの閉埋め込みなので，これは可能である．
$W'_w \subset U'_w$ を $f'_1, \ldots, f'_r$ によって切り出される閉部分スキームとする．
$W'_w \to \mathcal{X}'$ は滑らかであると主張する．
この主張により証明が終わる．実際，構成から
$W_w = \mathcal{X} \times_{\mathcal{X}'} W'_w$ である．

\medskip\noindent
この主張を確かめるには，基底変換
$W'_w \times_{\mathcal{X}'} X' \to X'$ が滑らかであることを，
任意のアフィンスキーム $X'$ であって $\mathcal{X}'$ 上滑らかなものに対して
確かめれば十分である．\'etale 射
$$
Y' \to U'_w \times_{\mathcal{X}'} X'
$$
を，$Y'$ がアフィンとなるように選ぶ．
$U'_w \times_{\mathcal{X}'} X'$ はそのような射の像で被覆されるので，
$Z'$ を，$Y'$ の中で $f'_1, \ldots, f'_r$ によって切り出される
閉部分スキームとする．これが $X'$ 上滑らかであることを示せば十分である．図式は
$$
\xymatrix{
Z' \ar[r] \ar[d] & Y' \ar[d] \\
W'_w \times_{\mathcal{X}'} X' \ar[d] \ar[r] &
U'_w \times_{\mathcal{X}'} X' \ar[d] \ar[r] & X' \\
W'_w = V(f'_1, \ldots, f'_r) \ar[r] & U'_w
}
$$
$X = \mathcal{X} \times_{\mathcal{X}'} X'$，
$Y = X \times_{X'} Y' = \mathcal{X} \times_{\mathcal{X}'} Y'$，および
$Z = Y \times_{Y'} Z' = X \times_{X'} Z' =
\mathcal{X} \times_{\mathcal{X}'} Z'$ とおく．
この列に現れる二つの射
$(Z \subset Z') \to (Y \subset Y') \subset (X \subset X')$ はいずれも，
アフィンスキームの厚化の（カルテジアンな）射であり，
$Z \to X$ と $Y' \to X'$ は滑らかであることが与えられている．
最後に，関数列 $f'_1, \ldots, f'_r$ は
$\Gamma(Y', \mathcal{O}_{Y'})$ の Koszul 正則列へ写る．これは
More on Algebra, Lemma
\ref{more-algebra-lemma-koszul-regular-flat-base-change} による．実際，
$Y' \to U'_w$ は滑らかであり，したがって平坦である．
More on Algebra, Lemma \ref{more-algebra-lemma-cut-by-koszul}
（また，More on Algebra, Lemmas
\ref{more-algebra-lemma-regular-koszul-regular}，
\ref{more-algebra-lemma-koszul-regular-H1-regular}，および
\ref{more-algebra-lemma-H1-regular-quasi-regular} により，Koszul 正則列が
準正則列であるという事実）から，望むとおり $Z' \to X'$ は滑らかであると結論される．
\end{proof}

\begin{lemma}
\label{stacks-more-morphisms-lemma-etale-local-lifts-isomorphic}
$\mathcal{X} \subset \mathcal{X}'$ を代数スタックの厚化とする．
次の可換図式を考える．
$$
\xymatrix{
W'' \ar[d]_{x''} & W \ar[l] \ar[r] \ar[d]_x & W' \ar[d]^{x'} \\
\mathcal{X}' & \mathcal{X} \ar[l] \ar[r] & \mathcal{X}'
}
$$
ここで，二つの正方形はカルテジアンであり，$W', W, W''$ は代数空間，
縦の矢印は滑らかである．このとき，次のものが存在する．
\begin{enumerate}
\item \'etale 被覆 $\{f'_k : W'_k \to W'\}_{k \in K}$，
\item \'etale 射 $f''_k : W'_k \to W''$，および
\item $2$-射 $\gamma_k : x'' \circ f''_k \to x' \circ f'_k$．
\end{enumerate}
これらは (a) $(f'_k)^{-1}(W) = (f''_k)^{-1}(W)$，(b)
$f'_k|_{(f'_k)^{-1}(W)} = f''_k|_{(f''_k)^{-1}(W)}$ を満たし，さらに
(c) $\gamma_k$ を (a) の閉部分スキームへ引き戻したものは，
最初の図式の可換性によって与えられる $2$-射，すなわち $W$ 上のものと一致する．
\end{lemma}

\begin{proof}
与えられた二つの厚化をそれぞれ $i : W \to W'$ および
$i'' : W \to W''$ と記す．
補題の主張にある図式の可換性は，$2$-射
$\delta : x' \circ i' \to x'' \circ i''$ が存在することを意味する．
これが主張の (c) で言及された $2$-射である．代数空間
$$
I' = W' \times_{x', \mathcal{X}', x''} W''
$$
を，射影 $p' : I' \to W'$ および $q' : I' \to W''$ とともに考える．
「普遍」$2$-射 $\gamma : x' \circ p' \to x'' \circ q'$ が存在することに
注意する（これは後で用いる）．$\delta$ の選択により射
$$
\xymatrix{
W \ar[rr]_\delta & & I' \ar[ld]^{p'} \ar[rd]_{q'} \\
& W' & & W''
}
$$
が定まり，合成 $W \to I' \to W'$ と $W \to I' \to W''$ はそれぞれ
$i : W \to W'$ と $i' : W \to W''$ である．
$x''$ は滑らかなので，射 $p' : I' \to W'$ は $x''$ の基底変換として滑らかである．

\medskip\noindent
\'etale 被覆 $\{f'_k : W'_k \to W'\}$ と射 $\delta_k : W'_k \to I'$ であって，
$\delta_k$ の $W_k = (f'_k)^{-1}$ への制限が $\delta \circ f_k$ に等しいものを
見いだせると仮定する．ここで $f_k = f'_k|_{W_k}$ である．図式は
$$
\xymatrix{
W_k \ar[r]^{f_k} \ar[d] & W \ar[r]^\delta & I' \ar[d]^{p'} \\
W'_k \ar[rr]^{f'_k} \ar[rru]^{\delta_k} & & W'
}
$$
である．言い換えれば，与えられた切断 $\delta : W \to I'$（$p'$ の切断）を
$W'$ 上の切断へ延長したいのであり，必要なら $W'$ を \'etale 被覆で置き換えてよい．

\medskip\noindent
これができれば，$f''_k = q' \circ \delta_k$ および
$\gamma_k = \gamma \star \text{id}_{\delta_k}$（より簡潔には
$\gamma_k = \delta_k^*\gamma$）とおける．実際，この時点で残るのは
射 $f''_k$ が \'etale であることを示すことだけである．
構成により，射 $x' \circ p'$ は $2$-同型によって
$x'' \circ q'$ と結ばれる．したがって $x'' \circ f''_k$ は
$2$-同型によって $x' \circ f'_k$ と結ばれる．
ゆえに，合成
$$
W'_k \xrightarrow{f''_k} W'' \xrightarrow{x''} \mathcal{X}'
$$
は $x' \circ f'_k$ が滑らかであることから滑らかである．
$f_k$ は \'etale なので，$f''_k$ は \'etale であると結論される．
これは Lemma \ref{stacks-more-morphisms-lemma-deform-property-fp-over-ft} による．

\medskip\noindent
厚化が一次厚化ならば，任意の \'etale 被覆
$\{W'_k \to W'\}$ であって $W_k'$ がアフィンであるものを選べる．
実際，$p'$ は滑らかなので，無限小持ち上げ判定法
（More on Morphisms of Spaces, Lemma
\ref{spaces-more-morphisms-lemma-smooth-formally-smooth}）により，
$p'$ は形式的に滑らかである．$W_k$ はアフィンであり，
$W_k \to W'_k$ は一次厚化（$\mathcal{X} \to \mathcal{X}'$ の基底変換であるため．
Lemma \ref{stacks-more-morphisms-lemma-base-change-thickening} を参照）なので，求める $\delta_k$ を得る．

\medskip\noindent
一般の場合，被覆と射 $\delta_k$ の存在は More on Morphisms of Spaces, Lemma
\ref{spaces-more-morphisms-lemma-smooth-strong-lift} から従う．
\end{proof}

\begin{lemma}
\label{stacks-more-morphisms-lemma-gerbe-of-lifts}
圏 $p : \mathcal{C} \to W_{spaces, \etale}$ は，
Remark \ref{stacks-more-morphisms-remark-gerbe-of-lifts} で構成したものであり，ジェルブである．
\end{lemma}

\begin{proof}
Lemma \ref{stacks-more-morphisms-lemma-gerbe-of-lifts-stack} で，これは亜群のスタックであることを見た．
したがって，Stacks, Definition \ref{stacks-definition-gerbe} の条件 (2) と (3) を
確かめることだけが残る．条件 (2) は
Lemma \ref{stacks-more-morphisms-lemma-etale-local-lifts} から従う．
条件 (3) は Lemma \ref{stacks-more-morphisms-lemma-etale-local-lifts-isomorphic} から従う．
\end{proof}

\begin{lemma}
\label{stacks-more-morphisms-lemma-gerbe-of-lifts-first-order}
Remark \ref{stacks-more-morphisms-remark-gerbe-of-lifts} において，
$\mathcal{X} \subset \mathcal{X}'$ が一次厚化であると仮定する．このとき，
\begin{enumerate}
\item ジェルブ $p : \mathcal{C} \to W_{spaces, \etale}$
（Remark \ref{stacks-more-morphisms-remark-gerbe-of-lifts} で構成したもの）の対象の自己同型層は可換であり，
\item 群の層 $\mathcal{G}$（Stacks, Lemma
\ref{stacks-lemma-gerbe-abelian-auts} で構成されるもの）は準連接
$\mathcal{O}_W$-加群である．
\end{enumerate}
\end{lemma}

\begin{proof}
二つの主張を同時に証明する．すなわち，対象
$\xi = (U, U', a, i, x', \alpha)$ が与えられたとき，
$\mathit{Aut}(\xi)$ に準連接 $\mathcal{O}_U$-加群の構造を
$U_{spaces, \etale}$ 上で入れ，この構造が引き戻しと両立することを示す．
層の貼り合わせ（Sites, Section \ref{sites-section-glueing-sheaves}），
$\mathcal{G}$ の構成（Stacks, Lemma
\ref{stacks-lemma-gerbe-abelian-auts} の証明では，自己同型層
$\mathit{Aut}(\xi)$ の貼り合わせとして構成される），
および加群が準連接であることは \'etale 被覆へ移った後に確かめれば十分であること
（Properties of Spaces, Lemma
\ref{spaces-properties-lemma-characterize-quasi-coherent}）により，これで十分である．

\medskip\noindent
層 $\mathit{Aut}(\xi)$ を，Lemma
\ref{stacks-more-morphisms-lemma-etale-local-lifts-isomorphic} の証明で用いたものと
同じ方法により記述する．代数空間
$$
I' = U' \times_{x', \mathcal{X}', x'} U'
$$
を，射影 $p' : I' \to U'$ および $q' : I' \to U'$ とともに考える．
$I'$ 上には普遍 $2$-射 $\gamma : x' \circ p' \to x' \circ q'$ がある．
恒等射 $x' \to x'$ は対角射
$$
\xymatrix{
U' \ar[rr]_{\Delta'} & & I' \ar[ld]^{p'} \ar[rd]_{q'} \\
& U' & & U'
}
$$
を定め，合成 $U' \to I' \to U'$ と $U' \to I' \to U'$ は恒等射となる．
$U', I', p', q', \Delta'$ の $\mathcal{X}$ への基底変換を
$U, I, p, q, \Delta$ と記す．$W' \to \mathcal{X}'$ は滑らかなので，
$p' : I' \to U'$ は基底変換として滑らかである．

\medskip\noindent
$\mathit{Aut}(\xi)$ の $U$ 上の切断とは，射 $\delta' : U' \to I'$ であって，
$\delta'|_U = \Delta$ および $p' \circ \delta' = \text{id}_{U'}$ を満たすもののことである．
明示的には，対応する自己同型は
$(\text{id}_U, q' \circ \delta', (\delta')^*\gamma) : \xi \to \xi$
で与えられる．より一般に，$f : V \to U$ が \'etale 射ならば，
厚化 $j : V \to V'$ と \'etale 射 $f' : V' \to U'$ が存在する．
後者の $V$ への制限は $f$ であり，$f^*\xi$ は
$(V, V', a \circ f, j, x' \circ f', f^*\alpha)$ に対応する．
Lemma \ref{stacks-more-morphisms-lemma-gerbe-of-lifts-fibred} の証明を参照されたい．
$\mathit{Aut}(\xi)$ の $V$ 上の切断とは，射 $\delta' : V' \to I'$ であって，
$\delta'|_V = \Delta \circ f$ および
$p' \circ \delta' = f'$\footnote{対応する自己同型は
$(\text{id}_V, h', (\delta')^*\gamma)$ で与えられる．
ここで $h' : V' \to V'$ は，$h'|_V = \text{id}_V$ を満たし，かつ
$$
\xymatrix{
V' \ar[r]_{h'} \ar[rd]_{q' \circ \delta'} & V' \ar[d]^{f'} \\
& U'
}
$$
を可換にする一意的な（同型）射である．$h'$ の一意性と存在は
\'etale サイトの位相的不変性による．More on Morphisms of Spaces,
Theorem \ref{spaces-more-morphisms-theorem-topological-invariance} を参照されたい．
読者は，代わりに射
$\delta'' : V' \to V' \times_{\mathcal{X}'} V'$ であって，
$\delta'' \circ j = \Delta_{V'/\mathcal{X}'}$ および
$\text{pr}_1 \circ \delta'' = \text{id}_{V'}$ を満たすものを考えるべきだと
感じるかもしれない．これでもよい．実際，
$V' \times_{\mathcal{X}'} V' \to I'$ は \'etale なので，同じ位相的不変性により，
$\delta''$ を
$\delta' = (V' \times_{\mathcal{X}'} V' \to I') \circ \delta''$ へ送る写像は，
二つの射の集合の間の全単射である．} を満たすもののことである．

\medskip\noindent
以上から，集合の層としての $\mathit{Aut}(\xi)$ は，
More on Morphisms of Spaces, Remark
\ref{spaces-more-morphisms-remark-another-special-case} において，厚化
$(U \subset U')$ と $(I \subset I')$ に対して，$(U \subset U')$ 上の
$\text{id}_{U'}$ と $p'$ を介して定義された層と一致する．
対角射 $\Delta'$ はこの層の切断であり，More on Morphisms of Spaces, Lemma
\ref{spaces-more-morphisms-lemma-action-sheaf} を用いてこの切断に作用させることで，同型
\begin{equation}
\label{stacks-more-morphisms-equation-isomorphism}
\SheafHom_{\mathcal{O}_U}(\Delta^*\Omega_{I/U}, \mathcal{C}_{U/U'})
\longrightarrow
\mathit{Aut}(\xi)
\end{equation}
を $U_{spaces, \etale}$ 上に得る．なお，次の三点を確かめる必要がある．
\begin{enumerate}
\item (\ref{stacks-more-morphisms-equation-isomorphism}) の構成が \'etale 局所化と可換すること，
\item $\SheafHom_{\mathcal{O}_U}(\Delta^*\Omega_{I/U}, \mathcal{C}_{U/U'})$
が $U$ 上の準連接加群であること，
\item $\mathit{Aut}(\xi)$ における合成が，この準連接加群の切断の加法に
対応すること．
\end{enumerate}
これらを順に確かめる．

\medskip\noindent
(1) を見るには，$f : V \to U$ が \'etale ならば，
(\ref{stacks-more-morphisms-equation-isomorphism})，すなわち $\xi$ を用いて $U$ 上で
構成した同型が，写像 (\ref{stacks-more-morphisms-equation-isomorphism})
$$
\SheafHom_{\mathcal{O}_V}(
\Delta_V^*\Omega_{V \times_\mathcal{X} V/V}, \mathcal{C}_{V/V'}) \to
\mathit{Aut}(\xi|_V)
$$
へ制限されることを示さなければならない．ここで後者は，$\xi|_V$ を用いて
$V$ 上で構成した $V_{spaces, \etale}$ 上の写像である．
これは上の脚注の議論と
More on Morphisms of Spaces, Lemma
\ref{spaces-more-morphisms-lemma-action-by-derivations-etale-localization} から従う．

\medskip\noindent
(2) の証明．$p'$ は滑らかなので，射 $I \to U$ は滑らかであり，
したがって相対微分加群 $\Omega_{I/U}$ は有限局所自由である
（More on Morphisms of Spaces, Lemma
\ref{spaces-more-morphisms-lemma-smooth-omega-finite-locally-free}）．
一方，$\mathcal{C}_{U/U'}$ は準連接である
（More on Morphisms of Spaces, Definition
\ref{spaces-more-morphisms-definition-conormal-sheaf}）．
Properties of Spaces, Lemma
\ref{spaces-properties-lemma-properties-quasi-coherent} により結論を得る．

\medskip\noindent
(3) の証明．射 $c' : I' \times_{p', U', q'} I' \to I'$ が存在し，
$(U', I', p', q', c')$ は恒等射 $\Delta'$ をもつ代数空間の亜群となる．
たとえば Algebraic Stacks, Lemma
\ref{algebraic-lemma-map-space-into-stack} を参照されたい．
$\mathit{Aut}(\xi)$ における合成は，射 $c'$ によって次のように誘導される．
二つの射
$$
\delta'_1, \delta'_2 : U' \longrightarrow I'
$$
が，上で述べたように $\mathit{Aut}(\xi)$ の $U$ 上の切断に対応するとする．
言い換えれば，$\delta'_i|U = \Delta_U$ および
$p' \circ \delta'_i = \text{id}_{U'}$ が成り立つ．このとき
$\mathit{Aut}(\xi)$ における合成は
$$
\delta'_1 \circ \delta'_2 = c'(\delta'_1 \circ q' \circ \delta'_2, \delta'_2)
$$
である．詳しい確認は省略する\footnote{読者には，$\delta'_1$ に
$q' \circ \delta'_2$ を前合成して初めて，$U'$ に値をもつ
ファイバー積 $I' \times_{p', U', q'} I'$ の良定義された点が得られることが
直ちに分かるであろう．}．
したがって More on Groupoids in Spaces, Section
\ref{spaces-more-groupoids-section-groupoid-sections} で述べられた状況にあり，
求める結果は More on Groupoids in Spaces, Lemma
\ref{spaces-more-groupoids-lemma-composition-is-addition} から従う．
\end{proof}

\begin{proposition}[Emerton]
\label{stacks-more-morphisms-proposition-affine-smooth-lift-to-first-order}
\begin{reference}
2016年4月27日付の Matthew Emerton の電子メール．
\end{reference}
$\mathcal{X} \subset \mathcal{X}'$ を代数スタックの一次厚化とする．
$W$ をアフィンスキームとし，$W \to \mathcal{X}$ を滑らかな射とする．
このとき，カルテジアン図式
$$
\xymatrix{
W \ar[d] \ar[r] & W' \ar[d] \\
\mathcal{X} \ar[r] & \mathcal{X}'
}
$$
であって，$W' \to \mathcal{X}'$ が滑らかで $W'$ がアフィンであるものが存在する．
\end{proposition}

\begin{proof}
圏 $p : \mathcal{C} \to W_{spaces, \etale}$ を考える．これは
Remark \ref{stacks-more-morphisms-remark-gerbe-of-lifts} で導入したものである．
命題の主張は，$\mathcal{C}$ の対象で $W$ 上にあるものが存在するということである．
実際，そのような対象 $(W, W', a, i, y', \alpha)$ があれば，
$W = \mathcal{X} \times_{\mathcal{X}'} W'$ である．
したがって $W \to W'$ は代数空間の厚化であるから，
More on Morphisms of Spaces, Lemma
\ref{spaces-more-morphisms-lemma-thickening-scheme}
および More on Morphisms, Lemma
\ref{more-morphisms-lemma-thickening-affine-scheme} により $W'$ はアフィンである．

\medskip\noindent
Lemma \ref{stacks-more-morphisms-lemma-gerbe-of-lifts} によれば，$\mathcal{C}$ は
$W_{spaces, \etale}$ 上のジェルブである．これは，\'etale 局所的に解を
見いだすことができ，かつそれらの局所解が \'etale 局所的に同型であることを意味する．
この部分では，この厚化が一次厚化であるという仮定は必要ない．
Lemma \ref{stacks-more-morphisms-lemma-gerbe-of-lifts-first-order} により，このジェルブの対象の
自己同型層は可換であり，互いに貼り合わさって準連接加群
$\mathcal{G}$ をなす．これは $W_{spaces, \etale}$ 上の加群である．
Cohomology on Sites, Lemma \ref{sites-cohomology-lemma-existence} の
条件 (1) と (2) を検証し，$\mathcal{C}$ の対象で $W$ 上にあるものの存在を結論する．
条件 (1) は成り立つ．すなわち，\'etale 被覆 $\{W_i \to W\}$ であって，
各 $W_i$ がアフィンであるものは，被覆全体の族の中で共終である．
そのような被覆に対して，$W_i$ と $W_i \times_W W_j$ はアフィンであり，
$H^1(W_i, \mathcal{G})$ と $H^1(W_i \times_W W_j, \mathcal{G})$ は零である．
実際，たとえば Cohomology of Spaces, Proposition
\ref{spaces-cohomology-proposition-vanishing} により，アフィン代数空間上の
準連接加群のコホモロジーは零である．
最後に，条件 (2) は $H^2(W, \mathcal{G}) = 0$ となることである．
ここで $\mathcal{G}$ は上の準連接層であり，この消滅も Cohomology of Spaces,
Proposition \ref{spaces-cohomology-proposition-vanishing} から従う．
これで証明が完了した．
\end{proof}
\section{無限小変形}
\label{stacks-more-morphisms-section-inf}

\noindent
引き続き，Artin's Axioms, Section \ref{artin-section-inf} の議論を行う．

\begin{lemma}
\label{stacks-more-morphisms-lemma-inf-quasi-coherent}
$\mathcal{X}$ をスキーム $S$ 上の代数スタックとする．
$\mathcal{I}_\mathcal{X} \to \mathcal{X}$ は局所有限表示であると仮定する．
$A \to B$ を平坦な $S$-代数準同型とする．
$x$ を $\mathcal{X}$ の $A$ 上の対象とし，$y = x|_B$ とおく．
このとき $\text{Inf}_x(M) \otimes_A B = \text{Inf}_y(M \otimes_A B)$ である．
\end{lemma}

\begin{proof}
$\text{Inf}_x(M)$ は，$x$ の $A[M]$ への自明な変形の自己同型のうち，
$x$ の $A$ 上の恒等自己同型を誘導するもの全体の集合であった．
この自明な変形は，$x$ の $\Spec(A[M])$ への，
$\Spec(A[M]) \to \Spec(A)$ による引き戻しである．
$G \to \Spec(A)$ を $x$ の自己同型群代数空間とする
（これは $\mathcal{X}$ が代数空間なので存在する）．
$e : \Spec(A) \to G$ を単位元とする．
More on Morphisms of Spaces, Section
\ref{spaces-more-morphisms-section-action-by-derivations}
の議論から
$$
\text{Inf}_x(M) = \Hom_A(e^*\Omega_{G/A}, M)
$$
を得る．同様に
$$
\text{Inf}_y(M \otimes_A B) = \Hom_B(e_B^*\Omega_{G_B/B}, M \otimes_A B)
$$
である．仮定により $G \to \Spec(A)$ は局所有限表示なので，
$\Omega_{G/A}$ は局所有限表示である．More on Morphisms of Spaces, Lemma
\ref{spaces-more-morphisms-lemma-finite-presentation-differentials}
を参照されたい．したがって $e^*\Omega_{G/A}$ は有限表示 $A$-加群である．
さらに，$\Omega_{G_B/B}$ は $\Omega_{G/A}$ の引き戻しである．
これは More on Morphisms of Spaces, Lemma
\ref{spaces-more-morphisms-lemma-base-change-differentials} による．
ゆえに $e_B^*\Omega_{G_B/B} = e^*\Omega_{G/A} \otimes_A B$ である．
More on Algebra, Lemma
\ref{more-algebra-lemma-pseudo-coherence-and-base-change-ext}
から結論が従う．
\end{proof}

\begin{lemma}
\label{stacks-more-morphisms-lemma-sheaf-of-infinitesimal-lifts}
$\mathcal{X}$ を基礎スキーム $S$ 上の代数スタックとする．
$\mathcal{I}_\mathcal{X} \to \mathcal{X}$ は局所有限表示であると仮定する．
$(A' \to A, x)$ を変形状況とする．このとき関手
$$
F : B' \longmapsto
\{\text{次の写像の持ち上げ：}x|_{B' \otimes_{A'} A}\text{，持ち上げ先：} B'\}/\text{同型}
$$
はサイト $(\textit{Aff}/\Spec(A'))_{fppf}$ 上の層である．
これは Topologies, Definition \ref{topologies-definition-big-small-fppf}
のサイトである．
\end{lemma}

\begin{proof}
$\{T'_i \to T'\}_{i = 1, \ldots n}$ を $A'$ 上のアフィンスキームの
標準 fppf 被覆とする．$T' = \Spec(B')$ と書く．通常どおり
$$
T'_{i_0 \ldots i_p} =
T'_{i_0} \times_{T'} \ldots \times_{T'} T'_{i_p} = \Spec(B'_{i_0 \ldots i_p})
$$
と書く．ここで環は適切なテンソル積である．
$B = B' \otimes_{A'} A$ および
$B_{i_0 \ldots i_p} = B'_{i_0 \ldots i_p} \otimes_{A'} A$ とおく．
$y = x|_B$ および $y_{i_0 \ldots i_p} = x|_{B_{i_0 \ldots i_p}}$ と書く．
$\gamma_i \in F(B'_i)$ とし，$\gamma_{i_0}$ と $\gamma_{i_1}$ の
$F(B'_{i_0i_1})$ における像が同じであるとする．
唯一の $\gamma \in F(B')$ で，その像が $\gamma_i$
（$F(B'_i)$ の元）であるものを見つけなければならない．

\medskip\noindent
$y'_i$ という $\textit{Lift}(y_i, B'_i)$ の実際の対象を，
同型類 $\gamma_i$ の中から選ぶ．同型
$\varphi_{i_0i_1} : y'_{i_0}|_{B'_{i_0i_1}} \to y'_{i_1}|_{B'_{i_0i_1}}$
を圏 $\textit{Lift}(y_{i_0i_1}, B'_{i_0i_1})$ の中で選ぶ．
写像 $\varphi_{i_0i_1}$ がコサイクル条件を満たすなら，
所望の対象 $\gamma$ が得られる．これは $\mathcal{X}$ が
fppf 位相に関するスタックだからである．コサイクル条件とは，合成
$$
y'_{i_0}|_{B'_{i_0i_1i_2}}
\xrightarrow{\varphi_{i_0i_1}|_{B'_{i_0i_1i_2}}}
y'_{i_1}|_{B'_{i_0i_1i_2}}
\xrightarrow{\varphi_{i_1i_2}|_{B'_{i_0i_1i_2}}}
y'_{i_2}|_{B'_{i_0i_1i_2}}
\xrightarrow{\varphi_{i_2i_0}|_{B'_{i_0i_1i_2}}}
y'_{i_0}|_{B'_{i_0i_1i_2}}
$$
が恒等射であるという条件である．そうでなければ，これらの写像は元
$$
\delta_{i_0i_1i_2} \in
\text{Inf}_{y_{i_0i_1i_2}}(J_{i_0i_1i_2}) =
\text{Inf}_y(J) \otimes_B B_{i_0i_1i_2}
$$
を与える．ここで $J = \Ker(B' \to B)$ および
$J_{i_0 \ldots i_p} = \Ker(B'_{i_0 \ldots i_p} \to B_{i_0 \ldots i_p})$ である．
表示した式の等号は，Lemma \ref{stacks-more-morphisms-lemma-inf-quasi-coherent} を
$B' \to B'_{i_0 \ldots i_p}$ と $y$ および
$y_{i_0 \ldots i_p}$ に適用することから従う．
写像 $B' \to B'_{i_0 \ldots i_p}$ の平坦性は，さらに
$J_{i_0 \ldots i_p} = J \otimes_{B'} B'_{i_0 \ldots i_p}$ も保証する．
計算（省略する）により，$\delta_{i_0i_1i_2}$ は次の {\v C}ech 複体における
$2$-コサイクルをなす．
$$
\prod \text{Inf}_y(J) \otimes_B B_{i_0} \to
\prod \text{Inf}_y(J) \otimes_B B_{i_0i_1} \to
\prod \text{Inf}_y(J) \otimes_B B_{i_0i_1i_2} \to \ldots
$$
Descent, Lemma
\ref{descent-lemma-standard-covering-Cech-quasi-coherent}
により，この複体の正次数のコホモロジーは消滅し，
$H^0 = \text{Inf}_y(J)$ である．
$\text{Inf}_{y_{i_0i_1}}(J_{i_0i_1})$ は射に作用するので
（Artin's Axioms, Remark \ref{artin-remark-automorphisms}），
$\varphi_{i_0i_1}$ の選び方を変更して
$\delta_{i_0i_1i_2} = 0$ の場合に帰着できる．

\medskip\noindent
一意性．さらに，$\gamma$ であって，すべての $i$ について $\gamma_i$ に
制限されるものが高々一つであることを示さなければならない．
$y', z'$ を $\textit{Lift}(y, B')$ の対象とし，
$\psi_i : y'|_{B'_i} \to z'|_{B'_i}$ を
$\textit{Lift}(y_i, B'_i)$ における同型とする．すると
$$
\psi_{i_1}^{-1} \circ \psi_{i_0} \in
\text{Inf}_{y_{i_0i_1}}(J_{i_0i_1}) =
\text{Inf}_y(J) \otimes_B B_{i_0i_1}
$$
を考えることができる．先ほどと同様に論じると，
$y'$ と $z'$ の間の $B'$ 上の同型の存在に対する障害は，上に表示した
{\v C}ech 複体の $H^1$ の元であるが，これは零である．
\end{proof}

\begin{lemma}
\label{stacks-more-morphisms-lemma-T-quasi-coherent}
$\mathcal{X}$ をスキーム $S$ 上の代数スタックとし，その構造射
$\mathcal{X} \to S$ は局所有限表示であるとする．
$A \to B$ を平坦な $S$-代数準同型とする．
$x$ を $\mathcal{X}$ の $A$ 上の対象とする．
このとき $T_x(M) \otimes_A B = T_y(M \otimes_A B)$ である．
\end{lemma}

\begin{proof}
スキーム $U$ と全射かつ滑らかな射 $U \to \mathcal{X}$ を選ぶ．
まず補題を，$x$ が $U$ に持ち上がる場合へ帰着させる．
$T_x(M)$ は $x$ の $A[M]$ への持ち上げの同型類全体の集合であった．
したがって
Lemma \ref{stacks-more-morphisms-lemma-sheaf-of-infinitesimal-lifts}\footnote{この補題は適用できる．
$\Delta : \mathcal{X} \to \mathcal{X} \times_S \mathcal{X}$ は
Morphisms of Stacks, Lemma
\ref{stacks-morphisms-lemma-finite-presentation-permanence}
と $\mathcal{X} \to S$ が局所有限表示であるという仮定により，
局所有限表示である．したがって $\mathcal{I}_\mathcal{X} \to \mathcal{X}$ は
$\Delta$ の基底変換として局所有限表示である．}
により，対応
$$
A_1 \mapsto T_{x|_{A_1}}(M \otimes_A A_1)
$$
は $\Spec(A)$ の小 \'etale サイト上の層である．テンソル積は，
$A[M] \to A_1[M \otimes_A A_1]$ を平坦な環準同型とするために必要である．
忠実平坦な \'etale 環準同型 $A \to A_1$ であって，
$x|_{A_1}$ が射 $u_1 : \Spec(A_1) \to U$ に持ち上がるものを選ぶことができる．
例えば Sheaves on Stacks, Lemma
\ref{stacks-sheaves-lemma-surjective-flat-locally-finite-presentation}
を参照されたい．$A_2 = A_1 \otimes_A A_1$ と書き，
$B_1 = B \otimes_A A_1$ および $B_2 = B \otimes_A A_2$ とおく．
次の図式を考える．
$$
\xymatrix{
0 \ar[r] &
T_y(M \otimes_A B) \ar[r] &
T_{y|_{B_1}}(M \otimes_A B_1) \ar[r] &
T_{y|_{B_2}}(M \otimes_A B_2) \\
0 \ar[r] &
T_x(M) \ar[r] \ar[u] &
T_{x|_{A_1}}(M \otimes_A A_1) \ar[r] \ar[u] &
T_{x|_{A_2}}(M \otimes_A A_2) \ar[u]
}
$$
層条件により各行は完全である．
$M \otimes_A B_i = (M \otimes_A A_i) \otimes_{A_i} B_i$ である．
したがって中央と右の縦矢印について結果を証明すれば，
求める結果が従う．これで次の段落で議論する場合に帰着された．

\medskip\noindent
$x$ は射 $u : \Spec(A) \to U$ の像であると仮定する．
$T_u(M) \to T_x(M)$ は全射であることに注意する．これは
$U \to \mathcal{X}$ が滑らかで代数空間によって表現可能だからである．
Criteria for Representability, Lemma
\ref{criteria-lemma-representable-by-spaces-formally-smooth}
（説明についてはその直前の議論を参照されたい）および
More on Morphisms of Spaces, Lemma
\ref{spaces-more-morphisms-lemma-smooth-formally-smooth}
を参照されたい．$R = U \times_\mathcal{X} U$ とおく．
代数空間の亜群 $(U, R, s, t, c, e, i)$ が得られ，
$\mathcal{X} = [U/R]$ となることを思い出そう．
Artin's Axioms, Lemma \ref{artin-lemma-ses-inf-and-T}
により，完全列
$$
T_{e \circ u}(M) \to T_u(M) \oplus T_u(M) \to T_x(M) \to 0
$$
を得る．ここで右端の零は上で示した．$B$ への基底変換についても
同様の列が成り立つ．したがって，$T_u(M)$ と $T_{e \circ u}(M)$
について補題の結果を証明できれば，求める結果が従う．
これで次の段落で議論する場合に帰着された．

\medskip\noindent
$\mathcal{X} = X$ は $S$ 上局所有限表示な代数空間であると仮定する．
このとき
$$
T_x(M) = \Hom_A(x^*\Omega_{X/S}, M)
$$
である．これは More on Morphisms of Spaces, Section
\ref{spaces-more-morphisms-section-action-by-derivations}
の議論による．同様に
$$
T_y(M \otimes_A B) = \Hom_B(y^*\Omega_{X/S}, M \otimes_A B)
$$
である．$X \to S$ は局所有限表示なので，$\Omega_{X/S}$ は
局所有限表示である．More on Morphisms of Spaces, Lemma
\ref{spaces-more-morphisms-lemma-finite-presentation-differentials}
を参照されたい．したがって $x^*\Omega_{X/S}$ は有限表示 $A$-加群である．
明らかに $y^*\Omega_{X/S} = x^*\Omega_{X/S} \otimes_A B$ である．
More on Algebra, Lemma
\ref{more-algebra-lemma-pseudo-coherence-and-base-change-ext}
から結論が従う．
\end{proof}

\begin{lemma}
\label{stacks-more-morphisms-lemma-local-lift-enough}
$\mathcal{X}$ をスキーム $S$ 上の代数スタックとし，その構造射
$\mathcal{X} \to S$ は局所有限表示であるとする．
$(A' \to A, x)$ を変形状況とする．忠実平坦な有限表示 $A'$-代数
$B'$ と，対象 $y'$ であって $\mathcal{X}$ の $B'$ 上にあり，
$x|_{B' \otimes_{A'} A}$ を持ち上げるものが存在するなら，対象
$x'$ であって，$A'$ 上で $x$ を持ち上げるものが存在する．
\end{lemma}

\begin{proof}
$I = \Ker(A' \to A)$ とする．$B'_1 = B' \otimes_{A'} B'$ および
$B'_2 = B' \otimes_{A'} B' \otimes_{A'} B'$ とおく．
$J = IB'$，$J_1 = IB'_1$，$J_2 = IB'_2$ および
$B = B'/J$，$B_1 = B'_1/J_1$，$B_2 = B'_2/J_2$ とおく．
$y = x|_B$，$y_1 = x|_{B_1}$，$y_2 = x|_{B_2}$ とおく．
$F$ を Lemma \ref{stacks-more-morphisms-lemma-sheaf-of-infinitesimal-lifts} の fppf 層とする
（これは適用できる．Lemma \ref{stacks-more-morphisms-lemma-T-quasi-coherent} の証明中の
脚注を参照されたい）．したがって等化子図式
$$
\xymatrix{
F(A') \ar[r] &
F(B') \ar@<1ex>[r] \ar@<-1ex>[r] &
F(B'_1)
}
$$
を得る．一方，$F(B') = \text{Lift}(y, B')$，
$F(B'_1) = \text{Lift}(y_1, B'_1)$，$F(B'_2) = \text{Lift}(y_2, B'_2)$
である．ここでは Artin's Axioms, Section \ref{artin-section-inf}
の用語を用いた．これらの集合は空でなく，それぞれ $T_y(J)$，
$T_{y_1}(J_1)$，$T_{y_2}(J_2)$ の（標準的な）主等質空間である．
Artin's Axioms, Lemma \ref{artin-lemma-properties-lift-RS-star}
を参照されたい．したがって，$y'$ の二つの像の $F(B'_1)$ における差は元
$$
\delta_1 \in T_{y_1}(J_1) = T_x(I) \otimes_A B_1
$$
である．表示した式の等号は，Lemma \ref{stacks-more-morphisms-lemma-T-quasi-coherent} を
$A' \to B'_1$ と $x$ および $y_1$ に適用することから従う．
$A' \to B'_1$ の平坦性は，さらに
$J_1 = I \otimes_{A'} B'_1$ も保証する．$B'$ と $B'_2$ についても
同様の等式が成り立つ．計算（省略する）により，$\delta_1$ は
次の {\v C}ech 複体における $1$-コサイクルをなす．
$$
T_x(I) \otimes_A B \to
T_x(I) \otimes_A B_1 \to
T_x(I) \otimes_A B_2 \to \ldots
$$
Descent, Lemma
\ref{descent-lemma-standard-covering-Cech-quasi-coherent}
により，この複体の正次数のコホモロジーは消滅し，$H^0 = T_x(I)$ である．
したがって，$T_x(I) \otimes_A B = T_y(J)$ の中に元を選ぶことができ，
その境界は $\delta_1$ である．$y'$ を，この元を作用させた結果である
新たな選択 $y'$ に置き換えると，$\delta_1 = 0$ となる．したがって $y'$ は
二つの写像 $F(B') \to F(B'_1)$ の下で同じ元に写り，
層条件により $F(A')$ の元 o を得る．
\end{proof}











\section{形式的に滑らかな射}
\label{stacks-more-morphisms-section-formally-smooth}

\noindent
本節では，代数スタックの形式的に滑らかな射
$\mathcal{X} \to \mathcal{Y}$ の概念を導入する．このような射は，
$T$ がアフィンならば $\mathcal{X}$ の $T$-値点が $T$ の無限小厚化へ
持ち上がるという性質によって特徴づけられる．主要な結果は，形式的に
滑らかで局所有限表示な射は滑らかであるというものである．
Lemma \ref{stacks-more-morphisms-lemma-smooth-formally-smooth} を参照されたい．
この判定法はヤコビ判定法よりも使いやすいことが多い．

\begin{definition}
\label{stacks-more-morphisms-definition-formally-smooth}
代数スタックの射 $f : \mathcal{X} \to \mathcal{Y}$ が，
亜群にファイバー化された圏における $1$-射として対象について
形式的に滑らかであるとき，{\it 形式的に滑らか} であるという．
この意味は Criteria for Representability, Section
\ref{criteria-section-formally-smooth} で説明されている．
\end{definition}

\noindent
この定義の条件を，現在用いている言葉に翻訳しよう
（Properties of Stacks, Section \ref{stacks-properties-section-conventions}
を参照されたい）．$f : \mathcal{X} \to \mathcal{Y}$ を代数スタックの射とする．
次の $2$-可換な実線の図式を考える．
\begin{equation}
\label{stacks-more-morphisms-equation-diagram}
\vcenter{
\xymatrix{
T \ar[r]_-x \ar[d]_i & \mathcal{X} \ar[d]^f \\
T' \ar[r]^-y \ar@{..>}[ru] & \mathcal{Y}
}
}
\end{equation}
ここで $i : T \to T'$ はアフィンスキームの一次厚化である．
$$
\gamma : y \circ i \longrightarrow f \circ x
$$
を，この図式の $2$-可換性を与える $2$-射とする．
（記法については Categories, Sections \ref{categories-section-formal-cat-cat}
および \ref{categories-section-2-categories} を参照されたい．）
(\ref{stacks-more-morphisms-equation-diagram}) と $\gamma$ が与えられたとき，
{\it 点線矢印} とは，三つ組 $(x', \alpha, \beta)$ であって，射
$x' : T' \to \mathcal{X}$ と $2$-射
$\alpha : x' \circ i \to x$，$\beta : y \to f \circ x'$ からなり，
$\gamma = (\text{id}_f \star \alpha) \circ (\beta \star \text{id}_i)$
を満たすものをいう．言い換えれば，図式
$$
\xymatrix{
& f \circ x' \circ i \ar[rd]^{\text{id}_f \star \alpha} \\
y \circ i \ar[ru]^{\beta \star \text{id}_i} \ar[rr]^\gamma & &
f \circ x
}
$$
が可換であるということである．{\it 点線矢印の射}
$(x'_1, \alpha_1, \beta_1) \to (x'_2, \alpha_2, \beta_2)$ とは，
$2$-射 $\theta : x'_1 \to x'_2$ であって，
$\alpha_1 = \alpha_2 \circ (\theta \star \text{id}_i)$ および
$\beta_2 = (\text{id}_f \star \theta) \circ \beta_1$ を満たすものをいう．

\medskip\noindent
以上で記述した点線矢印の圏は，Categories, Definition
\ref{categories-definition-dotted-arrows} の特別な場合である．

\begin{lemma}
\label{stacks-more-morphisms-lemma-reformulate-formal-smoothness}
代数スタックの射 $f : \mathcal{X} \to \mathcal{Y}$ が形式的に滑らか
（Definition \ref{stacks-more-morphisms-definition-formally-smooth}）であるための必要十分条件は，
すべての図式 (\ref{stacks-more-morphisms-equation-diagram}) と $\gamma$ に対して，
点線矢印の圏が空でないことである．
\end{lemma}

\begin{proof}
異なる言葉の間の翻訳は省略する．
\end{proof}

\begin{lemma}
\label{stacks-more-morphisms-lemma-base-change-formal-smoothness}
代数スタックの形式的に滑らかな射を代数スタックの任意の射で
基底変換したものは形式的に滑らかである．
\end{lemma}

\begin{proof}
Categories, Lemma \ref{categories-lemma-cat-dotted-arrows-base-change}
と定義から従う．
\end{proof}

\begin{lemma}
\label{stacks-more-morphisms-lemma-composition-formal-smoothness}
代数スタックの形式的に滑らかな射の合成は形式的に滑らかである．
\end{lemma}

\begin{proof}
Categories, Lemma \ref{categories-lemma-cat-dotted-arrows-composition}
と定義から従う．
\end{proof}

\begin{lemma}
\label{stacks-more-morphisms-lemma-formal-smoothness-representable}
$f : \mathcal{X} \to \mathcal{Y}$ を代数空間によって表現可能な
代数スタックの射とする．このとき次の条件は同値である．
\begin{enumerate}
\item $f$ は形式的に滑らかである．
\item 任意のスキーム $T$ と射 $T \to \mathcal{Y}$ に対して，射
$\mathcal{X} \times_\mathcal{Y} T \to T$ は代数空間の射として
形式的に滑らかである．
\end{enumerate}
\end{lemma}

\begin{proof}
Categories, Lemma \ref{categories-lemma-cat-dotted-arrows-base-change}
と定義から従う．
\end{proof}

\begin{lemma}
\label{stacks-more-morphisms-lemma-lift-to-smooth}
$T \to T'$ をアフィンスキームの一次厚化とする．
$\mathcal{X}'$ を $T'$ 上の代数スタックとし，その構造射
$\mathcal{X}' \to T'$ は滑らかであるとする．
$x : T \to \mathcal{X}'$ を $T'$ 上の射とする．
このとき，射 $x' : T' \to \mathcal{X}'$ であって，
$T'$ 上にあり，$x'|_T = x$ を満たすものが存在する．
\end{lemma}

\begin{proof}
Lemma \ref{stacks-more-morphisms-lemma-local-lift-enough} の結果を適用できる．
したがって，全射かつ滑らかな射 $W' \to T'$ であって，$W'$ がアフィンであり，
$x|_{T \times_{W'} T'}$ が $W'$ に持ち上がるものを構成すれば十分である．
（すでに確立した代数空間に対する類似の結果を用いて，
この事実の証明を各自で見つけることを読者に勧める．）
スキーム $U'$ と全射かつ滑らかな射 $U' \to \mathcal{X}'$ を選ぶ．
$U' \to T'$ は滑らかであり，射影
$T \times_{\mathcal{X}'} U' \to T$ は全射かつ滑らかであることに注意する．
アフィンスキーム $W$ と \'etale 射
$W \to T \times_{\mathcal{X}'} U'$ であって，$W \to T$ が
全射となるものを選ぶ．すると $W \to T$ は
アフィンスキーム間の滑らかな射である．$W$ をアフィン主開部分の
非交和で置き換えることにより，アフィンスキームの滑らかな射
$W' \to T'$ が存在し，$W = T \times_{T'} W'$ を満たすと仮定してよい．
Algebra, Lemma \ref{algebra-lemma-lift-smooth} を参照されたい．
More on Morphisms of Spaces, Lemma
\ref{spaces-more-morphisms-lemma-smooth-formally-smooth}
により，射 $W' \to U'$ であって，$T'$ 上にあり，与えられた射
$W \to U'$ を持ち上げるものを見つけることができる．これで証明が完了する．
\end{proof}

\noindent
次の補題が本節の主要な結果である．これは Limits of Stacks, Proposition
\ref{stacks-limits-proposition-characterize-locally-finite-presentation}
と組み合わせると，代数スタックの射
$f : \mathcal{X} \to \mathcal{Y}$ が滑らかであるかどうかを，
亜群のスタックの $1$-射 $\mathcal{X} \to \mathcal{Y}$ の
``単純な'' 性質によって判定できることを意味する．

\begin{lemma}[無限小持ち上げ判定法]
\label{stacks-more-morphisms-lemma-smooth-formally-smooth}
$f : \mathcal{X} \to \mathcal{Y}$ を代数スタックの射とする．
次の条件は同値である．
\begin{enumerate}
\item 射 $f$ は滑らかである．
\item 射 $f$ は局所有限表示かつ形式的に滑らかである．
\end{enumerate}
\end{lemma}

\begin{proof}
$f$ は滑らかであると仮定する．すると Morphisms of Stacks, Lemma
\ref{stacks-morphisms-lemma-smooth-locally-finite-presentation}
により $f$ は局所有限表示である．したがって，図式
(\ref{stacks-more-morphisms-equation-diagram}) と $\gamma : y \circ i \to f \circ x$ が
与えられたときに点線矢印を見つければ十分である
（Lemma \ref{stacks-more-morphisms-lemma-reformulate-formal-smoothness} を参照されたい）．
ファイバー積を作ると
$$
\xymatrix{
T \ar[d] \ar[r] &
T' \times_\mathcal{Y} \mathcal{X} \ar[d] \ar[r] &
\mathcal{X} \ar[d] \\
T' \ar[r] & T' \ar[r] & \mathcal{Y}
}
$$
を得る．したがって，左の正方形に点線矢印を見つければ十分である．
$T' \times_\mathcal{Y} \mathcal{X} \to T'$ は滑らかなので
（Morphisms of Stacks, Lemma \ref{stacks-morphisms-lemma-base-change-smooth}），
左の正方形における点線矢印の存在は
Lemma \ref{stacks-more-morphisms-lemma-lift-to-smooth} によって保証される．

\medskip\noindent
逆に，$f$ は局所有限表示かつ形式的に滑らかであると仮定する．
スキーム $U$ と全射かつ滑らかな射 $U \to \mathcal{X}$ を選ぶ．
すると $a : U \to \mathcal{X}$ および $b : U \to \mathcal{Y}$ は
代数空間によって表現可能かつ局所有限表示である
（Morphisms of Stacks, Lemma
\ref{stacks-morphisms-lemma-composition-finite-presentation}
と，上で見た滑らかな射が局所有限表示であるという事実を用いる）．
Algebraic Stacks, Lemma
\ref{algebraic-lemma-representable-transformations-property-implication}
の一般原理を適用する．その際，More on Morphisms of Spaces, Lemma
\ref{spaces-more-morphisms-lemma-smooth-formally-smooth}
の同値性を入力として用いると同時に，
Criteria for Representability, Lemma
\ref{criteria-lemma-representable-by-spaces-formally-smooth}
の翻訳も用いる．まずこれを $a$ に適用すると，$a$ は対象について
形式的に滑らかであることが分かる．次に，仮定により $f$ は対象について
形式的に滑らかであること
（Lemma \ref{stacks-more-morphisms-lemma-reformulate-formal-smoothness} を参照されたい）と
Criteria for Representability, Lemma
\ref{criteria-lemma-composition-formally-smooth}
を用いると，$b = f \circ a$ は対象について形式的に滑らかであることが分かる．
そこでこの原理をもう一度適用すると，$b$ は滑らかであると結論できる．
代数スタックの射が滑らかであることの定義により，これは $f$ が
滑らかであることを意味し，証明が完了する．
\end{proof}







\section{ブローアップと平坦性}
\label{stacks-more-morphisms-section-blowup-flat}

\noindent
本節では，代数空間上の代数スタックについて，
More on Morphisms of Spaces, Sections
\ref{spaces-more-morphisms-section-blowup-flat} および
\ref{spaces-more-morphisms-section-applications-flattening-by-blowing-up}
から何が導けるかを手短に論じる．

\begin{lemma}
\label{stacks-more-morphisms-lemma-flatten-stack}
$f : \mathcal{X} \to Y$ を代数スタックから代数空間への射とする．
$V \subset Y$ を開部分空間とする．次を仮定する．
\begin{enumerate}
\item $Y$ は準コンパクトかつ準分離である．
\item $f$ は有限型かつ準分離である．
\item $V$ は準コンパクトである．
\item $\mathcal{X}_V$ は $V$ 上平坦かつ局所有限表示である．
\end{enumerate}
このとき，$V$-許容ブローアップ $Y' \to Y$ と閉部分スタック
$\mathcal{X}' \subset \mathcal{X}_{Y'}$ であって，
$\mathcal{X}'_V = \mathcal{X}_V$ であり，かつ
$\mathcal{X}' \to Y'$ が平坦かつ有限表示となるものが存在する．
\end{lemma}

\begin{proof}
$\mathcal{X}$ は準コンパクトであることに注意する．
アフィンスキーム $U$ と全射かつ滑らかな射
$U \to \mathcal{X}$ を選ぶ．
$R = U \times_\mathcal{X} U$ とおくと，代数空間の亜群
$(U, R, s, t, c)$ で $Y$ 上にあり，$\mathcal{X} = [U/R]$ を満たすものが得られる
（Algebraic Stacks, Lemma \ref{algebraic-lemma-stack-presentation}）．
More on Morphisms of Spaces, Lemma
\ref{spaces-more-morphisms-lemma-flat-after-blowing-up}
を $U \to Y$ と開部分空間 $V \subset Y$ に適用できる．
したがって，$V$-許容ブローアップ $Y' \to Y$ であって，狭義変換
$U' \subset U_{Y'}$ が $Y'$ 上平坦かつ有限表示となるものが得られる．
$R' \subset R_{Y'}$ を $R$ の狭義変換とする．
$s$ と $t$ は滑らか（特に平坦）なので，
Divisors on Spaces, Lemma
\ref{spaces-divisors-lemma-strict-transform-flat}
により，カルテジアン図式
$$
\vcenter{
\xymatrix{
R' \ar[r] \ar[d] & R_{Y'} \ar[d]^{s_{Y'}} \\
U' \ar[r] & U_{Y'}
}
}
\quad\text{および}\quad
\vcenter{
\xymatrix{
R' \ar[r] \ar[d] & R_{Y'} \ar[d]^{t_{Y'}} \\
U' \ar[r] & U_{Y'}
}
}
$$
を得る．言い換えれば，$U'$ は $R_{Y'}$-不変な
$U_{Y'}$ の閉部分空間である．したがって，$U'$ は閉部分スタック
$\mathcal{X}' \subset \mathcal{X}_{Y'}$ を定める．これは
Properties of Stacks, Lemma
\ref{stacks-properties-lemma-substacks-presentation} による．
射 $\mathcal{X}' \to Y'$ は平坦かつ局所有限表示である．これは
$U' \to Y'$ について成り立つからである．一方，すでに
$\mathcal{X}' \to Y'$ は準コンパクトかつ準分離であることが分かっている
（$f$ に関する仮定と，閉埋め込みについてこの性質が成り立つことによる）．
これで証明が完了する．
\end{proof}
\section{代数スタックに対する Chow の補題}
\label{stacks-more-morphisms-section-chow}

\noindent
本節では，代数スタックに対する Chow の補題を論じる．

\begin{lemma}
\label{stacks-more-morphisms-lemma-finite-cover-factor}
$Y$ を準コンパクトかつ準分離的な代数空間とする．
$V \subset Y$ を準コンパクト開部分空間とする．$f : \mathcal{X} \to V$
を全射，平坦かつ局所有限表示とする．
このとき，有限全射 $g : Y' \to Y$ であって，
$V' = g^{-1}(V) \to Y$ が Zariski 局所的に $f$ を経由して分解するものが存在する．
% P12-ERR-0092: preserve the frozen source codomain Y in the conclusion.
\end{lemma}

\begin{proof}
まず $Y$ がスキームである場合を証明する．
スキーム $U$ と全射平滑射
$U \to \mathcal{X}$ を選べる．このとき $\{U \to V\}$ はスキームの fppf 被覆である．
More on Morphisms, Lemma \ref{more-morphisms-lemma-dominate-fppf-finite}
により，有限全射
$V' \to V$ であって，$V' \to V$ が Zariski 局所的に
$U$ を経由して分解するものが存在する．
More on Morphisms, Lemma
\ref{more-morphisms-lemma-extend-finite-surjective-morphisms}
により，その $V$ への制限が所望の $V' \to V$ となる有限全射
$Y' \to Y$ を見いだせる．

\medskip\noindent
$Y$ が代数空間である場合には，$Y'$ がスキームであるような有限全射
$Y' \to Y$ によってまず有限基底変換を行えば，補題が成り立つことが分かる．
Limits of Spaces, Proposition
\ref{spaces-limits-proposition-there-is-a-scheme-finite-over}
を参照せよ．
\end{proof}

\begin{lemma}
\label{stacks-more-morphisms-lemma-make-section}
$f : \mathcal{X} \to Y$ を代数スタックから代数空間への射とする．
$V \subset Y$ を開部分空間とする．次を仮定する．
\begin{enumerate}
\item $f$ は分離的かつ有限型である．
\item $Y$ は準コンパクトかつ準分離的である．
\item $V$ は準コンパクトである．
\item $\mathcal{X}_V$ は $V$ 上の gerbe である．
\end{enumerate}
このとき，可換図式
$$
\xymatrix{
\overline{Z} \ar[rd]_{\overline{g}} &
Z \ar[l]^j \ar[d]_g \ar[r]_h & \mathcal{X} \ar[ld]^f \\
& Y
}
$$
であって，$j$ が開埋め込み，$\overline{g}$ と $h$ が固有であり，
$|V|$ が $|g|$ の像に含まれるものが存在する．
\end{lemma}

\begin{proof}
可換図式
$$
\xymatrix{
\mathcal{X}' \ar[d]_{f'} \ar[r] & \mathcal{X} \ar[d]^f \\
Y' \ar[r] & Y
}
$$
と準コンパクト開部分空間 $V' \subset Y'$ があり，
$Y' \to Y$ は代数空間の固有射，
$\mathcal{X}' \to \mathcal{X}$ は代数スタックの固有射，
$V' \subset Y'$ は $V$ の上へ全射的に写り，かつ
$\mathcal{X}'_{V'}$ は $V'$ 上の gerbe であると仮定する．
このとき，対 $(f' : \mathcal{X}' \to Y', V')$ について補題を証明すれば十分である．
詳細の一部は省略する．

\medskip\noindent
証明の全体的な方針を述べる．上の注意を繰り返し適用して，
$f$ の像が開であり，かつこの開集合上で $f$ が切断をもつ場合へ帰着する．
各段階は直接的であるが，段階数がかなり多いため証明は少し込み入っている．

\medskip\noindent
Limits of Spaces, Proposition
\ref{spaces-limits-proposition-there-is-a-scheme-finite-over}
を用いて，$Y$ がスキームである場合へ帰着する．
（$Y'$ がスキームであるような有限全射 $Y' \to Y$ をとり，
$\mathcal{X}' = \mathcal{X}_{Y'}$ と置いて証明冒頭の注意を適用する．）

\medskip\noindent
Lemma \ref{stacks-more-morphisms-lemma-flatten-stack}
（および gerbe が平坦かつ局所有限表示であることを示すための
Morphisms of Stacks, Lemma \ref{stacks-morphisms-lemma-gerbe-fppf}）
を用いて，$f$ が平坦かつ有限表示である場合へ帰着する．

\medskip\noindent
$f$ は平坦かつ局所有限表示なので，$|f|$ の像は開部分空間
$W \subset Y$ である．$f$ が有限型で $Y$ が準コンパクトであるため
$\mathcal{X}$ は準コンパクトであり，したがって $W$ は準コンパクトである．
Lemma \ref{stacks-more-morphisms-lemma-finite-cover-factor} により，有限全射 $g : Y' \to Y$
であって，$g^{-1}(W) \to Y$ が Zariski 局所的に
$\mathcal{X} \to Y$ を経由して分解するものを選べる．
$Y$ を $Y'$ で，$\mathcal{X}$ を $\mathcal{X} \times_Y Y'$ で
置き換えると，次段落の状況へ帰着する．

\medskip\noindent
$n \geq 0$，準コンパクト開部分空間
$W_i \subset Y$（$i = 1, \ldots, n$），および射
$x_i : W_i \to \mathcal{X}$ が存在し，次を満たすと仮定する．
(a) $f \circ x_i = \text{id}_{W_i}$，
(b) $W = \bigcup_{i = 1, \ldots, n} W_i$ は $V$ を含み，
(c) $W$ は $|f|$ の像である．
$n$ に関する帰納法を用いる．基底の場合は $n = 0$ である．このとき
$V = \emptyset$ であり，$\overline{Z} = \emptyset$ ととれる．
$n > 0$ のとき，$i = 1, \ldots, n$ に対し，台位相空間が
$Y \setminus W_i$ である被約閉部分スキーム $Y_i$ を考える．
有限射
$$
Y' = Y \amalg \coprod\nolimits_{i = 1, \ldots, n} Y_i \longrightarrow Y
$$
と準コンパクト開部分空間
$$
V' = (W_1 \cap \ldots \cap W_n \cap V) \amalg
\coprod_{i = 1, \ldots, n} (V \cap Y_i).
$$
を考える．証明冒頭の注意により，対
$$
(\mathcal{X} \to Y, W_1 \cap \ldots \cap W_n \cap V)
\quad\text{かつ}\quad
(\mathcal{X} \times_Y Y_i \to Y_i, V \cap Y_i),\quad
i = 1, \ldots, n
$$
について補題を証明できれば結論が従う．ここで集合論的等式
$V = (W_1 \cap \ldots \cap W_n \cap V) \cup
\bigcup\nolimits_{i = 1, \ldots n} (V \cap Y_i)$
を用いる．上の第2種の対には帰納法の仮定が適用できる．
したがって次段落の状況へ帰着する．

\medskip\noindent
$n \geq 0$，準コンパクト開部分空間
$W_i \subset Y$（$i = 1, \ldots, n$），および射
$x_i : W_i \to \mathcal{X}$ が存在し，次を満たすと仮定する．
(a) $f \circ x_i = \text{id}_{W_i}$，
(b) $W = \bigcup_{i = 1, \ldots, n} W_i$ は $V$ を含み，
(c) $W$ は $|f|$ の像であり，
(d) $V \subset W_1 \cap \ldots \cap W_n$ である．
射
$$
T_{ij} = \mathit{Isom}_\mathcal{X}(x_i|_{W_i \cap W_j \cap V},
x_j|_{W_i \cap W_j \cap V}) \longrightarrow W_i \cap W_j \cap V
$$
は全射，平坦かつ局所有限表示である
（Morphisms of Stacks, Lemma \ref{stacks-morphisms-lemma-gerbe-isom-fppf}）．
各準コンパクト開部分空間 $W_i \cap W_j \cap V$ と射
$T_{ij} \to W_i \cap W_j \cap V$ に Lemma \ref{stacks-more-morphisms-lemma-finite-cover-factor}
を適用して，有限全射 $Y'_{ij} \to Y$ を得る．
$Y$ をすべての $Y'_{ij}$ の $Y$ 上のファイバー積で置き換えると，
次段落の状況へ帰着する．

\medskip\noindent
$n \geq 0$，準コンパクト開部分空間
$W_i \subset Y$（$i = 1, \ldots, n$），および射
$x_i : W_i \to \mathcal{X}$ が存在し，次を満たすと仮定する．
(a) $f \circ x_i = \text{id}_{W_i}$，
(b) $W = \bigcup_{i = 1, \ldots, n} W_i$ は $V$ を含み，
(c) $W$ は $|f|$ の像であり，
(d) $V \subset W_1 \cap \ldots \cap W_n$ であり，
(e) $x_i$ と $x_j$ は $W_i \cap W_j \cap V$ 上 Zariski 局所的に同型である．
$y \in V$ を任意にとる．準コンパクト開近傍
$y \in V_y \subset V$ であって，対 $(\mathcal{X} \to Y, V_y)$ について
補題が成り立つものを見いだせたとする．その解を
$\overline{Z}_y, Z_y, \overline{g}_y, g_y, h_y$ とする．
$V$ は準コンパクトなので，$V = V_{y_1} \cup \ldots \cup V_{y_m}$
となる有限個の $y_1, \ldots, y_m$ を選べる．このとき
$$
\overline{Z} = \coprod \overline{Z}_{y_j},\quad
Z = \coprod Z_{y_j},\quad
\overline{g} = \coprod \overline{g}_{y_j},\quad
g = \coprod g_{y_j},\quad
h = \coprod h_{y_j}
$$
と置けば補題の解となる．$y$ が与えられれば，条件 (e) により
準コンパクト開近傍 $y \in V_y \subset V$ と同型
$\varphi_i : x_1|_{V_y} \to x_i|_{V_y}$（$i = 2, \ldots, n$）
を選べる．$\varphi_{ij} = \varphi_j \circ \varphi_i^{-1}$ と置く．
これにより次段落の状況へ至る．

\medskip\noindent
$n \geq 0$，準コンパクト開部分空間
$W_i \subset Y$（$i = 1, \ldots, n$），および射
$x_i : W_i \to \mathcal{X}$ が存在し，次を満たすと仮定する．
(a) $f \circ x_i = \text{id}_{W_i}$，
(b) $W = \bigcup_{i = 1, \ldots, n} W_i$ は $V$ を含み，
(c) $W$ は $|f|$ の像であり，
(d) $V \subset W_1 \cap \ldots \cap W_n$ であり，
(f) 同型 $\varphi_{ij} : x_i|_V \to x_j|_V$ が存在して
$\varphi_{jk} \circ \varphi_{ij} = \varphi_{ik}$ を満たす．
射
$$
I_{ij} = \mathit{Isom}_\mathcal{X}(x_i|_{W_i \cap W_j},
x_j|_{W_i \cap W_j}) \longrightarrow W_i \cap W_j
$$
は，$f$ が分離的なので固有である
（Morphisms of Stacks, Lemma
\ref{stacks-morphisms-lemma-separated-implies-isom}）．
$\varphi_{ij}$ は $V$ 上で $I_{ij} \to W_i \cap W_j$ の切断
$V \to I_{ij}$ を定めることに注意せよ．
More on Morphisms of Spaces, Lemma
\ref{spaces-more-morphisms-lemma-get-section-after-blowup}
により，$s_{ij}$ が $p_{ij}^{-1}(W_i \cap W_j)$ へ延長されるような
$V$-許容 blowup $p_{ij} : Y_{ij} \to Y$ を選べる．
% P12-ERR-0093: preserve the frozen undefined s_{ij}; the introduced section is phi_{ij}.
$Y$ をすべての $Y_{ij}$ の $Y$ 上のファイバー積で置き換えると，
次段落の状況に至る．

\medskip\noindent
$n \geq 0$，準コンパクト開部分空間
$W_i \subset Y$（$i = 1, \ldots, n$），および射
$x_i : W_i \to \mathcal{X}$ が存在し，次を満たすと仮定する．
(a) $f \circ x_i = \text{id}_{W_i}$，
(b) $W = \bigcup_{i = 1, \ldots, n} W_i$ は $V$ を含み，
(c) $W$ は $|f|$ の像であり，
(d) $V \subset W_1 \cap \ldots \cap W_n$ であり，
(g) 同型
$\varphi_{ij} : x_i|_{W_i \cap W_j} \to x_j|_{W_i \cap W_j}$
が存在して
$$
\varphi_{jk}|_V \circ \varphi_{ij}|_V = \varphi_{ik}|_V.
$$
を満たす．必要なら $Y$ をさらに別の $V$-許容 blowup で置き換えることにより，
$V$ が $Y$ において稠密かつスキーム論的に稠密であり，したがって $V$ を含む
$Y$ の任意の開部分空間でも同様であると仮定してよい．このように置き換えた後，
Morphisms of Spaces, Lemma
\ref{spaces-morphisms-lemma-equality-of-morphisms}
と $I_{ik} \to W_i \cap W_j$ が固有（したがって分離的）であることにより，
% P12-ERR-0094: preserve the frozen mismatched base W_i cap W_j for I_{ik}.
$$
\varphi_{jk}|_{W_i \cap W_j \cap W_k} \circ
\varphi_{ij}|_{W_i \cap W_j \cap W_k} =
\varphi_{ik}|_{W_i \cap W_j \cap W_k}
$$
を得る．もちろん，これは $(x_i, \varphi_{ij})$ が降下データであることを意味し，
% P12-ERR-0095: render the intended term while preserving all mathematics of the misspelled source prose.
$\mathcal{X}$ はスタックなので，各 $W_i$ 上で $x_i$ と一致する射
$x : W \to \mathcal{X}$ を得る．$x$ は分離射
$\mathcal{X} \to W$ の切断であるから，$x$ は固有である
（Morphisms of Stacks, Lemma \ref{stacks-morphisms-lemma-section-immersion}）．
したがって，$\overline{Z} = Y$，$Z = W$，
$\overline{g} = \text{id}_Y$，$g = \text{id}_W$，$h = x$
とすれば補題が成り立つ．
\end{proof}

\begin{theorem}[Chow の補題]
\label{stacks-more-morphisms-theorem-chow-finite-type}
\begin{reference}
これは Ofer Gabber による結果である．
\cite[Theorem 1.1]{olsson_proper} を参照せよ．
\end{reference}
$f : \mathcal{X} \to Y$ を代数スタックから代数空間への射とする．
次を仮定する．
\begin{enumerate}
\item $Y$ は準コンパクトかつ準分離的である．
\item $f$ は分離的かつ有限型である．
% P12-ERR-0096: Japanese supplies the conjunction omitted in the frozen English source.
\end{enumerate}
このとき，可換図式
$$
\xymatrix{
\mathcal{X} \ar[rd] & X \ar[l] \ar[d] \ar[r] & \overline{X} \ar[ld] \\
& Y
}
$$
であって，$X \to \mathcal{X}$ が固有全射，
$X \to \overline{X}$ が開埋め込み，かつ
$\overline{X} \to Y$ が代数空間の固有射であるものが存在する．
% P12-ERR-0097: Japanese supplies the article-free grammatical relation.
\end{theorem}

\begin{proof}
大まかな着想は，$\mathcal{X}$ が gerbe であるような稠密開部分をもつこと
（Morphisms of Stacks, Proposition
\ref{stacks-morphisms-proposition-open-stratum}）を用い，
Lemma \ref{stacks-more-morphisms-lemma-make-section} を適用することである．
これがそのままでは機能しない理由は，この開部分が準コンパクトとは限らず，
技術的な問題が生じるからである．そこでまず，Noether 的な場合への
（標準的な）帰着を行う．

\medskip\noindent
まず閉埋め込み $\mathcal{X} \to \mathcal{X}'$ であって，
$\mathcal{X}'$ が $Y$ 上分離的かつ有限型の代数スタックであるものを選ぶ．
Limits of Stacks, Lemma
\ref{stacks-limits-lemma-separated-closed-in-finite-presentation}
を参照せよ．$\mathcal{X}'$ について定理を証明すれば十分であることは明らかなので，
$\mathcal{X} \to Y$ は分離的かつ有限表示であると仮定してよい．

\medskip\noindent
$\mathcal{X} \to Y$ が分離的かつ有限表示であると仮定する．
Limits of Spaces, Proposition \ref{spaces-limits-proposition-approximate}
により，$Y = \lim Y_i$ を，アフィン遷移射をもつ Noether 代数空間系の
有向極限として書ける．Limits of Stacks, Lemma
\ref{stacks-limits-lemma-descend-a-stack-down} により，ある $i$ と，
代数スタックから $Y_i$ への有限表示射 $\mathcal{X}_i \to Y_i$ が存在して
$\mathcal{X} = Y \times_{Y_i} \mathcal{X}_i$ となる．
$i$ を大きくすれば $\mathcal{X}_i \to Y_i$ が分離的であると仮定できる．
Limits of Stacks, Lemma
\ref{stacks-limits-lemma-eventually-separated} を参照せよ．
したがって $\mathcal{X}_i \to Y_i$ について定理を証明すれば十分であり，
次段落の場合へ帰着する．

\medskip\noindent
$Y$ が Noether 的であると仮定する．$\mathcal{X}$ をその被約化で置き換えてよい
（Properties of Stacks, Definition
\ref{stacks-properties-definition-reduced-induced-stack}）．
これにより次段落の場合へ帰着する．

\medskip\noindent
$Y$ は Noether 的で，$\mathcal{X}$ は被約であると仮定する．
$\mathcal{X} \to Y$ は分離的で $Y$ は準分離的なので，
$\mathcal{X}$ は代数スタックとして準分離的である．したがって慣性射
$\mathcal{I}_\mathcal{X} \to \mathcal{X}$ は準コンパクトである．
よって Morphisms of Stacks, Proposition
\ref{stacks-morphisms-proposition-open-stratum} により，
gerbe である稠密開部分スタック $\mathcal{V} \subset \mathcal{X}$ が存在する．
$\mathcal{V} \to V$ を，$\mathcal{V}$ を代数空間 $V$ 上の gerbe として表す射とする．
$\mathcal{V} \to V$ の構成については Morphisms of Stacks, Lemma
\ref{stacks-morphisms-lemma-gerbe-over-iso-classes} を参照せよ．
特にこの構成から，射 $\mathcal{V} \to Y$ は
$\mathcal{V} \to V \to Y$ と分解することが分かる．図式は
$$
\xymatrix{
\mathcal{V} \ar[r] \ar[d] & \mathcal{X} \ar[d] \\
V \ar[r] & Y
}
$$
である．射 $\mathcal{V} \to V$ は全射，平坦かつ有限表示であり
（Morphisms of Stacks, Lemma \ref{stacks-morphisms-lemma-gerbe-fppf}），
$\mathcal{V} \to Y$ は局所有限表示なので，$V \to Y$ は局所有限表示である
（Morphisms of Stacks, Lemma
\ref{stacks-morphisms-lemma-flat-finite-presentation-permanence}）．
$\mathcal{V} \to V$ は普遍同相射であることに注意せよ
（Morphisms of Stacks, Lemma
\ref{stacks-morphisms-lemma-gerbe-bijection-points}）．
$\mathcal{V}$ は準コンパクトなので（Morphisms of Stacks, Lemma
\ref{stacks-morphisms-lemma-locally-closed-in-noetherian} を参照），
$V$ も準コンパクトである．最後に，$\mathcal{V} \to Y$ が分離的なので，
$\mathcal{V} \to V \to Y$ に Morphisms of Stacks, Lemma
\ref{stacks-morphisms-lemma-check-separated-on-ui-cover}
を適用すれば $V \to Y$ も分離的である
（その仮定が満たされることは既に確認した）．

\medskip\noindent
以上により，Limits of Spaces, Lemma
\ref{spaces-limits-lemma-embedding-into-affine-over-qs} の仮定が射
$V \to Y$ に適用できる．したがって，稠密開部分空間 $V' \subset V$ と
$Y$ 上の埋め込み $V' \to \mathbf{P}^n_Y$ を選べる．
$V$ を $V'$ で，$\mathcal{V}$ を $\mathcal{V}$ における $V'$ の逆像で
置き換えてよいことは明らかである（上で見たように
$|\mathcal{V}| = |V|$ であることを思い出せ）．
よって図式
$$
\xymatrix{
\mathcal{V} \ar[rr] \ar[d] & & \mathcal{X} \ar[d] \\
V \ar[r] & \mathbf{P}^n_Y \ar[r] & Y
}
$$
があり，矢印 $V \to \mathbf{P}^n_Y$ が埋め込みであると仮定してよい．
$\mathcal{X}'$ を射
$$
j : \mathcal{V} \longrightarrow \mathbf{P}^n_Y \times_Y \mathcal{X}
$$
のスキーム論的像とし，$Y'$ を射 $V \to \mathbf{P}^n_Y$ の
スキーム論的像とする．可換図式
$$
\xymatrix{
\mathcal{V} \ar[r] \ar[d] &
\mathcal{X}' \ar[r] \ar[d] &
\mathbf{P}^n_Y \times_Y \mathcal{X} \ar[d] \ar[r] &
\mathcal{X} \ar[d] \\
V \ar[r] &
Y' \ar[r] &
\mathbf{P}^n_Y \ar[r] &
Y
}
$$
を得る
（Morphisms of Stacks, Lemma \ref{stacks-morphisms-lemma-factor-factor} を参照）．
$\mathcal{V} = V \times_{Y'} \mathcal{X}'$ であり，かつ
射 $\mathcal{X}' \to Y'$ と開部分空間 $V \subset Y'$ に
Lemma \ref{stacks-more-morphisms-lemma-make-section} を適用できると主張する．
この主張が正しければ，
$$
\xymatrix{
\overline{X} \ar[rd]_{\overline{g}} &
X \ar[l] \ar[d]_g \ar[r]_h & \mathcal{X}' \ar[ld]^f \\
& Y'
}
$$
であって，$X \to \overline{X}$ が開埋め込み，
$\overline{g}$ と $h$ が固有で，$|V|$ が $|g|$ の像に含まれるものを得る．
このとき合成 $X \to \mathcal{X}' \to \mathcal{X}$ は固有射の合成として固有であり，
その像は $|\mathcal{V}|$ を含むから，この合成は全射である．
また，$\overline{X} \to Y' \to Y$ も固有射の合成として固有である．

\medskip\noindent
最後にこの主張を証明する．$\mathcal{X}' \to Y'$ は分離的かつ有限型であり，
$Y'$ は準コンパクトかつ準分離的であり，$V$ は準コンパクトであることに注意せよ
（詳細の完全な確認は省略する）．次に，Morphisms of Stacks, Lemma
\ref{stacks-morphisms-lemma-scheme-theoretic-image-of-partial-section}
により，$b : \mathcal{X}' \to \mathcal{X}$ は $\mathcal{V}$ 上で同型である．
特に $\mathcal{V}$ は $\mathcal{X}'$ の開部分スタックと同一視される．
射 $j$ は準コンパクトである（始域が準コンパクトで終域が準分離的だからである）．
したがって Morphisms of Stacks, Lemma
\ref{stacks-morphisms-lemma-existence-plus-flat-base-change}
により，$j$ のスキーム論的像の形成は平坦基底変換と可換である．
特に，$V \times_{Y'} \mathcal{X}'$ は射
$\mathcal{V} \to V \times_{Y'} \mathcal{X}'$ のスキーム論的像である．
一方，Morphisms of Stacks, Lemma
\ref{stacks-morphisms-lemma-universally-closed-permanence}
により，$|\mathcal{V}| \to |V \times_{Y'} \mathcal{X}'|$ の像は閉である
（上で見たとおり $\mathcal{V} \to V$ は普遍同相射，したがって普遍閉であることを用いる）．
さらに，この像は稠密である（直前の議論と Morphisms of Stacks, Lemma
\ref{stacks-morphisms-lemma-topology-scheme-theoretic-image} を組み合わせよ）．
したがって $|\mathcal{V}| = |V \times_{Y'} \mathcal{X}'|$ を得る．
% P12-ERR-0098: Japanese separates the source's run-on conclusion into a sentence.
よって $\mathcal{V} \to V \times_{Y'} \mathcal{X}'$ は同型であり，
主張の証明が完了する．
\end{proof}




\section{Noether 的付値判定法}
\label{stacks-more-morphisms-section-Noetherian-valuative-criterion}

\noindent
本節では Noether 的な設定において，離散付値環のみを用いる代数スタックの射の
（精密化された）付値判定法を論じる．さまざまな変種があり，
必要に応じて今後さらに追加する．

\medskip\noindent
$f : \mathcal{X} \to \mathcal{Y}$ を代数スタック（または代数空間，スキーム）の
射とする．{\it 精密化された}付値判定法とは，ある性質をもつ射
$\mathcal{U} \to \mathcal{X}$ が与えられ，次の形の実線図式における
点線矢印の存在または一意性のみを調べるものである．
$$
\xymatrix{
\Spec(K) \ar[d] \ar[r] & \mathcal{U} \ar[r] & \mathcal{X} \ar[d] \\
\Spec(A) \ar[rr] \ar@{..>}[rru] & & \mathcal{Y}
}
$$
以下では，これまでに得た結果を記述するためにこの用語を用いる．
% P12-ERR-0099: Japanese renders the frozen typo "sofar" as its intended adverb.

\medskip\noindent
代数スタックの射に対する非 Noether 的な付値判定法は次のとおりである．
\begin{enumerate}
\item Morphisms of Stacks, Section
\ref{stacks-morphisms-section-valuative-second}
（対角射の分離性について），
\item Morphisms of Stacks, Section
\ref{stacks-morphisms-section-valuative-diagonal}
（分離性について），
\item Morphisms of Stacks, Section
\ref{stacks-morphisms-section-valutive-criterion}
（普遍閉性について），
\item Morphisms of Stacks, Section
\ref{stacks-morphisms-section-valutive-criterion-properness}
（固有性について）．
\end{enumerate}
代数空間については次の付値判定法がある．
\begin{enumerate}
\item Morphisms of Spaces, Section
\ref{spaces-morphisms-section-valuative-criterion-universally-closed}
（普遍閉性について），
\item Morphisms of Spaces, Lemma
\ref{spaces-morphisms-lemma-refined-valuative-criterion-universally-closed}
（普遍閉性についての精密化），
\item Morphisms of Spaces, Section
\ref{spaces-morphisms-section-valuative-separatedness}
（分離性について），
\item Morphisms of Spaces, Section
\ref{spaces-morphisms-section-valuative-criterion-properness}
（固有性について），
\item Decent Spaces, Section
\ref{decent-spaces-section-valuative-criterion-universally-closed}
（decent 空間の普遍閉性について），
\item Decent Spaces, Lemma
\ref{decent-spaces-lemma-re-characterize-universally-closed}
（代数空間間の decent 射の普遍閉性について），
\item Cohomology of Spaces, Section
\ref{spaces-cohomology-section-Noetherian-valuative-criterion}
には次の Noether 的付値判定法が含まれる．
\begin{enumerate}
\item Cohomology of Spaces, Lemma
\ref{spaces-cohomology-lemma-check-separated-dvr}
（離散付値環を用いる分離性について），
\item Cohomology of Spaces, Lemma
\ref{spaces-cohomology-lemma-check-proper-dvr}
（離散付値環を用いる固有性について），
\item Cohomology of Spaces, Remark
\ref{spaces-cohomology-remark-variant}
（完備離散付値環の場合への帰着方法を論じる）．
\end{enumerate}
\item Limits of Spaces, Section
\ref{spaces-limits-section-Noetherian-valuative-criterion}
は Noether 的付値判定法を論じ，次を含む．
\begin{enumerate}
\item Limits of Spaces, Lemma
\ref{spaces-limits-lemma-Noetherian-dvr-valuative-separation}
（離散付値環と一般点を用いる分離性について），
\item Limits of Spaces, Lemma
\ref{spaces-limits-lemma-Noetherian-dvr-valuative-proper}
（離散付値環と一般点を用いる固有性について），
\item Limits of Spaces, Lemma
\ref{spaces-limits-lemma-check-universally-closed-Noetherian}
（離散付値環を用いる普遍閉性について）．
\end{enumerate}
\item Limits of Spaces, Section
\ref{spaces-limits-section-refined-valuative-criteria}
は精密化された Noether 的付値判定法を論じ，次を含む．
\begin{enumerate}
\item Limits of Spaces, Lemmas
\ref{spaces-limits-lemma-refined-valuative-criterion-proper} and
\ref{spaces-limits-lemma-refined-valuative-criterion-universally-closed}
（離散付値環を用いる固有性についての精密化），
\item Limits of Spaces, Lemma
\ref{spaces-limits-lemma-refined-valuative-criterion-separated}
（離散付値環を用いる分離性についての精密化）．
\end{enumerate}
\end{enumerate}
スキームについては次の付値判定法がある．
\begin{enumerate}
\item Schemes, Section
\ref{schemes-section-valuative-criterion-universal-closedness}
（普遍閉性について），
\item Schemes, Section \ref{schemes-section-valuative-separatedness}
（分離性について），
\item Morphisms, Section \ref{morphisms-section-valuative-criteria}
（固有性について），
\item Morphisms, Lemma
\ref{morphisms-lemma-refined-valuative-criterion-universally-closed}
（普遍閉性についての精密化），
\item Limits, Section \ref{limits-section-Noetherian-valuative-criterion}
は Noether 的付値判定法を論じ，次を含む．
\begin{enumerate}
\item Limits, Lemma \ref{limits-lemma-Noetherian-dvr-valuative-separation}
（離散付値環と一般点を用いる分離性について），
\item Limits, Lemma \ref{limits-lemma-Noetherian-dvr-valuative-proper}
（離散付値環と一般点を用いる固有性について），
\item Limits, Lemma \ref{limits-lemma-check-universally-closed-Noetherian}
（離散付値環を用いる普遍閉性について）．
\end{enumerate}
\item Limits, Section \ref{limits-section-refined-valuative-criteria}
は精密化された Noether 的付値判定法を論じ，次を含む．
\begin{enumerate}
\item Limits, Lemmas
\ref{limits-lemma-refined-valuative-criterion-proper} and
\ref{limits-lemma-refined-valuative-criterion-universally-closed}
（離散付値環を用いる固有性についての精密化），
\item Limits, Lemma \ref{limits-lemma-refined-valuative-criterion-separated}
（離散付値環を用いる分離性についての精密化）．
\end{enumerate}
\item Limits, Section \ref{limits-section-nagata-valuative}
は Noether 基底上の付値判定法を論じ，そこでは基底上本質的有限型の
離散付値環を得ることができる．
\end{enumerate}
以上で既存の結果の一覧を終える．

\medskip\noindent
本節の結果の多くは，次の補題を適用して証明できる
（また，おそらくそうすべきである）が，常にそうしたわけではない．

\begin{lemma}
\label{stacks-more-morphisms-lemma-reach-point-closure-Noetherian}
$f : \mathcal{X} \to \mathcal{Y}$ を代数スタックの射とする．
$f$ は有限型で，$\mathcal{Y}$ は局所 Noether 的であると仮定する．
% P12-ERR-0100: Japanese supplies the copula and preposition omitted in the frozen source.
$y \in |\mathcal{Y}|$ を $|f|$ の像の閉包に属する点とする．
このとき，代数スタックの可換図式
$$
\xymatrix{
\Spec(K) \ar[r] \ar[d] & \mathcal{X} \ar[d]^f \\
\Spec(A) \ar[r] & \mathcal{Y}
}
$$
であって，$A$ が離散付値環，$K$ がその分数体であり，
$\Spec(A)$ の閉点が $y$ に写るものが存在する．
\end{lemma}

\begin{proof}
アフィンスキーム $V$，点 $v \in V$，および $v$ を $y$ に写す平滑射
$V \to \mathcal{Y}$ を選ぶ．写像 $|V| \to |\mathcal{Y}|$ は開であり，
Properties of Stacks, Lemma \ref{stacks-properties-lemma-points-cartesian}
により，$|\mathcal{X} \times_\mathcal{Y} V| \to |V|$ の像は
$|f|$ の像の逆像である．したがって点 $v$ は
$|\mathcal{X} \times_\mathcal{Y} V| \to |V|$ の像の閉包に属する．
$\mathcal{X} \times_\mathcal{Y} V \to V$ と点 $v$ について補題を証明すれば，
$f$ と $y$ について補題が従う．このようにして次段落の状況へ帰着する．

\medskip\noindent
補題と同じ条件を満たす $f : \mathcal{X} \to Y$ と $y \in |Y|$ があり，
$Y$ が Noether アフィンスキームであると仮定する．$f$ は準コンパクトなので
$\mathcal{X}$ は準コンパクトである．したがってアフィンスキーム $W$ と
全射平滑射 $W \to \mathcal{X}$ を選べる．このとき $|f|$ の像は
$|W| \to |Y|$ の像と同じである．このようにしてスキームの場合へ帰着し，
これは Limits, Lemma \ref{limits-lemma-reach-point-closure-Noetherian}
である．
\end{proof}

\begin{lemma}
\label{stacks-more-morphisms-lemma-check-separated-dvr}
$f : \mathcal{X} \to \mathcal{Y}$ を代数スタックの射とする．次を仮定する．
\begin{enumerate}
\item $\mathcal{Y}$ は局所 Noether 的である．
\item $f$ は局所有限型かつ準分離的である．
\item 任意の可換図式
$$
\xymatrix{
\Spec(K) \ar[r]_x \ar[d]_j & \mathcal{X} \ar[d]^f \\
\Spec(A) \ar[r]^y \ar@{-->}[ru] & \mathcal{Y}
}
$$
で，$A$ が離散付値環，$K$ がその分数体であるもの，および任意の $2$-射
$\gamma : y \circ j \to f \circ x$ に対し，点線矢印の圏
（Morphisms of Stacks, Definition
\ref{stacks-morphisms-definition-fill-in-diagram}）は空であるか，
ちょうど一つの同型類をもつ setoid である．
\end{enumerate}
このとき $f$ は分離的である．
\end{lemma}

\begin{proof}
$f$ が分離的であることを証明するには，
$\Delta : \mathcal{X} \to \mathcal{X} \times_\mathcal{Y} \mathcal{X}$
が固有であることを示せばよい．$\Delta$ は代数空間によって表現可能で，
局所有限型であり（Morphisms of Stacks, Lemma
\ref{stacks-morphisms-lemma-properties-diagonal}），また準コンパクトかつ
準分離的であること（$f$ が準分離的であることの定義による）は既知である．
スキーム $U$ と全射平滑射
$U \to \mathcal{X} \times_\mathcal{Y} \mathcal{X}$ を選ぶ．
$$
V = \mathcal{X} \times_{\Delta, \mathcal{X} \times_\mathcal{Y} \mathcal{X}} U
$$
と置く．代数空間の射 $V \to U$ が固有であることを示せば十分である
（Properties of Stacks, Lemma
\ref{stacks-properties-lemma-check-property-covering}）．
$U$ は局所 Noether 的であること
（Morphisms of Stacks, Lemma
\ref{stacks-morphisms-lemma-locally-finite-type-locally-noetherian}
と $U \to \mathcal{Y}$ が局所有限型であることを用いる），また
$V \to U$ は有限型かつ準分離的であることに注意せよ
（これは $\Delta$ の基底変換であり，上に挙げた $\Delta$ の性質を用いる）．
Cohomology of Spaces, Lemma \ref{spaces-cohomology-lemma-check-proper-dvr}
を適用すると，次を示せば十分である．すなわち，可換図式
$$
\xymatrix{
\Spec(K) \ar[r]_v \ar[d]_j &
V \ar[d]^g \ar[r] &
\mathcal{X} \ar[d]^\Delta \\
\Spec(A) \ar[r]^u \ar@{-->}[ru] \ar@{..>}[rru] &
U \ar[r] &
\mathcal{X} \times_\mathcal{Y} \mathcal{X}
}
$$
が与えられ，$A$ が離散付値環，$K$ がその分数体であるとき，
図式を可換にする破線矢印が一意に存在することを示せばよい．
Morphisms of Stacks, Lemma
\ref{stacks-morphisms-lemma-cat-dotted-arrows-base-change}
により，破線矢印の圏と点線矢印の圏は同値である．仮定 (3) により，
同型を除いて一意な点線矢印が存在する．Morphisms of Stacks, Lemma
\ref{stacks-morphisms-lemma-helper-diagonal} を参照せよ．
したがって所望の一意な破線矢印が存在する．
\end{proof}

\begin{lemma}
\label{stacks-more-morphisms-lemma-refined-valuative-criterion-proper}
$f : \mathcal{X} \to \mathcal{Y}$ と $h : \mathcal{U} \to \mathcal{X}$
を代数スタックの射とする．$\mathcal{Y}$ は局所 Noether 的，
$f$ と $h$ は有限型，$f$ は分離的であり，
$|h| : |\mathcal{U}| \to |\mathcal{X}|$ の像は $|\mathcal{X}|$ において
稠密であると仮定する．任意の $2$-可換図式
$$
\xymatrix{
\Spec(K) \ar[r]_-u \ar[d]_j & \mathcal{U} \ar[r]_h & \mathcal{X} \ar[d]^f \\
\Spec(A) \ar[rr]^-y & & \mathcal{Y}
}
$$
で，$A$ が分数体 $K$ をもつ離散付値環であり，
$\gamma : y \circ j \to f \circ h \circ u$ であるものが与えられたとき，
体の拡大 $K'/K$ と，$A$ を支配する付値環 $A' \subset K'$ が存在し，
誘導される図式
% P12-ERR-0101: Japanese separates the hypothesis from the existential conclusion.
$$
\xymatrix{
\Spec(K') \ar[r]_-{x'} \ar[d]_{j'} & \mathcal{X} \ar[d]^f \\
\Spec(A') \ar[r]^-{y'} \ar@{..>}[ru] & \mathcal{Y}
}
$$
において，誘導される $2$-射
$\gamma' : y' \circ j' \to f \circ x'$ に関する点線矢印の圏が
空でないならば（Morphisms of Stacks, Definition
\ref{stacks-morphisms-definition-fill-in-diagram}），$f$ は固有である．
\end{lemma}

\begin{proof}
$f$ が普遍閉であることを証明すれば十分である．
$V$ をアフィンスキームとし，$V \to \mathcal{Y}$ を平滑射とする．
Properties of Stacks, Lemma
\ref{stacks-properties-lemma-points-cartesian}
により，$|\mathcal{U} \times_\mathcal{Y} V| \to
|\mathcal{X} \times_\mathcal{Y} V|$ の像 $I$ は $|h|$ の像の逆像である．
$|\mathcal{X} \times_\mathcal{Y} V| \to |\mathcal{X}|$ は開なので
（Morphisms of Stacks, Lemma \ref{stacks-morphisms-lemma-fppf-open}），
$I$ は $|\mathcal{X} \times_\mathcal{Y} V|$ において稠密である．
また，点線矢印の圏は基底変換に関してよく振る舞うので
（Morphisms of Stacks, Lemma
\ref{stacks-morphisms-lemma-cat-dotted-arrows-base-change}），
拡大後に点線矢印が存在するという仮定は射
$\mathcal{U} \times_\mathcal{Y} V \to
\mathcal{X} \times_\mathcal{Y} V \to V$
に継承される．したがって，射
$\mathcal{U} \times_\mathcal{Y} V \to \mathcal{X} \times_\mathcal{Y} V \to V$
は補題の仮定を満たす．
よって $\mathcal{Y}$ はアフィンスキームであると仮定してよい．

\medskip\noindent
$\mathcal{Y} = Y$ がアフィンスキームであると仮定する．
（以後は $2$-射 $\gamma$ と $\gamma'$ を考慮する必要はない．
Morphisms of Stacks, Lemma
\ref{stacks-morphisms-lemma-cat-dotted-arrows-independent} を参照せよ．）
このとき $\mathcal{U}$ は準コンパクトである．アフィンスキーム $U$ と
全射平滑射 $U \to \mathcal{U}$ を選ぶ．ここで $\mathcal{U}$ を $U$ で
置き換えてよく，実際そうする．したがって $\mathcal{U}$ は
アフィンスキームであると仮定してよい．

\medskip\noindent
$\mathcal{Y} = Y$ と $\mathcal{U} = U$ がアフィンスキームであると仮定する．
Chow の補題（Theorem \ref{stacks-more-morphisms-theorem-chow-finite-type}）により，
$X$ が代数空間であるような全射固有射 $X \to \mathcal{X}$ を選べる．
以下では，$X \to Y$ が分離射の合成として分離的であることを用いる．
代数空間 $W = X \times_\mathcal{X} U$ を考える．射影 $W \to X$ は有限型である．
$X$ を $W \to X$ のスキーム論的像で置き換えてよい．したがって
$|X|$ における $|W|$ の像は $|X|$ において稠密であると仮定してよい
（ここでは $|h|$ の像が $|\mathcal{X}|$ において稠密であることを用いるので，
この置換後も射 $X \to \mathcal{X}$ は全射である）．
任意の実線可換図式
$$
\xymatrix{
\Spec(K) \ar[r] \ar[d] & W \ar[r] & X \ar[d] \\
\Spec(A) \ar[rr] \ar@{..>}[rru] & & Y
}
$$
で，$A$ が分数体 $K$ をもつ離散付値環であるものに対し，
図式を可換にする点線矢印が存在すると主張する．
まず，$K$ を拡大で，$A$ をその拡大内で $A$ を支配する付値環で置き換えた後に
点線矢印が存在することを証明すれば十分である．Morphisms of Spaces, Lemma
\ref{spaces-morphisms-lemma-push-down-solution} を参照せよ．
補題の仮定により，拡大 $K'/K$，$A$ を支配する付値環
$A' \subset K'$，および合成 $\Spec(A') \to \Spec(A) \to Y$ を持ち上げ，
合成 $\Spec(K') \to \Spec(K) \to W \to X$ と両立する矢印
$\Spec(A') \to \mathcal{X}$ を得る．$X \to \mathcal{X}$ は固有なので，
固有性の付値判定法
（Morphisms of Stacks, Lemma \ref{stacks-morphisms-lemma-criterion-proper}）
を用いて，拡大 $K''/K'$，$A'$ を支配する付値環 $A'' \subset K''$，および
合成 $\Spec(A'') \to \Spec(A') \to \mathcal{X}$ を持ち上げ，
合成 $\Spec(K'') \to \Spec(K') \to \Spec(K) \to X$ と両立する射
$\Spec(A'') \to X$ を見いだせる．このとき $K''/K$，$A'' \subset K''$，
および射 $\Spec(A'') \to X$ は，上で述べた問題の解をなす．したがって主張が従う．
% P12-ERR-0102: Japanese gives the source's three-part subject a joint predicate.

\medskip\noindent
この主張と代数空間に対する補題の場合
（Limits of Spaces, Lemma
\ref{spaces-limits-lemma-refined-valuative-criterion-proper}）
により，射 $X \to Y$ は固有である．Morphisms of Stacks, Lemma
\ref{stacks-morphisms-lemma-image-proper-is-proper}
により，所望のとおり $\mathcal{X} \to Y$ が固有であることを得る．
\end{proof}

\begin{lemma}
\label{stacks-more-morphisms-lemma-refined-valuative-criterion-separated}
$f : \mathcal{X} \to \mathcal{Y}$ と $h : \mathcal{U} \to \mathcal{X}$
を代数スタックの射とする．$\mathcal{Y}$ は局所 Noether 的，
$f$ は局所有限型かつ準分離的，$h$ は有限型であり，
$|h| : |\mathcal{U}| \to |\mathcal{X}|$ の像は $|\mathcal{X}|$ において
稠密であると仮定する．任意の $2$-可換図式
$$
\xymatrix{
\Spec(K) \ar[r]_-u \ar[d]_j & \mathcal{U} \ar[r]_h & \mathcal{X} \ar[d]^f \\
\Spec(A) \ar[rr]^-y \ar@{..>}[rru] & & \mathcal{Y}
}
$$
で，$A$ が分数体 $K$ をもつ離散付値環であり，
$\gamma : y \circ j \to f \circ h \circ u$ であるものが与えられたとき，
点線矢印の圏が空であるか，ちょうど一つの同型類をもつ setoid ならば，
$f$ は分離的である．
\end{lemma}

\begin{proof}
$\Delta$ が固有射であることを証明しなければならない．
まず $\Delta$ が分離的であると仮定する．このとき，射
$\mathcal{U} \to \mathcal{X}$ と
$\Delta : \mathcal{X} \to \mathcal{X} \times_\mathcal{Y} \mathcal{X}$
に Lemma \ref{stacks-more-morphisms-lemma-refined-valuative-criterion-proper} を適用できる．
$f$ が準分離的なので $\Delta$ は準コンパクトであることに注意せよ．
もちろん $\Delta$ は局所有限型である（これは任意の対角射について成り立つ．
Morphisms of Stacks, Lemma
\ref{stacks-morphisms-lemma-properties-diagonal} を参照せよ）．
最後に，$2$-可換図式
$$
\xymatrix{
\Spec(K) \ar[r]_-u \ar[d]_j &
\mathcal{U} \ar[r]_h &
\mathcal{X} \ar[d]^\Delta \\
\Spec(A) \ar[rr]^-y \ar@{..>}[rru] & &
\mathcal{X} \times_\mathcal{Y} \mathcal{X}
}
$$
で，$A$ が分数体 $K$ をもつ離散付値環であり，
$\gamma : y \circ j \to \Delta \circ h \circ u$ であるものが
与えられたとする．Morphisms of Stacks, Lemma
\ref{stacks-morphisms-lemma-helper-diagonal} と補題の仮定により，
一意な点線矢印が存在する．これは Lemma
\ref{stacks-more-morphisms-lemma-refined-valuative-criterion-proper} の最後の仮定が
成り立つことを示し，結論が従う．

\medskip\noindent
一般の場合，$\Delta$ が分離的であることを証明すれば十分である．
そうすれば前の場合に戻るからである．実際，射
$$
\mathcal{U} \to \mathcal{X}
\quad\text{かつ}\quad
\Delta :
\mathcal{X} \to
\mathcal{X} \times_\mathcal{Y} \mathcal{X}
$$
について補題の仮定が成り立つと主張する．すなわち，$\Delta$ は代数空間によって
表現可能なので，前段落のような図式に対する点線矢印の圏は setoid である
（例えば Morphisms of Stacks, Lemma
\ref{stacks-morphisms-lemma-cat-dotted-arrows} を参照）．
前段落の議論により，これらの圏は空であるか，一つの同型類をもつ．
したがって $\Delta$ は分離的である．
\end{proof}

\begin{lemma}
\label{stacks-more-morphisms-lemma-valuative-criterion-universally-closed}
$f : \mathcal{X} \to \mathcal{Y}$ を代数スタックの射とする．
$\mathcal{Y}$ は局所 Noether 的で，$f$ は有限型であると仮定する．
任意の $2$-可換図式
$$
\xymatrix{
\Spec(K) \ar[r]_-x \ar[d]_j & \mathcal{X} \ar[d]^f \\
\Spec(A) \ar[r]^-y & \mathcal{Y}
}
$$
で，$A$ が分数体 $K$ をもつ離散付値環であり，
$\gamma : y \circ j \to f \circ x$ であるものが与えられたとき，
体の拡大 $K'/K$ と，$A$ を支配する付値環 $A' \subset K'$ が存在し，
誘導される図式
% P12-ERR-0103: Japanese separates the hypothesis from the existential conclusion.
$$
\xymatrix{
\Spec(K') \ar[r]_-{x'} \ar[d]_{j'} & \mathcal{X} \ar[d]^f \\
\Spec(A') \ar[r]^-{y'} \ar@{..>}[ru] & \mathcal{Y}
}
$$
において，誘導される $2$-射
$\gamma' : y' \circ j' \to f \circ x'$ に関する点線矢印の圏が
空でないならば（Morphisms of Stacks, Definition
\ref{stacks-morphisms-definition-fill-in-diagram}），$f$ は普遍閉である．
\end{lemma}

\begin{proof}
$V$ をアフィンスキームとし，$V \to \mathcal{Y}$ を平滑射とする．
点線矢印の圏は基底変換に関してよく振る舞う
（Morphisms of Stacks, Lemma
\ref{stacks-morphisms-lemma-cat-dotted-arrows-base-change}）．
したがって，拡大後に点線矢印が存在するという仮定は射
$\mathcal{X} \times_\mathcal{Y} V \to V$ に継承される．
よって射 $\mathcal{X} \times_\mathcal{Y} V \to V$ は
補題の仮定を満たす．したがって $\mathcal{Y}$ は
アフィンスキームであると仮定してよい．

\medskip\noindent
$\mathcal{Y} = Y$ が Noether アフィンスキームであると仮定する．
（以後は $2$-射 $\gamma$ と $\gamma'$ を考慮する必要はない．
Morphisms of Stacks, Lemma
\ref{stacks-morphisms-lemma-cat-dotted-arrows-independent} を参照せよ．）
$f$ が普遍閉であることを証明するには，Limits of Stacks, Lemma
\ref{stacks-limits-lemma-test-universally-closed} により，すべての $n$ に対し
$|\mathcal{X} \times \mathbf{A}^n| \to |Y \times \mathbf{A}^n|$
が閉写像であることを示せば十分である．補題の仮定は積射
$\mathcal{X} \times \mathbf{A}^n \to Y \times \mathbf{A}^n$
に継承されるので（詳細は省略する），$|\mathcal{X}| \to |Y|$ が
閉写像であることの証明へ帰着する．

\medskip\noindent
$Y$ が Noether アフィンスキームであると仮定する．
$T \subset |\mathcal{X}|$ を閉部分集合とする．$|Y|$ における $T$ の像が
閉であることを示さなければならない．$\mathcal{X}$ を $T$ 上に誘導される
被約閉部分空間構造で置き換えてよい．点線矢印の存在に関する性質が
この置換によって保存されることの確認は省略する．
こうして $|\mathcal{X}| \to |Y|$ の像が閉であることの証明へ帰着する．

\medskip\noindent
$y \in |Y|$ を $|\mathcal{X}| \to |Y|$ の像の閉包に属する点とする．
Lemma \ref{stacks-more-morphisms-lemma-reach-point-closure-Noetherian} により，可換図式
$$
\xymatrix{
\Spec(K) \ar[r] \ar[d] & \mathcal{X} \ar[d]^f \\
\Spec(A) \ar[r] & Y
}
$$
で，$A$ が離散付値環，$K$ がその分数体であり，$\Spec(A)$ の閉点を $y$ に
写すものを選べる．補題の仮定から直ちに，$y$ は
$|\mathcal{X}| \to |Y|$ の像に属することが従い，証明が完了する．
\end{proof}

\section{モジュライ空間}
\label{stacks-more-morphisms-section-moduli-spaces}

\noindent
本節では，代数スタックから代数空間への射
$f : \mathcal{X} \to Y$ を論じる．適切な仮定のもとで，
$Y$ を $\mathcal{X}$ の {\it モジュライ空間}という．
$\mathcal{X} = [U/R]$ が一つの表示ならば，$R$-不変射 $U \to Y$ が得られ，
適切な仮定のもとで $Y$ は亜群 $(U, R, s, t, c)$ の {\it 商}である．
異なる種類の商については Quotients of Groupoids, Section
\ref{groupoids-quotients-section-introduction} 以降で論じられている．

\begin{definition}
\label{stacks-more-morphisms-definition-categorical-quotient}
$\mathcal{X}$ を代数スタックとし，$f : \mathcal{X} \to Y$ を
代数空間 $Y$ への射とする．
\begin{enumerate}
\item 代数空間 $W$ への任意の射 $\mathcal{X} \to W$ が $f$ を経由して
一意に分解するとき，$f$ を {\it 圏論的モジュライ空間}という．
\item 代数空間の任意の平坦射 $Y' \to Y$ に対し，基底変換
$f' : Y' \times_Y \mathcal{X} \to Y'$ が圏論的モジュライ空間となるとき，
$f$ を {\it 一様圏論的モジュライ空間}という．
\end{enumerate}
$\mathcal{C}$ を代数空間の圏の充満部分圏とする．
\begin{enumerate}
\item[(3)] $Y \in \Ob(\mathcal{C})$ であり，$W \in \Ob(\mathcal{C})$ を
満たす任意の射 $\mathcal{X} \to W$ が $f$ を経由して一意に分解するとき，
$f$ を {\it $\mathcal{C}$ における圏論的モジュライ空間}という．
\item[(4)] $Y \in \Ob(\mathcal{C})$ であり，$\mathcal{C}$ における任意の平坦射
$Y' \to Y$ に対し，基底変換 $f' : Y' \times_Y \mathcal{X} \to Y'$ が
$\mathcal{C}$ における圏論的モジュライ空間となるとき，
これを {\it $\mathcal{C}$ における一様圏論的モジュライ空間}という．
% P12-ERR-0104: Japanese uses an anaphoric subject for the one omitted in the frozen source.
\end{enumerate}
\end{definition}

\noindent
Yoneda の補題により，圏論的モジュライ空間は存在すれば一意である．
これを商について導入した用語と対応させよう．

\begin{lemma}
\label{stacks-more-morphisms-lemma-quotient-compare}
$(U, R, s, t, c)$ を代数空間の亜群とし，
$s, t : R \to U$ は平坦かつ局所有限表示であるとする．
代数スタック $\mathcal{X} = [U/R]$ を考える．
代数空間 $Y$ が与えられたとき，射 $f : \mathcal{X} \to Y$ と
$R$-不変射 $\phi : U \to Y$ との間には $1$ 対 $1$ の対応がある．
\end{lemma}

\begin{proof}
Criteria for Representability, Theorem
\ref{criteria-theorem-flat-groupoid-gives-algebraic-stack}
により，$\mathcal{X}$ は代数スタックである．射
$f : \mathcal{X} \to Y$ が与えられたとき，$\phi : U \to Y$ を合成
$U \to \mathcal{X} \to Y$ とする．
$R = U \times_\mathcal{X} U$ なので（Groupoids in Spaces, Lemma
\ref{spaces-groupoids-lemma-quotient-stack-2-cartesian}），
$\phi$ が $R$-不変であることは直ちに従う．逆に，$\phi : U \to Y$ が
代数空間への $R$-不変射ならば，Groupoids in Spaces, Lemma
\ref{spaces-groupoids-lemma-quotient-stack-2-coequalizer}
により射 $f : \mathcal{X} \to Y$ を得る．
Groupoids in Spaces, Lemma
\ref{spaces-groupoids-lemma-quotient-stack-objects}
にある商スタックの明示的記述を用いて，$\phi$ から $f$ を構成することもできる．
\end{proof}

\begin{lemma}
\label{stacks-more-morphisms-lemma-categorical-quotient-compare}
Lemma \ref{stacks-more-morphisms-lemma-quotient-compare} と同じ仮定および記法を用いる．
% P12-ERR-0105: Japanese supplies the determiner/plural structure missing in the frozen source.
このとき，$f$ が（一様）圏論的モジュライ空間であることと，
$\phi$ が（一様）圏論的商であることとは同値である．
充満部分圏におけるモジュライ空間についても同様である．
\end{lemma}

\begin{proof}
Lemma \ref{stacks-more-morphisms-lemma-quotient-compare} で確立した $1$ 対 $1$ の対応から，
$f$ が圏論的モジュライ空間であることと $\phi$ が圏論的商であることとは
直ちに同値となる（Quotients of Groupoids, Definition
\ref{groupoids-quotients-definition-categorical}）．
$Y' \to Y$ が射ならば，$U \to \mathcal{X}$ の基底変換として
$U' = Y' \times_Y U \to Y' \times_Y \mathcal{X} = \mathcal{X}'$
は全射，平坦かつ局所有限表示である
（Criteria for Representability, Lemma
\ref{criteria-lemma-flat-quotient-flat-presentation}）．
また，ファイバー積の結合則により
$R' = Y' \times_Y R$ は $U' \times_{\mathcal{X}'} U'$ に等しい．
したがって $\mathcal{X}' = [U'/R']$ である．Algebraic Stacks, Remark
\ref{algebraic-remark-flat-fp-presentation} を参照せよ．
よって，この状況の $Y'$ への基底変換は，補題の主張と同じ型の状況である．
これより，$f$ が一様圏論的モジュライ空間であることと，
$\phi$ が一様圏論的商であることとの同値が直ちに従う．
\end{proof}

\begin{lemma}
\label{stacks-more-morphisms-lemma-check-uniform-categorical-quotient-on-affines}
$f : \mathcal{X} \to Y$ を代数スタックから代数空間への射とする．
任意のアフィンスキーム $Y'$ と平坦射 $Y' \to Y$ に対し，基底変換
$f' : Y' \times_Y \mathcal{X} \to Y'$ が圏論的モジュライ空間ならば，
$f$ は一様圏論的モジュライ空間である．
\end{lemma}

\begin{proof}
$Y_i$ がアフィンスキームであるような \'etale 被覆
$\{Y_i \to Y\}$ を選ぶ．各 $i$ および $j$ に対し，アフィン開被覆
$Y_i \times_Y Y_j = \bigcup Y_{ijk}$ を選ぶ．
% P12-ERR-0106: Japanese is unaffected by the frozen source's article error.
$\mathcal{X}_i = Y_i \times_Y \mathcal{X}$ および
$\mathcal{X}_{ijk} = Y_{ijk} \times_Y \mathcal{X}$ と置く．
$g : \mathcal{X} \to W$ を代数空間への射とする．図式
$$
\xymatrix{
\mathcal{X}_i \ar[r] \ar[d] & \mathcal{X} \ar[d] \ar[r]_g & W \\
Y_i \ar[r] \ar@{..>}[rru] & Y
}
$$
を考える．$\mathcal{X}_i \to Y_i$ が圏論的モジュライ空間であるという仮定から，
一意な点線矢印 $h_i : Y_i \to W$ が得られる．
$\mathcal{X}_{ijk} \to Y_{ijk}$ が圏論的モジュライ空間であるという仮定から，
$h_i$ と $h_j$ の $Y_{ijk}$ への制限は等しい．
% P12-ERR-0107: Japanese explicitly gives the plural restrictions the source intends.
したがって $h_i$ と $h_j$ は $Y_i \times_Y Y_j$ 上で一致する．
$Y = \coprod Y_i / \coprod Y_i \times_Y Y_j$ なので
（Spaces, Section \ref{spaces-section-presentations}），$g$ が経由して分解する
一意な射 $Y \to W$ が存在する．したがって $f$ は圏論的モジュライ空間である．
同じ議論は平坦基底変換後にも適用できるので，$f$ は一様圏論的モジュライ空間である．
\end{proof}

\section{Keel--Mori の定理}
\label{stacks-more-morphisms-section-Keel-Mori}

\noindent
本節では，代数スタックの枠組みにおける Keel--Mori の定理の議論を
始める．関連文献については，Guide to Literature, Subsection
\ref{guide-subsection-coarse-moduli-spaces} を参照されたい．

\begin{definition}
\label{stacks-more-morphisms-definition-well-nigh-affine}
代数スタック $\mathcal{X}$ をとる．$\mathcal{X}$ が
{\it ほとんどアフィン}であるとは，アフィンスキーム $U$ と
全射，平坦，有限かつ有限表示な射
$U \to \mathcal{X}$ が存在することをいう．
\end{definition}

\noindent
この性質をそれほど頻繁に使うつもりはないので，
いくぶん滑稽な名前を付けた．

\begin{lemma}
\label{stacks-more-morphisms-lemma-well-nigh-affine}
代数スタック $\mathcal{X}$ をとる．次の条件は同値である．
\begin{enumerate}
\item $\mathcal{X}$ はほとんどアフィンである．
\item 亜群スキーム $(U, R, s, t, c)$ であって，$U$ と
$R$ がアフィンであり，$s, t : R \to U$ が有限局所自由で，
$\mathcal{X} = [U/R]$ を満たすものが存在する．
\end{enumerate}
これらの条件が成り立つならば，$\mathcal{X}$ は準コンパクト，準 DM かつ分離である．
\end{lemma}

\begin{proof}
$\mathcal{X}$ がほとんどアフィンであると仮定する．アフィンスキーム $U$ と
全射，平坦，有限かつ有限表示な射
$U \to \mathcal{X}$ をとる．$R = U \times_\mathcal{X} U$ とおく．このとき，
代数空間における亜群 $(U, R, s, t, c)$ と同型
$[U/R] \to \mathcal{X}$ を得る．Algebraic Stacks, Lemma
\ref{algebraic-lemma-map-space-into-stack}
および Remark \ref{algebraic-remark-flat-fp-presentation} を参照されたい．
$s, t : R \to U$ は
% P12-ERR-0108: Preserve the frozen closing parenthesis inside the math span.
（$U \to \mathcal{X})$ の基底変換として）平坦，有限かつ有限表示な射なので，
$s, t$ は有限局所自由である
（Morphisms, Lemma \ref{morphisms-lemma-finite-flat}）．
したがって $R$ はアフィンであり（有限射はアフィン射だからである），
(2) が成り立つ．

\medskip\noindent
% P12-ERR-0109: Japanese supplies normal grammar for the frozen extra copula.
亜群スキーム $(U, R, s, t, c)$ が与えられ，$U$ と
$R$ がアフィンであり，$s, t : R \to U$ が有限局所自由であるとする．
$\mathcal{X} = [U/R]$ とおく．このとき $\mathcal{X}$ は，
Criteria for Representability, Theorem
\ref{criteria-theorem-flat-groupoid-gives-algebraic-stack} により代数スタックである
（厳密にいえば，ここではこの事実は必要ないが，この事実はいくら強調してもしすぎることはない）．
射 $U \to \mathcal{X}$ は，Criteria for Representability, Lemma
\ref{criteria-lemma-flat-quotient-flat-presentation} により，
全射，平坦かつ局所有限表示である．
したがって，$U \to \mathcal{X}$ が有限であるかどうかは，射影
$U \times_\mathcal{X} U \to U$ が有限であるかどうかを調べれば判定できる．
Properties of Stacks, Lemma
\ref{stacks-properties-lemma-check-property-covering} を参照されたい．
$U \times_\mathcal{X} U = R$ である
（Groupoids in Spaces, Lemma
\ref{spaces-groupoids-lemma-quotient-stack-2-cartesian}）ので，実際に有限である．したがって
$\mathcal{X}$ はほとんどアフィンである．

\medskip\noindent
最後の主張を証明する．$\mathcal{X}$ は準コンパクトである
（Properties of Stacks, Lemma
\ref{stacks-properties-lemma-quasi-compact-stack}）．また，
$\mathcal{X} = [U/R]$ は，Morphisms of Stacks, Lemma
\ref{stacks-morphisms-lemma-properties-diagonal-from-presentation} により
準 DM かつ分離である．
\end{proof}

\begin{lemma}
\label{stacks-more-morphisms-lemma-affine-over-well-nigh-affine}
代数スタック $\mathcal{X}$ はほとんどアフィンであるとする．
\begin{enumerate}
\item $\mathcal{X}$ が代数空間ならば，これはアフィンである．
\item $\mathcal{X}' \to \mathcal{X}$ が代数スタックのアフィン射ならば，
$\mathcal{X}'$ はほとんどアフィンである．
\end{enumerate}
\end{lemma}

\begin{proof}
% P12-ERR-0110: Japanese removes the frozen duplicated preposition.
(1) は直ちに Limits of Spaces, Lemma
\ref{spaces-limits-lemma-affine} から従う．
しかし，これは必要以上に強い結果を用いている．(1) は Lemma
\ref{stacks-more-morphisms-lemma-well-nigh-affine} と Groupoids, Proposition
\ref{groupoids-proposition-finite-flat-equivalence} を組み合わせても従う．

\medskip\noindent
(2) を証明するため，アフィンスキーム $U$ と全射，平坦，有限かつ
有限表示な射 $U \to \mathcal{X}$ をとる．
このとき $U' = \mathcal{X}' \times_\mathcal{X} U$ から $U$ への射はアフィンである
（Morphisms of Stacks, Lemma
\ref{stacks-morphisms-lemma-base-change-affine}）．
したがって $U'$ はアフィンスキームである．もちろん
$U' \to \mathcal{X}'$ は $U \to \mathcal{X}$ の基底変換として，
全射，平坦，有限かつ有限表示である．
\end{proof}

\begin{lemma}
\label{stacks-more-morphisms-lemma-well-nigh-affine-moduli-space}
代数スタック $\mathcal{X}$ はほとんどアフィンであるとする．
アフィンスキームの圏における一様圏論的モジュライ空間
$$
f : \mathcal{X} \longrightarrow M
$$
が存在する．さらに，$f$ は分離かつ準コンパクトであり，普遍同相射である．
\end{lemma}

\begin{proof}
$\mathcal{X} = [U/R]$ と $(U, R, s, t, c)$ を Lemma
\ref{stacks-more-morphisms-lemma-well-nigh-affine} のように書く．$C$ を
$R$ に関して不変な $U$ 上の関数の環とする．Groupoids, Section
\ref{groupoids-section-finite-flat} を参照されたい．
$M = \Spec(C)$ とおく．$R$-不変な射
$U \to M$ は，射 $f : \mathcal{X} \to M$ に対応する
（Lemma \ref{stacks-more-morphisms-lemma-quotient-compare}）．
Schemes, Lemma \ref{schemes-lemma-morphism-into-affine} による
% P12-ERR-0111: Preserve the frozen source's previously undefined phi.
アフィンスキームへの射の特徴づけから，$\phi : U \to M$ は
アフィンスキームの圏における圏論的商であることが直ちに保証される．
したがって $f$ は，アフィンスキームの圏における
圏論的モジュライ空間である
（Lemma \ref{stacks-more-morphisms-lemma-categorical-quotient-compare}）．

\medskip\noindent
$\mathcal{X}$ は分離である（Lemma \ref{stacks-more-morphisms-lemma-well-nigh-affine}）ので，
$f$ は分離である（Morphisms of Stacks, Lemma
\ref{stacks-morphisms-lemma-compose-after-separated}）．

\medskip\noindent
$U \to \mathcal{X}$ は全射であり，$U \to M$ は準コンパクトなので，
Morphisms of Stacks, Lemma
\ref{stacks-morphisms-lemma-surjection-from-quasi-compact} により，$f$ は準コンパクトである．

\medskip\noindent
Groupoids, Lemma \ref{groupoids-lemma-integral-over-invariants} により，合成
$$
U \to \mathcal{X} \to M
$$
はアフィンスキームの整射である．とくに，これは普遍閉である
（Morphisms, Lemma \ref{morphisms-lemma-integral-universally-closed}）．
$U \to \mathcal{X}$ は全射なので，$\mathcal{X} \to M$ は普遍閉である
（Morphisms of Stacks, Lemma
\ref{stacks-morphisms-lemma-image-proper-is-proper}）．
$\mathcal{X} \to M$ が普遍同相射であると結論するには，
これが普遍全単射，すなわち全射かつ普遍単射であることを示せば十分である．

\medskip\noindent
$|\mathcal{X}| = |U|/|R|$ である
（Morphisms of Stacks, Lemma \ref{stacks-morphisms-lemma-points-presentation}）．
したがって $|f|$ は全射であり，さらに全単射である
（Groupoids, Lemma \ref{groupoids-lemma-points}）．

\medskip\noindent
環準同型 $C \to C'$ をとる．$(U', R', s', t', c')$ を
$(U, R, s, t, c)$ の $M' = \Spec(C') \to M$ による基底変換とする．
$\mathcal{X}' = [U'/R']$ とおく．このとき
$M' \times_M \mathcal{X} = \mathcal{X}'$ である
（Quotients of Groupoids, Lemma
\ref{groupoids-quotients-lemma-base-change-quotient-stack}）．
$C^1$ を $R'$ に関して不変な $U'$ 上の関数の環とする．
$M^1 = \Spec(C^1)$ とおき，次の図式を考える．
$$
\xymatrix{
\mathcal{X}' \ar[d]^{f'} \ar[r] & \mathcal{X} \ar[dd]^f \\
M^1 \ar[d] \\
M' \ar[r] & M
}
$$
Groupoids, Lemma \ref{groupoids-lemma-invariants-base-change} および
Algebra, Lemma \ref{algebra-lemma-universally-bijective} により，
射 $M^1 \to M'$ は同相写像である．
一方，前の段落を $(U', R', s', t', c')$ に適用すると，
$|f'|$ は全単射であることがわかる．
したがって，$f$ はアフィンスキームによる任意の基底変換の後に
点の集合上の全単射を誘導する．ゆえに Morphisms of Stacks, Lemma
\ref{stacks-morphisms-lemma-universally-injective-local} により，
$f$ は普遍単射である．

\medskip\noindent
% P12-ERR-0112: Preserve the frozen omission of `categorical` in this sentence.
最後に，$f$ がアフィンスキームの圏における一様モジュライ空間であることを
示す必要がある．これは上の議論と，環準同型
$C \to C'$ が平坦ならば $C' \to C^1$ が同型であるという事実
（Groupoids, Lemma \ref{groupoids-lemma-invariants-base-change}）から従う．
\end{proof}

\begin{lemma}
\label{stacks-more-morphisms-lemma-well-nigh-affine-moduli-space-etale}
$h : \mathcal{X}' \to \mathcal{X}$ を代数スタックの射とする．
$\mathcal{X}'$ と $\mathcal{X}$ はほとんどアフィンであり，$h$ は \'etale で，
$h$ は自己同型群上に同型を誘導すると仮定する
（Morphisms of Stacks, Remark
\ref{stacks-morphisms-remark-identify-automorphism-groups}）．
このとき，カルテジアン図式
$$
\xymatrix{
\mathcal{X}' \ar[d] \ar[r] & \mathcal{X} \ar[d] \\
M' \ar[r] & M
}
$$
が存在する．ここで $M' \to M$ は \'etale であり，縦の矢印は
Lemma \ref{stacks-more-morphisms-lemma-well-nigh-affine-moduli-space} で構成したモジュライ空間への射である．
\end{lemma}

\begin{proof}
$h$ は代数空間によって表現可能であることに注意する．これは
Morphisms of Stacks, Lemmas
\ref{stacks-morphisms-lemma-aut-iso-unramified} および
\ref{stacks-morphisms-lemma-stabilizer-preserving} による．
アフィンスキーム $U$ と，全射，平坦，有限かつ有限表示な射
$U \to \mathcal{X}$ を選ぶ．このとき
$U' = \mathcal{X}' \times_\mathcal{X} U$ は代数空間であり，
$U' \to \mathcal{X}'$ は有限（とくにアフィン）である．
Lemma \ref{stacks-more-morphisms-lemma-affine-over-well-nigh-affine} により，$U'$ はアフィンである．
$R = U \times_\mathcal{X} U$ および
$R' = U' \times_{\mathcal{X}'} U'$ とおくと，亜群
$(U, R, s, t, c)$ と $(U', R', s', t', c')$ が得られ，
$\mathcal{X} = [U/R]$ および $\mathcal{X}' = [U'/R']$ となる．
Lemma \ref{stacks-more-morphisms-lemma-well-nigh-affine} の証明を参照されたい．
次の図式がカルテジアンであることが分かる．
$$
\xymatrix{
R' \ar[d]_{s'} \ar[r]_f & R \ar[d]^s \\
U' \ar[r]^f & U
}
\quad
\quad
\xymatrix{
R' \ar[d]_{t'} \ar[r]_f & R \ar[d]^t \\
U' \ar[r]^f & U
}
\quad
\quad
\xymatrix{
G' \ar[d] \ar[r]_f & G \ar[d] \\
U' \ar[r]^f & U
}
$$
ここで $G$ と $G'$ は安定化群スキームである．
最初の二つについてはファイバー積の推移性から従い，最後のものについては，同型
$\mathcal{I}_{\mathcal{X}'} \to
\mathcal{X}' \times_\mathcal{X} \mathcal{I}_\mathcal{X}$
の引き戻しであることから従う（この同型は，既に用いた
Morphisms of Stacks, Lemma
\ref{stacks-morphisms-lemma-aut-iso-unramified} による）．
$M$，resp.\ $M'$ は Lemma
\ref{stacks-more-morphisms-lemma-well-nigh-affine-moduli-space} において，それぞれ
$R$-不変な $U$ 上の関数の環，resp.\ $R'$-不変な $U'$ 上の関数の環の
スペクトルとして構成されたことを思い出そう．したがって
$M' \to M$ は \'etale であり，$U' = M' \times_M U$ である．これは
Groupoids, Lemma \ref{groupoids-lemma-etale} による．これより
% P12-ERR-0114: Preserve the frozen repeated-U base-change identity.
$R' = M' \times_M U$ となる．言い換えれば，亜群
$(U', R', s', t', c')$ は $(U, R, s, t, c)$ の $M' \to M$ による
基底変換である．これは，補題の図式がカルテジアンであることを意味する．
Quotients of Groupoids, Lemma
\ref{groupoids-quotients-lemma-base-change-quotient-stack} による．
\end{proof}

\begin{lemma}
\label{stacks-more-morphisms-lemma-moduli-space-finite-affine}
代数スタック $\mathcal{X}$ はほとんどアフィンであるとする．
$$
f : \mathcal{X} \longrightarrow M
$$
という Lemma \ref{stacks-more-morphisms-lemma-well-nigh-affine-moduli-space} の射は
一様圏論的モジュライ空間である．
\end{lemma}

\begin{proof}
% P12-ERR-0115: Preserve the frozen attribution of the property to M rather than f.
$M$ はアフィンスキームの圏における一様圏論的モジュライ空間であることを
既に知っている．Lemma
\ref{stacks-more-morphisms-lemma-check-uniform-categorical-quotient-on-affines} により，基底変換
$f' : M' \times_M \mathcal{X} \to M'$ が，アフィンスキームの任意の平坦射
$M' \to M$ に対して圏論的モジュライ空間であることを
示せば十分である．$\mathcal{X}' = M' \times_M \mathcal{X}$ は Lemma
\ref{stacks-more-morphisms-lemma-affine-over-well-nigh-affine} によりほとんどアフィンであることに
注意する．したがって $\mathcal{X}$ を $\mathcal{X}'$ で，$M$ を $M'$ で
置き換えることにより，$f$ が圏論的モジュライ空間であることの証明に帰着される．

\medskip\noindent
$g : \mathcal{X} \to Y$ を射とし，$Y$ は代数空間であるとする．
$g = h \circ f$ となる一意な射 $h : M \to Y$ が存在することを
示さなければならない．

\medskip\noindent
一意性．二つの射 $h_i : M \to Y$ があり，
$g = h_1 \circ f = h_2 \circ f$ を満たすとする．
$M' \subset M$ を $h_1$ と $h_2$ の等化子とする．
このとき $M' \to M$ はモノ射であり，$f : \mathcal{X} \to M$ は $M'$ を経由して
分解する．したがって $M' \to M$ は普遍同相射である．
$M'$ はアフィンであると結論される
（Morphisms, Lemma \ref{morphisms-lemma-universal-homeomorphism}）．
しかしこのとき，$f : \mathcal{X} \to M$ はアフィンスキームの圏における
圏論的モジュライ空間なので，$M' = M$ であることが分かる．

\medskip\noindent
存在性．以下で，すべての $p \in M$ に対し，カルテジアンな正方形
$$
\xymatrix{
\mathcal{X}' \ar[r] \ar[d] & \mathcal{X} \ar[d] \\
M' \ar[r] & M
}
$$
が存在することを示す．ここで $M' \to M$ はアフィンスキーム間の \'etale 射で，
その像は $p$ を含み，合成 $\mathcal{X}' \to \mathcal{X} \to Y$ は $M'$ を経由して
分解する．これは，写像 $h : M \to Y$ を $M$ 上 \'etale 局所的に構成できることを
意味する．$Y$ は \'etale 位相に関する層であり，上で一意性を示したので，
これで十分である（細部は省略する）．

\medskip\noindent
$y \in |Y|$ を $p$ の像とする．
$(V, v) \to (Y, y)$ を，$V$ がアフィンであるような \'etale 射とする．
$\mathcal{X}' = V \times_Y \mathcal{X}$ を考える．
$\mathcal{X}' \to \mathcal{X}$ は分離的かつ \'etale であることに注意する．
これは $V \to Y$ の基底変換だからである．さらに，
$\mathcal{X}' \to \mathcal{X}$ は自己同型群上に同型を誘導する
（Morphisms of Stacks, Remark
\ref{stacks-morphisms-remark-identify-automorphism-groups}）．
実際，これは $V \to Y$ について成り立つ．Morphisms of Stacks, Lemma
\ref{stacks-morphisms-lemma-base-change-stabilizer-preserving} を参照されたい．
表示 $\mathcal{X} = [U/R]$ を Lemma
\ref{stacks-more-morphisms-lemma-well-nigh-affine} と同様に選ぶ．
$U' = \mathcal{X}' \times_\mathcal{X} U = V \times_Y U$ とおき，
$u' \in U'$ を，$p$ と $v$ に写るように選ぶ（これは
Properties of Spaces, Lemma
\ref{spaces-properties-lemma-points-cartesian} により可能である）．
$U' \to U$ は分離的かつ \'etale なので，$U'$ の任意の有限個の点は
あるアフィン開集合に含まれる．More on Morphisms, Lemma
\ref{more-morphisms-lemma-separated-locally-quasi-finite-over-affine} を参照されたい．
一方，射 $U' \to \mathcal{X}'$ は全射，有限，平坦かつ局所有限表示である．
$R' = U' \times_{\mathcal{X}'} U'$ とおくと，
$s', t' : R' \to U'$ は有限局所自由であることが分かる．
Groupoids, Lemma \ref{groupoids-lemma-find-invariant-affine} により，
$R'$-不変なアフィン開部分スキーム $U'' \subset U'$ であって
$u'$ を含むものが存在する．
対応する開部分スタックを $\mathcal{X}'' \subset \mathcal{X}'$ とする．
このとき $\mathcal{X}''$ はほとんどアフィンである．Lemma
\ref{stacks-more-morphisms-lemma-well-nigh-affine-moduli-space-etale} により，カルテジアンな正方形
$$
\xymatrix{
\mathcal{X}'' \ar[r] \ar[d] & \mathcal{X} \ar[d] \\
M'' \ar[r] & M
}
$$
が得られる．ここで $M'' \to M$ は \'etale である．
$\mathcal{X}'' \to M''$ はアフィンスキームの圏における
圏論的モジュライ空間なので，射 $M'' \to V$ であって，合成
$\mathcal{X}'' \to \mathcal{X}' \to V$ と合成
$\mathcal{X}'' \to M'' \to V$ とが等しくなるものを得る．
これで主張が証明され，証明が完了する．
\end{proof}

\begin{lemma}
\label{stacks-more-morphisms-lemma-etale-separated-over-well-nigh-affine}
$h : \mathcal{X}' \to \mathcal{X}$ を代数スタックの射とする．
$\mathcal{X}$ はほとんどアフィンであり，$h$ は \'etale で，$h$ は分離的で，
$h$ は自己同型群上に同型を誘導すると仮定する
（Morphisms of Stacks, Remark
\ref{stacks-morphisms-remark-identify-automorphism-groups}）．
このとき，カルテジアン図式
$$
\xymatrix{
\mathcal{X}' \ar[d] \ar[r] & \mathcal{X} \ar[d] \\
M' \ar[r] & M
}
$$
が存在する．ここで $M' \to M$ はスキーム間の分離的 \'etale 射であり，
$\mathcal{X} \to M$ は Lemma
\ref{stacks-more-morphisms-lemma-well-nigh-affine-moduli-space} で構成したモジュライ空間である．
\end{lemma}

\begin{proof}
アフィンスキーム $U$ と，全射，平坦，有限かつ局所有限表示な射
$U \to \mathcal{X}$ を選ぶ．$h$ は代数空間によって表現可能なので
（Morphisms of Stacks, Lemmas
\ref{stacks-morphisms-lemma-aut-iso-unramified} および
\ref{stacks-morphisms-lemma-stabilizer-preserving}），
$U' = \mathcal{X}' \times_\mathcal{X} U$ は代数空間であることが分かる．
$U' \to U$ は分離的かつ \'etale なので，$U'$ はスキームであり，
$U'$ の任意の有限個の点はあるアフィン開集合に含まれる．
Morphisms of Spaces, Lemma
\ref{spaces-morphisms-lemma-locally-quasi-finite-separated-representable}
および More on Morphisms, Lemma
\ref{more-morphisms-lemma-separated-locally-quasi-finite-over-affine} を参照されたい．
$R' = U' \times_{\mathcal{X}'} U'$ とおくと，
$s', t' : R' \to U'$ は有限局所自由であることが分かる．
Groupoids, Lemma \ref{groupoids-lemma-find-invariant-affine} により，
開被覆 $U' = \bigcup U'_i$ が存在し，その被覆は
$R'$-不変なアフィン開部分スキーム $U'_i \subset U'$ からなる．
対応する開部分スタックを
$\mathcal{X}'_i \subset \mathcal{X}'$ とする．これらはほとんどアフィンである．
実際，$U'_i \to \mathcal{X}'_i$ は全射，平坦，有限かつ有限表示である．
Lemma \ref{stacks-more-morphisms-lemma-well-nigh-affine-moduli-space-etale} により，
カルテジアン図式
$$
\xymatrix{
\mathcal{X}'_i \ar[r] \ar[d] & \mathcal{X} \ar[d] \\
M'_i \ar[r] & M
}
$$
が得られる．ここで $M'_i \to M$ はアフィンスキーム間の \'etale 射であり，
縦の矢印は Lemma \ref{stacks-more-morphisms-lemma-well-nigh-affine-moduli-space} のものと同じである．
$\mathcal{X}'_{ij} = \mathcal{X}'_i \cap \mathcal{X}'_j$ は
% P12-ERR-0117: Preserve the frozen `open subspace` classification.
$\mathcal{X}'_i$ および $\mathcal{X}'_j$ の開部分空間であることに注意する．
したがって，対応する開部分スキーム
$V_{ij} \subset M'_i$ と $V_{ji} \subset M'_j$ を得る．
Lemma \ref{stacks-more-morphisms-lemma-moduli-space-finite-affine} の結果により，
$\mathcal{X}'_{ij} \to V_{ij}$ と
$\mathcal{X}'_{ji} \to V_{ji}$ はともに圏論的モジュライ空間である！
したがって，一意な同型 $\varphi_{ij} : V_{ij} \to V_{ji}$ であって，図式
$$
\xymatrix{
\mathcal{X}'_i \ar[d] & &
\mathcal{X}'_i \cap \mathcal{X}'_j \ar[rr] \ar[ll] \ar[ld] \ar[rd] & &
\mathcal{X}'_j \ar[d] \\
M'_i &
V_{ij} \ar[l] \ar[rr]^{\varphi_{ij}} & &
V_{ji} \ar[r] &
M'_j
}
$$
を可換にするものを得る．
これらの同型は Schemes, Section
\ref{schemes-section-glueing-schemes} のコサイクル条件を満たす．
これは計算（および前の補題のもう一度の適用）によるが，省略する．
したがってアフィンスキームを貼り合わせてスキーム $M'$ を得る．
Schemes, Lemma \ref{schemes-lemma-glue} を参照されたい．
$M'_i$ を $M'$ におけるその像と同一視しよう．
射 $\mathcal{X}' \to M'$ であって，カルテジアン図式
$$
\xymatrix{
\mathcal{X}'_i \ar[r] \ar[d] & \mathcal{X}' \ar[d] \\
M'_i \ar[r] & M'
}
$$
に入るものが存在すると主張する．
これは，貼り合わせたスキーム $M'$ への射についての Schemes, Lemma
\ref{schemes-lemma-glue} の記述と，射 $\mathcal{X}' \to M'$ を与えることが，
射 $T \to M'$ を任意の射 $T \to \mathcal{X}'$ に対して与えることと同じである
という事実から明らかである．同様に，射 $M' \to M$ であって，与えられた射
$M'_i \to M$ に $M'_i$ 上で制限されるものが存在する．
射 $M' \to M$ は \'etale である（\'etale 被覆の各要素上で \'etale だからである）．
また，ファイバー積の性質は（アフィン）開被覆
$M' = \bigcup M'_i$ の各要素上で確認できるので成り立つ．
最後に，$M' \to M$ は分離的である．実際，合成
$U' \to \mathcal{X}' \to M'$ は全射かつ普遍閉なので，
Morphisms, Lemma
\ref{morphisms-lemma-image-universally-closed-separated} を適用できる．
\end{proof}

\begin{lemma}
\label{stacks-more-morphisms-lemma-etale-local-finite-inertia}
$\mathcal{X}$ を代数スタックとする．
$\mathcal{I}_\mathcal{X} \to \mathcal{X}$ は有限であると仮定する．
このとき集合 $I$ と，各 $i \in I$ に対して代数スタックの射
$$
g_i : \mathcal{X}_i \longrightarrow \mathcal{X}
$$
であって，次の性質をもつものが存在する．
\begin{enumerate}
\item $|\mathcal{X}| = \bigcup |g_i|(|\mathcal{X}_i|)$ である，
\item $\mathcal{X}_i$ はほとんどアフィンである，
\item $\mathcal{I}_{\mathcal{X}_i} \to
\mathcal{X}_i \times_\mathcal{X} \mathcal{I}_\mathcal{X}$
は同型であり，
\item $g_i : \mathcal{X}_i \to \mathcal{X}$ は代数空間によって表現可能で，
分離かつ \'etale である．
\end{enumerate}
\end{lemma}

\begin{proof}
任意の $x \in |\mathcal{X}|$ に対し，
$g : \mathcal{U} \to \mathcal{X}$，$\mathcal{U} = [U/R]$，および $u$ を
Morphisms of Stacks,
Lemma \ref{stacks-morphisms-lemma-etale-local-quasi-DM-at-x}
におけるように選ぶことができる．
すると Morphisms of Stacks, Lemma
\ref{stacks-morphisms-lemma-stabilizer-preserving-unramified}
により，開部分スタック
$\mathcal{U}' \subset \mathcal{U}$ であって，$u$ を含み，
$\mathcal{I}_{\mathcal{U}'} \to
\mathcal{U}' \times_\mathcal{X} \mathcal{I}_\mathcal{X}$
が同型となるものが存在することが分かる．
$U' \subset U$ を，$R$-不変な開部分であって，
開部分スタック $\mathcal{U}'$ に対応するものとする．
$u' \in U'$ を $U'$ の点で $u$ に写るものとする．
$t(s^{-1}(\{u'\}))$ は有限であることに注意する．これは
$s : R \to U$ が有限だからである．
Properties, Lemma \ref{properties-lemma-ample-finite-set-in-affine}
および Groupoids, Lemma \ref{groupoids-lemma-find-invariant-affine}
により，$R$-不変なアフィン開部分
$U'' \subset U'$ であって，$u'$ を含むものを見つけることができる．
$R''$ を $R$ の $U''$ への制限とする．
すると $\mathcal{U}'' = [U''/R'']$ は
$\mathcal{U}'$ の開部分スタックであって $u$ を含み，ほとんどアフィンであり，
$\mathcal{I}_{\mathcal{U}''} \to
\mathcal{U}'' \times_\mathcal{X} \mathcal{I}_\mathcal{X}$
は同型であり，また $\mathcal{U}'' \to \mathcal{X}$ は
代数空間によって表現可能かつ \'etale である．
最後に，$\mathcal{U}'' \to \mathcal{X}$ は分離である．実際，
$\mathcal{U}''$ は分離であり（Lemma \ref{stacks-more-morphisms-lemma-well-nigh-affine}），
$\mathcal{X}$ の対角射は分離であり
（Morphisms of Stacks, Lemma \ref{stacks-morphisms-lemma-diagonal-diagonal}），
分離性は Morphisms of Stacks, Lemma
\ref{stacks-morphisms-lemma-compose-after-separated} から従う．
点 $x \in |\mathcal{X}|$ は任意であったから，証明は完了する．
\end{proof}

\begin{theorem}[Keel--Mori]
\label{stacks-more-morphisms-theorem-keel-mori}
$\mathcal{X}$ を代数スタックとする．
$\mathcal{I}_\mathcal{X} \to \mathcal{X}$ は有限であると仮定する．
このとき一様圏論的モジュライ空間
$$
f : \mathcal{X} \longrightarrow M
$$
が存在し，$f$ は分離，準コンパクト，かつ普遍同相射である．
\end{theorem}

\begin{proof}
集合 $I$\footnote{集合論的な問題をなお追っている読者は，
$I$ が大きすぎないことを確認すべきである．}
と，各 $i \in I$ に対して代数スタックの射
$g_i : \mathcal{X}_i \to \mathcal{X}$ を Lemma
\ref{stacks-more-morphisms-lemma-etale-local-finite-inertia} におけるように選ぶ．以下では，
この補題に列挙したすべての性質を断りなく用いる．
$$
f_i : \mathcal{X}_i \to M_i
$$
を Lemma \ref{stacks-more-morphisms-lemma-well-nigh-affine-moduli-space} におけるものとする．
スタック
$$
\mathcal{X}_{ij} = \mathcal{X}_i \times_{g_i, \mathcal{X}, g_j} \mathcal{X}_j
$$
を $i, j \in I$ に対して考える．
射影 $\mathcal{X}_{ij} \to \mathcal{X}_i$ および
$\mathcal{X}_{ij} \to \mathcal{X}_j$ は
Morphisms of Stacks, Lemma
\ref{stacks-morphisms-lemma-base-change-separated}
により分離であり，また \'etale である．これは
Morphisms of Stacks, Lemma
\ref{stacks-morphisms-lemma-base-change-etale}
による．さらに自己同型群上に同型を誘導する
（Morphisms of Stacks, Remark
\ref{stacks-morphisms-remark-identify-automorphism-groups}
における意味で）．これは Morphisms of Stacks, Lemma
\ref{stacks-morphisms-lemma-base-change-stabilizer-preserving}
による．したがって Lemma
\ref{stacks-more-morphisms-lemma-etale-separated-over-well-nigh-affine}
を適用して，カルテジアンな正方形からなる可換図式
$$
\xymatrix{
\mathcal{X}_i \ar[d]_{f_i} &
\mathcal{X}_{ij} \ar[d]_{f_{ij}} \ar[l] \ar[r] &
\mathcal{X}_j \ar[d]_{f_j} \\
M_i &
M_{ij} \ar[l] \ar[r] &
M_j
}
$$
であって，$M_{ij} \to M_i$ と $M_{ij} \to M_j$ が
スキームの分離な \'etale 射となるものを見つけることができる．
% P12-ERR-0123: Preserve source terminology `uniform categorical quotient` here.
ここではさらに，Lemma \ref{stacks-more-morphisms-lemma-moduli-space-finite-affine}
により $f_i$ が一様圏論的商であることも用いた．
次を主張する：
$$
\coprod M_{ij} \longrightarrow \coprod M_i \times \coprod M_i
$$
は \'etale 同値関係である．

\medskip\noindent
主張の証明．$R = \coprod M_{ij}$ および $U = \coprod M_i$ とおく．
$t : R \to U$ と $s : R \to U$ が \'etale であることはすでに見た．
射 $c : R \times_{s, U, t} R \to R$ であって，
$\text{pr}_{13} : U \times U \times U \to U \times U$
と両立するものを構成しよう．
すなわち，$i, j, k \in I$ に対して
$$
\mathcal{X}_{ijk} =
\mathcal{X}_i \times_{g_i, \mathcal{X}, g_j} \mathcal{X}_j
\times_{g_j, \mathcal{X}, g_k} \mathcal{X}_k =
\mathcal{X}_{ij} \times_{\mathcal{X}_j} \mathcal{X}_{jk}
$$
を考える．直前の段落とまったく同様に論じると，
$M_{ijk} = M_{ij} \times_{M_j} M_{jk}$ は
$\mathcal{X}_{ijk}$ の圏論的モジュライ空間であることが分かる．
特に，標準的な射
$M_{ijk} = M_{ij} \times_{M_j} M_{jk} \to M_{ik}$ が存在する．
これは射影 $\mathcal{X}_{ijk} \to \mathcal{X}_{ik}$ から生じる．
これらの射をまとめると，射 $c$ が得られる．
同様に，射 $e : U \to R$ であって
$\Delta : U \to U \times U$ と両立するもの，および
$i : R \to R$ であって因子を入れ替える写像 $U \times U \to U \times U$
と両立するものを構成する．
$k$ を代数閉体とする．すると
$$
\Mor(\Spec(k), \mathcal{X}_i) \to \Mor(\Spec(k), M_i) = M_i(k)
$$
は同型類上全単射であり，射 $M' \to M$ による任意の基底変換の後にも
% P12-ERR-0124: Preserve the frozen premature/wrong base-change target M.
同じことが成り立つ．これは $f_i$ の選び方と
Morphisms of Stacks, Lemmas
\ref{stacks-morphisms-lemma-universally-injective} および
\ref{stacks-morphisms-lemma-universally-injective-point}
から従う．$2$-ファイバー積の構成により，図式
$$
\xymatrix{
\Mor(\Spec(k), \mathcal{X}_{ij}) \ar[d] \ar[r] &
\Mor(\Spec(k), \mathcal{X}_j) \ar[d] \\
\Mor(\Spec(k), \mathcal{X}_i) \ar[r] &
\Mor(\Spec(k), \mathcal{X})
}
$$
は圏のファイバー積である．$g_i$ の選び方により，
この図式中の関手は自己同型群上に全単射を誘導する．
したがって，この図式は同型類の集合上にファイバー積図式を誘導する！
これにより
$$
R(k) = U(k) \times_{|\Mor(\Spec(k), \mathcal{X})|} U(k)
$$
が分かる．ここで $|\Mor(\Spec(k), \mathcal{X})|$ は
同型類全体の集合を表す．特に，任意の代数閉体 $k$ に対して，
% P12-ERR-0125: Japanese uses the collective point-set reading of frozen singular `point`.
$k$-値点上の写像は同値関係である．
Groupoids, Lemma \ref{groupoids-lemma-etale-equivalence-relation}
により，主張が成り立つと結論する．

\medskip\noindent
$M = U/R$ を，上の \'etale 同値関係による商である代数空間とする．
Spaces, Theorem \ref{spaces-theorem-presentation} を参照されたい．
標準的な射 $f : \mathcal{X} \to M$ が存在し，これは可換図式
\begin{equation}
\label{stacks-more-morphisms-equation-fundamental-diagram}
\xymatrix{
\mathcal{X}_i \ar[r]_{g_i} \ar[d]_{f_i} & \mathcal{X} \ar[d]^f \\
M_i \ar[r] & M
}
\end{equation}
に入る．
すなわち，このような射 $f$ は関手
$$
f : \Mor(T, \mathcal{X}) \longrightarrow \Mor(T, M)
$$
によって与えられる．これは任意のスキーム $T$ に対して定まり，
基底変換と両立する．$a : T \to \mathcal{X}$ を左辺の対象とする．
\'etale 被覆 $\{T_i \to T\}$ であって
$T_i = \mathcal{X}_i \times_\mathcal{X} T$ となるものと，射
$a_i : T_i \to \mathcal{X}_i$ が得られる．すると
$b_i = f_i \circ a_i : T_i \to M_i$ が得られる．
$T_i \times_T T_j = \mathcal{X}_{ij} \times_\mathcal{X} T$
なので，さらに射 $a_{ij} : T_i \times_T T_j \to \mathcal{X}_{ij}$
が得られる．$b_{ij} = f_{ij} \circ a_{ij}$ とおくと，
$b_i \times b_j$ はモノ射 $M_{ij} \to M_i \times M_j$ を経由することが分かる．
したがって射
$$
T_i \xrightarrow{b_i} M_i \to M
$$
は $T_i \times_T T_j$ 上で一致する．$M$ は \'etale 位相に関する層なので，
これらの射は一意な射 $b = f(a) : T \to M$ に貼り合わさることが分かる．
この構成が基底変換と両立することの確認，および図式
(\ref{stacks-more-morphisms-equation-fundamental-diagram}) が可換であることの確認は省略する．

\medskip\noindent
主張：図式 (\ref{stacks-more-morphisms-equation-fundamental-diagram}) はカルテジアンである．
これを見るため，誘導される射
$$
h_i : \mathcal{X}_i \longrightarrow M_i \times_M \mathcal{X}
$$
を調べる．これは \'etale な $\mathcal{X}$ 上のスタック間の射であり，
したがって $h_i$ は \'etale である
（Morphisms of Stacks, Lemma
\ref{stacks-morphisms-lemma-etale-permanence}）．
$g_i$ は分離なので，$h_i$ は分離であることが分かる
（Morphisms of Stacks, Lemma
\ref{stacks-morphisms-lemma-compose-after-separated} と，
上で見た $\mathcal{X}$ の対角射が分離であるという事実を用いる）．
射 $h_i$ は自己同型群上に同型を誘導する
（Morphisms of Stacks, Remark
\ref{stacks-morphisms-remark-identify-automorphism-groups}）．
これは $g_i$ について成り立つからである．
代数閉体 $k$ に対して，図式
$$
\xymatrix{
\Mor(\Spec(k), M_i \times_M \mathcal{X}) \ar[r] \ar[d] &
\Mor(\Spec(k), \mathcal{X}) \ar[d] \\
M_i(k) \ar[r] &
M(k)
}
$$
% P12-ERR-0126
はカルテジアンな圏の図式であり，上の矢印は自己同型群上に
全単射を誘導する．一方，
$$
M(k) = U(k)/R(k) = U(k)/
U(k) \times_{|\Mor(\Spec(k), \mathcal{X})|} U(k) =
|\Mor(\Spec(k), \mathcal{X})|
$$
が上で述べたことから成り立つ．したがって，上のカルテジアン図式の
右の縦矢印は同型類上の全単射である．
よって
$|\Mor(\Spec(k), M_i \times_M \mathcal{X})| \to M_i(k)$ は全単射である．
% P12-ERR-0127: Japanese removes the frozen extraneous article in the property list.
確認すると，$h_i$ は分離かつ \'etale であり，自己同型群上に同型を誘導し
（Morphisms of Stacks, Remark
\ref{stacks-morphisms-remark-identify-automorphism-groups} における意味で），
代数閉体上のファイバー圏の間に同値を誘導する．
したがって Morphisms of Stacks, Lemma
\ref{stacks-morphisms-lemma-etale-iso} により同型である．

\medskip\noindent
この主張から，特に次が得られる：
全射な \'etale 射 $U \to M$ であって，$f$ の基底変換が
分離，準コンパクト，かつ普遍同相射となるものが存在する．
したがって $f$ は分離，準コンパクト，かつ普遍同相射である．
これについては Morphisms of Stacks, Lemma
\ref{stacks-morphisms-lemma-check-separated-covering}，
\ref{stacks-morphisms-lemma-check-quasi-compact-covering}，および
% P12-ERR-0128: Preserve the frozen missing terminal punctuation.
\ref{stacks-morphisms-lemma-check-universal-homeomorphism-covering}
を参照されたい

\medskip\noindent
証明を終えるには，$f : \mathcal{X} \to M$ が
一様圏論的モジュライ空間であることを示さなければならない．
これを証明するには，代数空間の平坦な射 $M' \to M$ が与えられたとき，
基底変換
$$
M' \times_M \mathcal{X} \longrightarrow M'
$$
が圏論的モジュライ空間であることを示せば十分である．そこで，射
$$
\theta : M' \times_M \mathcal{X} \longrightarrow E
$$
を考える．ここで $E$ は代数空間である．各 $i$ に対して，
$f_i$ は一様圏論的モジュライ空間であることが分かっている．したがって
$$
\xymatrix{
M' \times_M \mathcal{X}_i \ar[d] \ar[r] &
M' \times_M \mathcal{X} \ar[d]^\theta \\
M' \times_M M_i \ar[r]^{\psi_i} &
E
}
$$
を得る．$\{M' \times_M M_i \to M'\}$ は \'etale 被覆なので，
所望の射 $\psi : M' \to E$ を得るには，$\psi_i$ と $\psi_j$ が
$M' \times_M M_i \times_M M_j = M' \times_M M_{ij}$ 上で
一致することを示せば十分である．これは
% P12-ERR-0129: Preserve source terminology `uniform categorical quotient` here.
$f_{ij} : \mathcal{X}_{ij} =
\mathcal{X}_i \times_\mathcal{X} \mathcal{X}_j \to M_{ij}$
が一様圏論的商であるという事実から容易に従う．詳細は省略する．
最後に，$\psi$ が可換図式
$$
\xymatrix{
M' \times_M \mathcal{X} \ar[d] \ar[rd]^\theta \\
M' \ar[r]^\psi &
E
}
$$
に入ることを示す．これは ``$\{M' \times_M \mathcal{X}_i \to M' \times_M \mathcal{X}\}$
が \'etale 被覆である'' ことと，構成により射 $\psi_i$ が
対応する可換図式に入ることから従う．
これで Keel--Mori の定理の証明が完了する．
\end{proof}

\noindent
次の補題は，この構成の \'etale 局所的な性質を強調する．

\begin{lemma}
\label{stacks-more-morphisms-lemma-etale-separated-over-keel-mori}
$h : \mathcal{X}' \to \mathcal{X}$ を代数スタックの射とする．次を仮定する．
\begin{enumerate}
\item $\mathcal{I}_\mathcal{X} \to \mathcal{X}$ は有限である．
\item $h$ は \'etale かつ分離的で，自己同型群上に同型を誘導する
（Morphisms of Stacks, Remark
\ref{stacks-morphisms-remark-identify-automorphism-groups}）．
\end{enumerate}
このとき，カルテジアン図式
$$
\xymatrix{
\mathcal{X}' \ar[d] \ar[r] &
\mathcal{X} \ar[d] \\
M' \ar[r] &
M
}
$$
であって，$M' \to M$ が代数空間の分離的 \'etale 射であり，
縦の矢印が Theorem \ref{stacks-more-morphisms-theorem-keel-mori} で構成された
モジュライ空間への射であるものが存在する．
\end{lemma}

\begin{proof}
Morphisms of Stacks, Lemma
\ref{stacks-morphisms-lemma-aut-iso-unramified} により，
$\mathcal{I}_{\mathcal{X}'} \to
\mathcal{X}' \times_\mathcal{X} \mathcal{I}_\mathcal{X}$
は同型である．したがって
$\mathcal{I}_{\mathcal{X}'} \to \mathcal{X}'$ は
$\mathcal{I}_\mathcal{X} \to \mathcal{X}$ の基底変換として有限である．
$f' : \mathcal{X}' \to M'$ と $f : \mathcal{X} \to M$ を
Theorem \ref{stacks-more-morphisms-theorem-keel-mori} のものとする．
$f'$ は圏論的モジュライ空間なので，補題にある可換図式を得る．
$I$ と $g'_i : \mathcal{X}'_i \to \mathcal{X}'$ を
Lemma \ref{stacks-more-morphisms-lemma-etale-local-finite-inertia} のように選ぶ．
$g_i = h \circ g'_i$ は \'etale かつ分離的で，自己同型群上に同型を誘導することに
注意せよ（Morphisms of Stacks, Remark
\ref{stacks-morphisms-remark-identify-automorphism-groups}）．
$f'_i : \mathcal{X}'_i \to M'_i$ を Lemma
\ref{stacks-more-morphisms-lemma-well-nigh-affine-moduli-space} のものとする．
Theorem \ref{stacks-more-morphisms-theorem-keel-mori} の証明において，図式
$$
\xymatrix{
\mathcal{X}'_i \ar[d]_{f'_i} \ar[r]_{g'_i} &
\mathcal{X}' \ar[d]^{f'} \\
M'_i \ar[r] &
M'
}
\quad\text{および}\quad
\xymatrix{
\mathcal{X}'_i \ar[d]_{f'_i} \ar[r]_{g_i} &
\mathcal{X} \ar[d]^f \\
M'_i \ar[r] &
M
}
$$
がカルテジアンであり，$M'_i \to M'$ と $M'_i \to M$ が \'etale であることを
既に見た（これまでに挙げた射 $g'_i, g_i, f', f'_i, f$ の性質から，
% P12-ERR-0131: Japanese renders the intended adverb despite frozen `sofar`.
全く同じ議論によって直接にも従う）．
これから第一に $M' \to M$ が \'etale であり，第二に補題の図式が
カルテジアンであることが従う．なお $M' \to M$ が分離的であることを
示さなければならない．そのため図式
$$
\xymatrix{
\mathcal{X}' \ar[r] \ar[d] &
\mathcal{X}' \times_\mathcal{X} \mathcal{X}' \ar[d] \\
M' \ar[r] &
M' \times_M M'
}
$$
を考える．上の水平矢印は，$\mathcal{X}' \to \mathcal{X}$ が
分離的なので普遍閉である．縦の矢印は Theorem
\ref{stacks-more-morphisms-theorem-keel-mori} のもの（$\mathcal{X} \to M$ の平坦基底変換）なので，
普遍同相射である．したがって下の水平矢印は普遍閉である．
これは，それが代数空間の \'etale モノ射でもあることと合わせて，
所望のとおり閉埋め込みであることを示す．
\end{proof}

\section{モジュライ空間の性質}
\label{stacks-more-morphisms-section-properties-moduli-spaces}

\noindent
モジュライ空間の存在が証明されれば，そのモジュライ空間の性質を確立することが
可能になり，通常は容易である．

\begin{lemma}
\label{stacks-more-morphisms-lemma-keel-mori-finite-type}
$p : \mathcal{X} \to Y$ を代数スタックから代数空間への射とする．
次を仮定する．
\begin{enumerate}
\item $\mathcal{I}_\mathcal{X} \to \mathcal{X}$ は有限である．
\item $Y$ は局所 Noether 的である．
\item $p$ は局所有限型である．
\end{enumerate}
$f : \mathcal{X} \to M$ を Theorem \ref{stacks-more-morphisms-theorem-keel-mori} で構成された
モジュライ空間とする．このとき $M \to Y$ は局所有限型である．
\end{lemma}

\begin{proof}
$f$ は一様圏論的モジュライ空間なので，射 $M \to Y$ を得る．
$M \to Y$ が局所有限型であることを $M$ と $Y$ 上 \'etale 局所的に
確認すれば十分である．$f$ は一様圏論的モジュライ空間なので，まず $Y$ を
$Y$ 上 \'etale なアフィンスキームで置き換えてよい．
次に $I$ と $g_i : \mathcal{X}_i \to \mathcal{X}$ を Lemma
\ref{stacks-more-morphisms-lemma-etale-local-finite-inertia} のように選べる．
すると Lemma \ref{stacks-more-morphisms-lemma-etale-separated-over-keel-mori} により，
$\mathcal{X} = \mathcal{X}_i$ の場合へ帰着する．言い換えれば，
$\mathcal{X}$ はほとんどアフィンであると仮定してよい．
この場合 $Y = \Spec(A_0)$ であり，$\mathcal{X} = [U/R]$ で，
$U = \Spec(A)$ であり，$M = \Spec(C)$ である．ここで $C \subset A$ は
$R$-不変な $U$ 上の関数全体である．
Lemmas \ref{stacks-more-morphisms-lemma-well-nigh-affine} and
\ref{stacks-more-morphisms-lemma-well-nigh-affine-moduli-space} を参照せよ．
このとき $A_0$ は Noether 環であり，$A_0 \to A$ は有限型である．
さらに Groupoids, Lemma \ref{groupoids-lemma-integral-over-invariants}
により $A$ は $C$ 上整であり，したがって $C$ 上有限である
（$A_0$ 上有限型だからである）．最後に Algebra, Lemma
\ref{algebra-lemma-Artin-Tate} を適用して結論を得る．
\end{proof}

\begin{lemma}
\label{stacks-more-morphisms-lemma-keel-mori-diagonal}
$\mathcal{X}$ を代数スタックとし，
$\mathcal{I}_\mathcal{X} \to \mathcal{X}$ は有限であると仮定する．
$f : \mathcal{X} \to M$ を Theorem \ref{stacks-more-morphisms-theorem-keel-mori} で構成された
モジュライ空間とする．
\begin{enumerate}
\item $\mathcal{X}$ が準分離的ならば，$M$ は準分離的である．
\item $\mathcal{X}$ が分離的ならば，$M$ は分離的である．
\item ここにさらに追加する．例えば，上記の相対版などである．
\end{enumerate}
\end{lemma}

\begin{proof}
これを証明するため，図式
$$
\xymatrix{
\mathcal{X} \ar[d]_f \ar[r]_{\Delta_\mathcal{X}} &
\mathcal{X} \times \mathcal{X} \ar[d]^{f \times f} \\
M \ar[r]^{\Delta_M} &
M \times M
}
$$
を考える．$f$ は普遍同相射なので，$f \times f$ も普遍同相射である．

\medskip\noindent
$\mathcal{X}$ が分離的ならば $\Delta_\mathcal{X}$ は固有，したがって
$\Delta_\mathcal{X}$ は普遍閉であり，
ゆえに $\Delta_M$ は普遍閉である．したがって Morphisms of Spaces, Lemma
\ref{spaces-morphisms-lemma-separated-diagonal-proper}
により $M$ は分離的である．

\medskip\noindent
$\mathcal{X}$ が準分離的ならば $\Delta_\mathcal{X}$ は準コンパクトであり，
したがって $\Delta_M$ は準コンパクト，ゆえに $M$ は準分離的である．
\end{proof}

\begin{lemma}
\label{stacks-more-morphisms-lemma-keel-mori-proper}
$p : \mathcal{X} \to Y$ を代数スタックから代数空間への射とする．
次を仮定する．
\begin{enumerate}
\item $\mathcal{I}_\mathcal{X} \to \mathcal{X}$ は有限である．
\item $p$ は固有である．
\item $Y$ は局所 Noether 的である．
\end{enumerate}
$f : \mathcal{X} \to M$ を Theorem \ref{stacks-more-morphisms-theorem-keel-mori} で構成された
モジュライ空間とする．このとき $M \to Y$ は固有である．
\end{lemma}

\begin{proof}
Lemma \ref{stacks-more-morphisms-lemma-keel-mori-finite-type} により，$M \to Y$ は局所有限型である．
Lemma \ref{stacks-more-morphisms-lemma-keel-mori-diagonal} により，$M \to Y$ は分離的である．
もちろん $M \to Y$ は準コンパクトかつ普遍閉である．これらは位相的性質であり，
$\mathcal{X} \to Y$ はこれらの性質をもち，$\mathcal{X} \to M$ は
普遍同相射だからである．
\end{proof}

\section{スタックと fpqc 被覆}
\label{stacks-more-morphisms-section-fpqc}

\noindent
ある種の代数スタックは fpqc 降下を満たす．代数空間について本節に対応するものは
Properties of Spaces, Section \ref{spaces-properties-section-fpqc} である．


\begin{proposition}
\label{stacks-more-morphisms-proposition-stack-fpqc}
\begin{reference}
Anatoly Preygel による ``Intro to Algebraic Stacks'' の Proposition 3.3.6．
\end{reference}
$\mathcal{X}$ を，対角射が準アフィン\footnote{ind-準アフィンと仮定するだけで十分である．}
である代数スタックとする．このとき $\mathcal{X}$ は fpqc 被覆に関する降下を満たす．
\end{proposition}

\begin{proof}
ここでの約束では，$\mathcal{X}$ は亜群のスタック
$p : \mathcal{X} \to (\Sch/S)_{fppf}$ であり，これは基底スキーム $S$ 上の
スキームの圏に fppf 位相を入れたものの上で考える．主張の意味は次のとおりである．
fpqc 被覆 $\mathcal{U} = \{U_i \to U\}_{i \in I}$ が
$S$ 上のスキームについて与えられたとき，関手
$$
\mathcal{X}_U \longrightarrow DD(\mathcal{U})
$$
は同値である．左辺は $\mathcal{X}$ の $U$ 上の対象の圏，右辺は
$\mathcal{X}$ 内の $\mathcal{U}$ に関する降下データの圏である．
Stacks, Section \ref{stacks-section-descent-data} の議論を参照せよ．

\medskip\noindent
充満忠実性．二つの対象 $x, y$ を $\mathcal{X}$ に取り，これらは $U$ 上にあるとする．
このとき $I = \mathit{Isom}(x, y)$ は $U$ 上の代数空間である．
$I$ の $U_i$ 上の切断の族であって，$U_i \times_U U_j$ への制限が一致するものは，
代数空間に対する命題の類似物，すなわち Properties of Spaces, Proposition
\ref{spaces-properties-proposition-sheaf-fpqc} により，$U$ 上の一意な切断から来る．
したがってこの関手は充満忠実である．

\medskip\noindent
本質的全射性．ここでは，対象 $x_i$ が $U_i$ 上に，また同型
$\varphi_{ij} : \text{pr}_0^*x_i \to \text{pr}_1^*x_j$ が
$U_i \times_U U_j$ 上に与えられ，
これらは $U_i \times_U U_j \times_U U_k$ 上でコサイクル条件を満たす．

\medskip\noindent
$W$ をアフィンスキームとし，$W \to \mathcal{X}$ を射とする．
各 $i$ に対し
$$
W_i = U_i \times_{x_i, \mathcal{X}} W
$$
を形成できる．射影 $W_i \to U_i$ は，$\mathcal{X}$ の対角射が
準アフィンなので準アフィンである．各組 $i, j \in I$ に対し，同型
$\varphi_{ij}$ は同型
$$
W_i \times_U U_j =
(U_i \times_U U_j) \times_{x_i \circ \text{pr}_0, \mathcal{X}} W
\to
(U_i \times_U U_j) \times_{x_j \circ \text{pr}_1, \mathcal{X}} W =
U_i \times_U W_j
$$
を誘導する．さらに，これらの同型は
$U_i \times_U U_j \times_U U_k$ 上でコサイクル条件を満たす．
言い換えれば，これらの同型はスキーム $W_i/U_i$ 上に
$\mathcal{U}$ に関する降下データを定める．Descent, Lemma
\ref{descent-lemma-quasi-affine} により，この降下データは有効である
\footnote{対角射が ind-準アフィンの場合には More on Groupoids, Lemma
\ref{more-groupoids-lemma-ind-quasi-affine} を用いてもよい．}．
したがって，準アフィン射 $W' \to U$ と，両正方形がカルテジアンである可換図式
$$
\xymatrix{
W' \ar[d] &
W_i \ar[l] \ar[d] \ar[r] &
W \ar[d] \\
U &
U_i \ar[l] \ar[r]^{x_i} &
\mathcal{X}
}
$$
が存在する．$\{W_i \to W'\}_{i \in I}$ は $\mathcal{U}$ の
$W' \to U$ による基底変換なので，fpqc 被覆である．$W$ は fpqc 被覆に関する
層条件を満たすから，一意な射 $W' \to W$ であって，
$W_i \to W' \to W$ が与えられた射 $W_i \to W$ となるものを得る．
言い換えれば，可換図式
$$
\xymatrix{
W_i \ar[d] \ar[r] &
W' \ar[d] \ar[r] &
W \ar[d] \\
U_i \ar[r] \ar@/_1pc/[rr]_{x_i} &
U &
\mathcal{X}
}
$$
をもち，これは同型 $\varphi_{ij}$ と両立し，その正方形と長方形は
カルテジアンである．

\medskip\noindent
アフィンスキームの族 $W_\alpha$（$\alpha \in A$）と平滑射
$W_\alpha \to \mathcal{X}$ を，$\coprod W_\alpha \to \mathcal{X}$ が
全射となるように選ぶ．前段落の手順により図式
$$
\xymatrix{
W_{\alpha, i} \ar[d] \ar[r] &
W_\alpha' \ar[d] \ar[r] &
W_\alpha \ar[d] \\
U_i \ar[r] \ar@/_1pc/[rr]_{x_i} &
U &
\mathcal{X}
}
$$
を各 $\alpha$ に対して作る．このとき射 $W_\alpha' \to U$ は平滑で，
合わせて全射である．

\medskip\noindent
$x_\alpha$ を，$\mathcal{X}$ の $W_\alpha'$ 上の対象であって，
$W_\alpha' \to W_\alpha \to \mathcal{X}$ に対応するものと書く．
$\mathcal{X}$ は fppf スタックで，$\{W_\alpha' \to U\}$ は fppf 被覆なので，
同型 $\text{pr}_0^*x_\alpha \to \text{pr}_1^*x_\beta$ であって，
$W_\alpha' \times_U W'_\beta$ 上でコサイクル条件を満たすものが
存在することを示せば十分である．
一方，$W_{\alpha, i}$ へ引き戻した後には，
$W_{\alpha, i} \times_{U_i} W_{\beta, i} =
U_i \times_U (W_\alpha' \times_U W'_\beta)$ 上でそのような同型が存在する．
実際，$x_\alpha$ の $W_{\alpha, i}$ への引き戻しは，$x_i$ の
$W_{\alpha, i}$ への引き戻しと同型である．
$\{U_i \times_U (W_\alpha' \times_U W'_\beta) \to
W_\alpha' \times_U W'_\beta\}_{i \in I}$ は fpqc 被覆であり，
さらに上の図式と $\varphi_{ij}$ との既述の両立性により，これらの同型は
$W_\alpha' \times_U W'_\beta$ へ降下する．これで証明が完了する．
\end{proof}

\section{テンソル関手}
\label{stacks-more-morphisms-section-tensor-functors}

\noindent
$f : \mathcal{Y} \to \mathcal{X}$ を Noether 的代数スタックの射とする．
引き戻し関手
$$
f^* :
\textit{Coh}(\mathcal{O}_\mathcal{X})
\longrightarrow
\textit{Coh}(\mathcal{O}_\mathcal{Y})
$$
は右完全テンソル関手である．すなわち，加法的かつ右完全であり，
連接加群のテンソル積と可換する．任意の右完全テンソル関手
$F : \textit{Coh}(\mathcal{O}_\mathcal{X}) \to
\textit{Coh}(\mathcal{O}_\mathcal{Y})$ が，どの程度まで射
$f : \mathcal{Y} \to \mathcal{X}$ から生じるかを問うことができる．
この種の非常に一般的な結果については
\cite{Hall-Rydh-coherent} を参照されたい．
本節の目的は，これらの考え方への導入として Theorem \ref{stacks-more-morphisms-theorem-main}
の短い証明を与えることである．

\medskip\noindent
いくつかの補題から始める．

\begin{lemma}
\label{stacks-more-morphisms-lemma-extend}
$\mathcal{X}$ と $\mathcal{Y}$ を Noether 的代数スタックとする．
任意の右完全テンソル関手 $F : \textit{Coh}(\mathcal{O}_\mathcal{X}) \to
\textit{Coh}(\mathcal{O}_\mathcal{Y})$ は，すべての余極限と可換する
右完全テンソル関手
$F : \QCoh(\mathcal{O}_\mathcal{X}) \to \QCoh(\mathcal{O}_\mathcal{Y})$
へ一意に拡張される．
\end{lemma}

\begin{proof}
拡張の存在と一意性は一般的事実である．Categories, Lemma
\ref{categories-lemma-extend-functor-by-colim} を参照せよ．
この補題が適用できることを見るため，局所 Noether 的代数スタック上の
連接加群は定義により有限表示加群であることに注意する．Cohomology of Stacks,
Definition \ref{stacks-cohomology-definition-coherent} を参照せよ．
したがって $\mathcal{X}$ 上の連接加群は
$\QCoh(\mathcal{O}_\mathcal{X})$ の圏論的コンパクト対象である．これは
Cohomology of Stacks, Lemma
\ref{stacks-cohomology-lemma-finite-presentation-quasi-compact-colimit} による．
最後に，Cohomology of Stacks, Lemma
\ref{stacks-cohomology-lemma-directed-colimit-coherent}
により，任意の準連接加群はその連接部分加群のフィルター余極限である．

\medskip\noindent
$F$ は加法的なので，$F$ の拡張も加法的である（詳細は省略する）．
$F$ はテンソル関手であり，加群の余極限はテンソル積の形成と可換するので，
$F$ の拡張もテンソル関手である（詳細は省略する）．

\medskip\noindent
この段落では，拡張が任意の直和と可換することを示す．
$\mathcal{F} = \bigoplus_{j \in J} \mathcal{H}_j$ であり，
$\mathcal{H}_j$ が準連接ならば，
$\mathcal{F} = \colim_{J' \subset J\text{ 有限}}
\bigoplus_{j \in J'} \mathcal{H}_j$ である．
$F$ の拡張も $F$ と書けば，
\begin{align*}
F(\mathcal{F})
& =
\colim_{J' \subset J\text{ 有限}}
F(\bigoplus\nolimits_{j \in J'} \mathcal{H}_j) \\
& =
\colim_{J' \subset J\text{ 有限}}
\bigoplus\nolimits_{j \in J'} F(\mathcal{H}_j) \\
& =
\bigoplus\nolimits_{j \in J} F(\mathcal{H}_j)
\end{align*}
を得る．したがって $F$ は任意の直和と可換する．

\medskip\noindent
この段落では，拡張が右完全であることを示す．
$0 \to \mathcal{F} \to \mathcal{F}' \to \mathcal{F}'' \to 0$
を準連接 $\mathcal{O}_\mathcal{X}$-加群の短完全列とする．
このとき $\mathcal{F}' = \bigcup \mathcal{F}'_i$ と，その連接部分加群の
合併として書く（上で挙げた参照箇所を見よ）．
% P12-ERR-0136: Japanese supplies the two prepositions omitted in the frozen source.
$\mathcal{F}''_i \subset \mathcal{F}''$ を $\mathcal{F}'_i$ の像と書き，さらに
$\mathcal{F}_i = \mathcal{F} \cap \mathcal{F}'_i =
\Ker(\mathcal{F}'_i \to \mathcal{F}''_i)$ と書く．すると明らかに
$\mathcal{F} = \bigcup \mathcal{F}_i$ および
$\mathcal{F}'' = \bigcup \mathcal{F}''_i$ であり，短完全列
$$
0 \to \mathcal{F}_i \to \mathcal{F}_i' \to \mathcal{F}_i'' \to 0
$$
を得る．拡張はフィルター余極限と可換するので，
$F(\mathcal{F}) = \colim_{i \in I} F(\mathcal{F}_i)$，
$F(\mathcal{F}') = \colim_{i \in I} F(\mathcal{F}'_i)$，および
$F(\mathcal{F}'') = \colim_{i \in I} F(\mathcal{F}''_i)$ である．
% P12-ERR-0137: Japanese is unaffected by the frozen source's agreement error.
加群の層のフィルター余極限は完全なので，$F$ の拡張は右完全であると結論する．

\medskip\noindent
すべての余積と可換する右完全関手はすべての余極限と可換するので，
証明は完了する．Categories, Lemma
\ref{categories-lemma-colimits-coproducts-coequalizers} を参照せよ．
\end{proof}

\begin{lemma}
\label{stacks-more-morphisms-lemma-affine}
$\mathcal{X}$ をアフィン対角射をもつ代数スタックとする．
$B$ を環とする．$F : \QCoh(\mathcal{O}_\mathcal{X}) \to \text{Mod}_B$
を，直和と可換する右完全テンソル関手とする．
$g : U \to \mathcal{X}$ を $U = \Spec(A)$ がアフィンである射とする．このとき
\begin{enumerate}
\item $C = F(g_{\QCoh, *}\mathcal{O}_U)$ は可換 $B$-代数であり，
\item 環準同型 $A \to C$ が存在する
\end{enumerate}
ので，$F \circ g_{\QCoh, *} : \text{Mod}_A \to \text{Mod}_B$ は
$M$ を $M \otimes_A C$ へ送り，これを $B$-加群とみなす．
\end{lemma}

\begin{proof}
$g$ は準コンパクトかつ準分離であることに注意する．Morphisms of Stacks,
Lemma \ref{stacks-morphisms-lemma-quasi-compact-quasi-separated-permanence}
を参照せよ．Cohomology of Stacks, Proposition
\ref{stacks-cohomology-proposition-direct-image-quasi-coherent}
において関手
$g_{\QCoh, *} : \QCoh(\mathcal{O}_U) \to \QCoh(\mathcal{O}_\mathcal{X})$
を構成した．Cohomology of Stacks, Remarks
\ref{stacks-cohomology-remark-direct-image-quasi-coherent-tensor} および
\ref{stacks-cohomology-remark-QCoh-tensor}
により乗法
$$
\mu :
g_{\QCoh, *}\mathcal{O}_U
\otimes_{\mathcal{O}_\mathcal{X}}
g_{\QCoh, *}\mathcal{O}_U
\longrightarrow
g_{\QCoh, *}\mathcal{O}_U
$$
を得る．これにより $g_{\QCoh, *}\mathcal{O}_U$ は可換
$\mathcal{O}_\mathcal{X}$-代数となる．したがって
$C = F(g_{\QCoh, *}\mathcal{O}_U)$ は $\text{Mod}_B$ における可換代数対象，
すなわち $C$ は可換 $B$-代数である．写像
$\kappa : A \to \text{End}_{\mathcal{O}_\mathcal{X}}(g_{\QCoh, *}\mathcal{O}_U)$
があり，任意の $a \in A$ に対して図式
$$
\xymatrix{
g_{\QCoh, *}\mathcal{O}_U \otimes_{\mathcal{O}_\mathcal{X}}
g_{\QCoh, *}\mathcal{O}_U
\ar[d]_{\kappa(r) \otimes 1} \ar[rr]_-\mu & &
g_{\QCoh, *}\mathcal{O}_U \ar[d]^{\kappa(r)} \\
g_{\QCoh, *}\mathcal{O}_U \otimes_{\mathcal{O}_\mathcal{X}}
g_{\QCoh, *}\mathcal{O}_U
\ar[rr]^-\mu & &
g_{\QCoh, *}\mathcal{O}_U
}
$$
% P12-ERR-0138: Frozen diagram retains undefined r rather than the quantified a.
は可換である．したがって写像
$\kappa' = F(\kappa) : A \to \text{End}_B(C)$
を得て，$\kappa'(a)(c) c' = \kappa'(a)(cc')$ が成り立つ．もちろんこれは，
$a \mapsto \kappa'(a)(1)$ が環準同型 $A \to C$ であることを意味する．

\medskip\noindent
射 $g : U \to \mathcal{X}$ はアフィンである．Morphisms of Stacks,
Lemma \ref{stacks-morphisms-lemma-affine-permanence} を参照せよ．
したがって $g_{\QCoh, *}$ は完全であり，直和と可換する．これは
Cohomology of Stacks, Lemma
\ref{stacks-cohomology-lemma-quasi-coherent-pushforward-affine} による．
ゆえに $F \circ g_{\QCoh, *} : \text{Mod}_A \to \text{Mod}_B$
は直和と可換する右完全関手であり，$A$ を $C$ へ送る．
Functors and Morphisms, Lemma \ref{functors-lemma-functor}
により，関手 $F \circ g_{\QCoh, *}$ は $A$-加群 $M$ を，
$M \otimes_A C$ へ送り，これを $B$-加群とみなすことが分かる．
\end{proof}

\begin{lemma}
\label{stacks-more-morphisms-lemma-universally-injective}
Lemma \ref{stacks-more-morphisms-lemma-affine} と同じ記法を用いる．$\mathcal{X}$ は
Noether 的であり，$g$ は全射かつ平坦であると仮定する．
このとき $B \to C$ は普遍単射である．
\end{lemma}

\begin{proof}
自然な写像
$1 : \mathcal{O}_\mathcal{X} \to g_{\QCoh, *}\mathcal{O}_U$
を $\QCoh(\mathcal{O}_\mathcal{X})$ において考える．これを $U$ へ引き戻し，
随伴を用いると，合成
$$
\mathcal{O}_U =
g^*\mathcal{O}_\mathcal{X} \xrightarrow{g^*1}
g^*g_{\QCoh, *}\mathcal{O}_U \to \mathcal{O}_U
$$
は $\QCoh(\mathcal{O}_U)$ における恒等写像であることがわかる．
$g_{\QCoh, *}\mathcal{O}_U = \colim \mathcal{F}_i$ を連接
$\mathcal{O}_\mathcal{X}$-加群のフィルター余極限として書く．Cohomology of Stacks, Lemma
\ref{stacks-cohomology-lemma-directed-colimit-coherent} を参照されたい．
十分大きい $i$ に対し，写像
$1 : \mathcal{O}_\mathcal{X} \to g_{\QCoh, *}\mathcal{O}_U$ は
$\mathcal{F}_i$ を経由して分解する．Cohomology of Stacks, Lemma
\ref{stacks-cohomology-lemma-finite-presentation-quasi-compact-colimit} を参照されたい．
$s : \mathcal{O}_\mathcal{X} \to \mathcal{F}_i$ をこの分解とする．このとき
$$
\mathcal{O}_U \xrightarrow{g^*s} g^*\mathcal{F}_i \to
g^*g_{\QCoh, *}\mathcal{O}_U \to \mathcal{O}_U
$$
は恒等写像である．言い換えると，$s$ は $U$ へ引き戻すと
直和因子の包含になる．
$\mathcal{F}_i^\vee = hom(\mathcal{F}_i, \mathcal{O}_\mathcal{X})$
とおく．記法は
Cohomology of Stacks, Lemma
\ref{stacks-cohomology-lemma-internal-hom-fp-into-qcoh} と同様である．
とくに評価写像
$ev : \mathcal{F}_i \otimes_{\mathcal{O}_\mathcal{X}} \mathcal{F}_i^\vee
\to \mathcal{O}_\mathcal{X}$ がある．
$s$ における評価は写像
$s^\vee : \mathcal{F}_i^\vee \to \mathcal{O}_\mathcal{X}$ を定める．
$s$ に関する主張の双対として，$g^*(s^\vee)$ は全射であることがわかる．
Cohomology of Stacks, Section
\ref{stacks-cohomology-section-further-remarks} を参照されたい．% P12-ERR-0139
ここでは $hom$ と $\otimes$ の
$U$ への制限との両立性を用いる．
$g$ は全射かつ平坦なので，$s^\vee$ は全射であると結論できる
（同所を参照）．$F$ は右完全なので，
$F(\mathcal{F}_i^\vee) \to F(\mathcal{O}_\mathcal{X}) = B$ は全射である．
$\lambda \in F(\mathcal{F}_i^\vee)$ を，その像が $1 \in B$ となるように選ぶ．
% P12-ERR-0140: see ../control/ERRATA_CANDIDATES.jsonl
また，$e = F(s)(1) \in F(\mathcal{F}_i)$ と記す．これは $1$ を写像
$F(s) : B = F(\mathcal{O}_\mathcal{X}) \to F(\mathcal{F}_i)$ で送った
像である．このとき写像
$$
F(ev) :
F(\mathcal{F}_i) \otimes_B F(\mathcal{F}_i^\vee) =
F(\mathcal{F}_i \otimes_{\mathcal{O}_\mathcal{X}} \mathcal{F}_i^\vee)
\longrightarrow
F(\mathcal{O}_\mathcal{X}) = B
$$
は構成により $e \otimes \lambda$ を $1$ へ送る．したがって，写像
$B \to F(\mathcal{F}_i)$，$b \mapsto be$ は普遍単射である．実際，片側逆写像
$F(\mathcal{F}_i) \to B$，$\xi \mapsto F(ev)(\xi \otimes \lambda)$ が存在する．
これは十分大きいすべての $i$ について成り立つので，結論を得る．
\end{proof}

\begin{lemma}
\label{stacks-more-morphisms-lemma-flat}
$B \to C$ を環準同型とする．次を仮定する．
\begin{enumerate}
\item 余積への二つの標準射 $C \to C \otimes_B C$ は平坦である．
\item $B \to C$ は普遍単射である．
\end{enumerate}
このとき $B \to C$ は忠実平坦である．
\end{lemma}

\begin{proof}
写像 $\Spec(C) \to \Spec(B)$ は，$B \to C$ が普遍単射なので全射である．
したがって，$B \to C$ が平坦であることを示せば十分であり，これは
Descent, Theorem \ref{descent-theorem-descend-module-properties} から従う．
\end{proof}

\noindent
% P12-ERR-0141: see ../control/ERRATA_CANDIDATES.jsonl
次に示す K\"unneth の非常に簡単な形は，代数スタック上の準連接加群の
コホモロジーに対する K\"unneth の定理についての節を書けば，
不要になるはずである．

\begin{lemma}
\label{stacks-more-morphisms-lemma-easy-kunneth}
$a : \mathcal{Y} \to \mathcal{X}$ および $b : \mathcal{Z} \to \mathcal{X}$ は，
スキームによって表現可能で，準コンパクト，準分離かつ平坦であるとする．
このとき $a_{\QCoh, *}\mathcal{O}_\mathcal{Y}
\otimes_{\mathcal{O}_\mathcal{X}}
b_{\QCoh, *}\mathcal{O}_\mathcal{Z} =
f_{\QCoh, *}\mathcal{O}_{\mathcal{Y} \times_\mathcal{X} \mathcal{Z}}$ である．ここで
$f : \mathcal{Y} \times_\mathcal{X} \mathcal{Z} \to \mathcal{X}$ は
明らかな射である．
\end{lemma}

\begin{proof}
$\mathcal{P} = \mathcal{Y} \times_\mathcal{X} \mathcal{Z}$ と略記する．
$a \circ \text{pr}_1 = f$ かつ $b \circ \text{pr}_2 = f$ なので，写像
$a_*\mathcal{O}_\mathcal{Y} \to f_*\mathcal{O}_\mathcal{P}$ および
$b_*\mathcal{O}_\mathcal{Z} \to f_*\mathcal{O}_\mathcal{P}$ を得る
（相対引き戻し写像を用いる．Sites, Section \ref{sites-section-pullback} を参照されたい）．
したがって相対カップ積
$$
\mu :
a_*\mathcal{O}_\mathcal{Y}
\otimes_{\mathcal{O}_\mathcal{X}}
b_*\mathcal{O}_\mathcal{Z} \longrightarrow
f_*\mathcal{O}_{\mathcal{Y} \times_\mathcal{X} \mathcal{Z}}
$$
を得る．$Q$ を適用し，テンソル積との両立性
（Cohomology of Stacks, Remark \ref{stacks-cohomology-remark-QCoh-tensor}）を用いると，
射
$Q(\mu) : a_{\QCoh, *}\mathcal{O}_\mathcal{Y}
\otimes_{\mathcal{O}_\mathcal{X}}
b_{\QCoh, *}\mathcal{O}_\mathcal{Z} \to
f_{\QCoh, *}\mathcal{O}_{\mathcal{Y} \times_\mathcal{X} \mathcal{Z}}$
を $\QCoh(\mathcal{O}_\mathcal{X})$ において得る．次に，スキーム $U$ と全射な平滑射
$U \to \mathcal{X}$ をとる．$Q(\mu)$ の $U_\etale$ への制限が
同型であることを示せば十分である．Cohomology of Stacks, Section
\ref{stacks-cohomology-section-further-remarks} を参照されたい．
さらに同じ節の内容により，$\mu$ の $U_\etale$ への制限が
同型であることを示せば十分である（ここでは $\mu$ の始域と終域が
基底変換性をもつ局所準連接加群であることを用いる）．
さらに，\'etale 位相で順像を計算してよい．Cohomology of Stacks, Proposition
\ref{stacks-cohomology-proposition-loc-qcoh-flat-base-change} を参照されたい．
このとき最終的に
$a_*\mathcal{O}_\mathcal{Y}|_{U_\etale} = (V \to U)_{small, *}\mathcal{O}_V$
であることがわかる．ここで $V = U \times_\mathcal{X} \mathcal{Y}$ である．
$b_*$ と $f_*$ についても同様である．したがって結果は，スキームの平坦，
準コンパクトかつ準分離な射に対する K\"unneth の公式から従う．Derived Categories of Schemes,
Lemma \ref{perfect-lemma-kunneth} を参照されたい．
\end{proof}

\begin{lemma}
\label{stacks-more-morphisms-lemma-fully-faithful}
$\mathcal{X}$ をアフィン対角射をもつ代数スタックとする．
$B$ を環とする．$f_i : \Spec(B) \to \mathcal{X}$，$i = 1, 2$ を
二つの射とする．$t : f_1^* \to f_2^*$ をテンソル関手
$f_i^* : \QCoh(\mathcal{O}_\mathcal{X}) \to \text{Mod}_B$ の同型とする．
このとき，$2$-射 $f_1 \to f_2$ で $t$ を誘導するものが存在する．
\end{lemma}

\begin{proof}
アフィンスキーム $U = \Spec(A)$ と全射な平滑射
$g : U \to \mathcal{X}$ をとる．Properties of Stacks, Lemma
\ref{stacks-properties-lemma-quasi-compact-stack} を参照されたい．
$\mathcal{X}$ の対角射はアフィンなので，
$U_i = \Spec(B) \times_{f_i, \mathcal{X}, g} U$ はアフィンである．
$U_i = \Spec(C_i)$ と書く．このとき $C_i$ は $B$-代数であり，環準同型
$A \to C_i$ とともに Lemma \ref{stacks-more-morphisms-lemma-affine} において
関手 $F = f_i^*$ を用いて構成されたものである．したがって $t$ は，
同型 $C_1 \to C_2$ を $B$-代数の同型として誘導し，これは環準同型
$A \to C_1$ および $A \to C_2$ と両立する．
% P12-ERR-0142: see ../control/ERRATA_CANDIDATES.jsonl
言い換えると，次の図式はいずれも可換である．
$$
\xymatrix{
U_i \ar[r] \ar[d] & U \ar[d]^g \\
\Spec(B) \ar[r]^{f_i} & \mathcal{X}
}
\quad\quad
\xymatrix{
&
U_2 \ar[ld] \ar[d]^{\cong} \ar[rd] \\
\Spec(B) &
U_1 \ar[l] \ar[r] &
U
}
$$
これにより，対象 $f_1$ と $f_2$ は，$\mathcal{X}$ の
$\Spec(B)$ 上の対象として，平滑被覆
$\{U_1 \to \Spec(B)\}$ に沿って引き戻すと同型になることがすでにわかる．
この局所同型を $f_1$ と $f_2$ の $\Spec(B)$ 上の同型に降下できることを示すには，
この同型（上の可換図式から得られるもの）が
$U_1 \times_{\Spec(B)} U_1$ 上の降下データと両立することを示す必要がある．
このため，$U \times_\mathcal{X} U$ もアフィンであり，
射 $g' : U \times_\mathcal{X} U \to \mathcal{X}$ が存在し，さらに
$$
U_i \times_{\Spec(B)} U_i =
\Spec(B) \times_{f_i, \mathcal{X}, g'}
(U \times_\mathcal{X} U)
$$
であることに注意する．したがって，同型 $C_1 \otimes_B C_1 \to C_2 \otimes_B C_2$ は
同型 $C_1 \to C_2$ から得られ，射
$U_i \times_{\Spec(B)} U_i \to U \times_\mathcal{X} U$ と両立する．
いくつかの詳細は省略する．
\end{proof}

\begin{lemma}
\label{stacks-more-morphisms-lemma-main}
$\mathcal{X}$ をアフィン対角射をもつ Noether 的代数スタックとする．
$B$ を環とする．
$F : \QCoh(\mathcal{O}_\mathcal{X}) \to \text{Mod}_B$ を，
直和と可換な右完全テンソル関手とする．
このとき $F$ は一意な射 $\Spec(B) \to \mathcal{X}$ に由来する．
\end{lemma}

\begin{proof}
$g : U \to \mathcal{X}$ を，$U = \Spec(A)$ がアフィンとなる
全射な滑らかな射として選ぶ．Properties of Stacks, Lemma
\ref{stacks-properties-lemma-quasi-compact-stack} を参照されたい．
Lemma \ref{stacks-more-morphisms-lemma-affine} を適用して，有限型可換 $B$-代数
$C = F(g_{\QCoh, *}\mathcal{O}_U)$ と環準同型 $A \to C$ を得る．
Lemma \ref{stacks-more-morphisms-lemma-universally-injective} により，
環準同型 $B \to C$ は普遍単射である．
代数
$$
C \otimes_B C =
F(g_{\QCoh, *}\mathcal{O}_U
\otimes_{\mathcal{O}_\mathcal{X}}
g_{\QCoh, *}\mathcal{O}_U)
$$
を考える．$g$ は平坦，準コンパクトかつ準分離であるから，
Lemma \ref{stacks-more-morphisms-lemma-easy-kunneth} により，次の式の最初の等号が成り立つ：
$$
g_{\QCoh, *}\mathcal{O}_U
\otimes_{\mathcal{O}_\mathcal{X}}
g_{\QCoh, *}\mathcal{O}_U
=
f_{\QCoh, *}\mathcal{O}_{U \times_\mathcal{X} U} =
g_{\QCoh, *}(\text{pr}_{2, *}\mathcal{O}_{U \times_\mathcal{X} U})
$$
ここで
$f : U \times_\mathcal{X} U \to \mathcal{X}$ は自然な射であり，
$\text{pr}_2 : U \times_\mathcal{X} U \to U$ は第二射影である．
第二の等号は Cohomology of Stacks, Lemma
\ref{stacks-cohomology-lemma-relative-leray} と
$f = g \circ \text{pr}_2$ から従う．$\mathcal{X}$ の対角射はアフィンなので，
$U \times_\mathcal{X} U = \Spec(R)$ はアフィンである．
$\text{pr}_2 : A \to R$ により $R$ を $A$-代数とみなすことにしよう．
以上を総合すると
$$
C \otimes_B C =
F(g_{\QCoh, *}\mathcal{O}_U
\otimes_{\mathcal{O}_\mathcal{X}}
g_{\QCoh, *}\mathcal{O}_U) =
F(g_{\QCoh, *}(\text{pr}_{2, *}\mathcal{O}_{U \times_\mathcal{X} U})) =
R \otimes_A C
$$
を得る．最後の等号は Lemma \ref{stacks-more-morphisms-lemma-affine} の最後の主張から従う．
$A \to R$ は平坦なので（これは $\text{pr}_2$ が
$U \to \mathcal{X}$ の基底変換として平坦だからである），
$C \otimes_B C$ は $C$ 上平坦であると結論される．
Lemma \ref{stacks-more-morphisms-lemma-flat} により，$B \to C$ は忠実平坦である．

\medskip\noindent
実線部分が可換となる次の図式が存在すると主張する：
$$
\xymatrix{
\Spec(C \otimes_B C) \ar@<1ex>[d] \ar@<-1ex>[d] \ar[r] &
U \times_\mathcal{X} U \ar@<1ex>[d] \ar@<-1ex>[d] \\
\Spec(C) \ar[d] \ar[r] &
U \ar[d] \\
\Spec(B) \ar@{..>}[r] &
\mathcal{X}
}
$$
矢印 $\Spec(C) \to U = \Spec(A)$ は，環準同型 $A \to C$ から得られる．
この環準同型は Lemma \ref{stacks-more-morphisms-lemma-affine} の主張に現れる．
同様に，矢印 $\Spec(C \otimes_B C) \to U \times_\mathcal{X} U$ は
環準同型 $R \to C \otimes_B C$ から得られる．
上の正方形が可換であることを確かめるには Lemma
\ref{stacks-more-morphisms-lemma-fully-faithful} を用いればよい；詳細は省略する．
点線の矢印 $\Spec(B) \to \mathcal{X}$ が得られることを
Proposition \ref{stacks-more-morphisms-proposition-stack-fpqc} により結論する．

\medskip\noindent
$F$ が点線の矢印による引き戻しに対応する関手であるという主張も，
以上と Lemma \ref{stacks-more-morphisms-lemma-affine} の対応する主張から明らかである．
詳細は省略する．
\end{proof}

\noindent
% P12-ERR-0143: Japanese supplies the grammatical notation relation omitted in source.
環 $B$ に対し，$\text{Mod}^{fg}_B$ で有限生成 $B$-加群
（別名，有限 $B$-加群）の圏を表すことにする．

\begin{theorem}
\label{stacks-more-morphisms-theorem-main}
$\mathcal{X}$ をアフィン対角射をもつ Noether 的代数スタックとする．
$B$ を Noether 環とする．
$F : \text{Coh}(\mathcal{O}_\mathcal{X}) \to \text{Mod}^{fg}_B$ を
右完全テンソル関手とする．
このとき $F$ は一意な射 $\Spec(B) \to \mathcal{X}$ に由来する．
\end{theorem}

\begin{proof}
Lemma \ref{stacks-more-morphisms-lemma-extend} により，$F$ を，すべての直和と可換な
% P12-ERR-0144: Japanese renders the intended term despite frozen `susms`.
右完全テンソル関手
$F : \QCoh(\mathcal{O}_\mathcal{X}) \to \text{Mod}_B$ に一意に拡張できる．
そこで Lemma \ref{stacks-more-morphisms-lemma-main} を適用できる．
\end{proof}
